\section{NN-A001: geometric formulation and dictionary}
\begin{geometricstatement}[NN-A001]
On the smooth locus of the analytic Cartier divisor $L_t(s)=0$, projection to the $t$-disc is locally an isomorphism and its section has tangent vector $\partial_t+v\partial_s$. With real normal coordinate $\Re s-1/2$, the normal component is $\Re v$. Its sign is the displayed Boolean observable; it is neither a group character of $\mathbb R/2\mathbb Z$ nor intrinsic to the reduced support without this deformation. This is a statement about analytic divisors and their normal direction, proved by the local implicit-function calculation.
\end{geometricstatement}
\begin{proof} The displayed dictionary transports the proof of NN-A001. Any additional realization, effectiveness or uniformity hypothesis is retained explicitly.\end{proof}


\section{NN-A002: geometric formulation and dictionary}
\begin{geometricstatement}[NN-A002]
For the pushforward empirical measure on $\mathbb R/2\mathbb Z$, the continuous characters are $x\mapsto e^{\pi ikx}$. The square wave has only odd Fourier modes. Thus the sufficient gate is the stated cancellation of odd characters together with zero mass at its discontinuities. A Gaussian limit after dilation lives on a different measure space and does not identify this compact pushforward. No Frobenius equidistribution theorem has been supplied for it.
\end{geometricstatement}
\begin{proof} The displayed dictionary transports the proof of NN-A002. Any additional realization, effectiveness or uniformity hypothesis is retained explicitly.\end{proof}


\section{NN-A003: geometric formulation and dictionary}
\begin{geometricstatement}[NN-A003]
The finite object is the projection kernel of $\operatorname{span}\{1,z,\ldots,z^{N-1}\}$ in $L^2(S^1)$. Marking an eigenangle replaces that projection by a rank-one smaller projection. The proved compression-determinant estimate and scaling limit identify the reduced Palm sine process and an independent Haar phase. This is a theorem in unitary harmonic analysis, with no asserted comparison morphism from an arithmetic zero sheaf.
\end{geometricstatement}
\begin{proof} The displayed dictionary transports the proof of NN-A003. Any additional realization, effectiveness or uniformity hypothesis is retained explicitly.\end{proof}


\section{NN-A004: geometric formulation and dictionary}
\begin{geometricstatement}[NN-A004]
In the complete local analytic ring at a simple zero, $L_t(s)=0$ defines a section $s=\rho+\sum_{n\ge1}c_nt^n$. All $c_n$ are determined by the displayed recursion. Multiplying $L$ by a local nonvanishing factor $G$ changes the first section jet by $G(\rho+1)/G(\rho)$ under this shifted deformation; hence the reduced divisor and its factorization alone do not determine the jet. The first nonzero real coefficient is the normal-order invariant relative to the chosen real line.
\end{geometricstatement}
\begin{proof} The displayed dictionary transports the proof of NN-A004. Any additional realization, effectiveness or uniformity hypothesis is retained explicitly.\end{proof}


\section{NN-A005: geometric formulation and dictionary}
\begin{geometricstatement}[NN-A005]
At multiplicity $m$, the local zero algebra is a finite flat analytic algebra of rank $m$ over the parameter disc after restricting to a small isolating contour. The base change $t=u^m$ separates the leading branches. The symmetric power sums, equivalently traces of multiplication by polynomials in $s$, descend to holomorphic functions of $t$ and have the displayed residue derivatives. This retains the nonreduced length $m$; one tangent vector per support point would lose it.
\end{geometricstatement}
\begin{proof} The displayed dictionary transports the proof of NN-A005. Any additional realization, effectiveness or uniformity hypothesis is retained explicitly.\end{proof}


\section{NN-A006: geometric formulation and dictionary}
\begin{geometricstatement}[NN-A006]
At a real heat collision of length $m$, the base change $t-t_0=u^2$ yields $m$ local sections with distinct leading slopes $\sqrt2 h_{m,j}$. The heat equation identifies the associated leading polynomial with a Hermite polynomial. Under A no such length-$m\ge2$ real collision occurs at positive time. This analytic ramification statement is not a claim of an algebraic arithmetic family or a pure Weil object.
\end{geometricstatement}
\begin{proof} The displayed dictionary transports the proof of NN-A006. Any additional realization, effectiveness or uniformity hypothesis is retained explicitly.\end{proof}


\section{NN-A007: geometric formulation and dictionary}
\begin{geometricstatement}[NN-A007]
The Toeplitz specialization sends elementary symmetric functions to the positive coefficients $c_k(t)$. Its minors are the corresponding skew Schur tests. For $t\ge0$ the positive product gives a totally nonnegative specialization; for $t<0$ at least one finite test is negative although all generators $c_k(t)$ remain positive. The distinction is between positivity of generators and positivity of the full Schur cone, not between two motives.
\end{geometricstatement}
\begin{proof} The displayed dictionary transports the proof of NN-A007. Any additional realization, effectiveness or uniformity hypothesis is retained explicitly.\end{proof}


\section{NN-A008: geometric formulation and dictionary}
\begin{geometricstatement}[NN-A008]
The trace-class density $D=\operatorname{diag}(\beta_j)$ gives the compatible characters $\operatorname{sgn}(\sigma)\operatorname{Tr}(D^{\otimes n}U_\sigma)$ of $S_n$, hence the Thoma point $(0,\beta,0)$. Exterior-power characters give $\prod_j\det((1-\beta_j)I+\beta_jU)$ on $U(\infty)$. For the finite positive Toeplitz flag chart, Rietsch Theorem~2.1, equations~(2.1)--(2.3), identifies $q_i=\Delta_{i-1}\Delta_{i+1}/\Delta_i^2>0$ with its quantum flag parameters. The finite Schur probability is $f^\lambda s_{\lambda^{\mathsf T}}(\beta)$; omitting $f^\lambda$ changes the measure. No integral positive basis or GCT obstruction is inferred.
\end{geometricstatement}
\begin{proof} The displayed dictionary transports the proof of NN-A008. Any additional realization, effectiveness or uniformity hypothesis is retained explicitly.\end{proof}


\section{NN-A009: geometric formulation and dictionary}
\begin{geometricstatement}[NN-A009]
The positive measure $\mu_a=\sum_\rho e^{-a\gamma_\rho^2}\delta_{\gamma_\rho^2}$ realizes the derivative matrices as Gram matrices of $1,x,\ldots,x^d$. Infinite distinct support makes every such matrix positive definite. Exponential moments make this measure determinate. This is an analytic positive Hilbert realization, distinct from the prior inverse-power determinant and not an independently constructed arithmetic cohomology.
\end{geometricstatement}
\begin{proof} The displayed dictionary transports the proof of NN-A009. Any additional realization, effectiveness or uniformity hypothesis is retained explicitly.\end{proof}


\section{NN-A010: geometric formulation and dictionary}
\begin{geometricstatement}[NN-A010]
The meromorphic family has principal parts at every $s=-2m$, with the displayed nonzero residues at $t=1$. Therefore its full analytic pole divisor cannot be identified with a principal-part object supported at only two points of $\mathbb P^1$. An algebraic two-point quotient may be defined, but a comparison with this family must explicitly discard or encode the remaining local terms.
\end{geometricstatement}
\begin{proof} The displayed dictionary transports the proof of NN-A010. Any additional realization, effectiveness or uniformity hypothesis is retained explicitly.\end{proof}


\section{NN-A011: geometric formulation and dictionary}
\begin{geometricstatement}[NN-A011]
The Weil form is the Hermitian pullback of $\sum_\rho m_\rho|\widehat h(\gamma_\rho)|^2$, while the branch form is the algebra trace on a quadratic algebra; the two forms are not identified. The fluid realization acts on the analytic zero divisor by a real translation, preserving each Artin germ and its multiplicity. Consequently the stated derivative singularity creates no nonreal zero under A. These concrete maps do not produce a finite pure motive or a universal obstruction to other maps.
\end{geometricstatement}
\begin{proof} The displayed dictionary transports the proof of NN-A011. Any additional realization, effectiveness or uniformity hypothesis is retained explicitly.\end{proof}


\section{NN-A012: geometric formulation and dictionary}
\begin{geometricstatement}[NN-A012]
For the printed finite-gerbe path integral, $H^0(S^4;\mathbb Z/n)$ has order $n$, $H^1$ is trivial and $H^2=0$, so conjugate quadratic actions both give $n$. The proposed tensor inverse therefore requires an Euler normalization. For adelic supports, the printed hypotheses yield a Boolean $\sigma$-algebra; arbitrary joins require an additional condition such as localizability. The explicit proper-curve counterexample proves that these two categorical upgrades cannot be made under the printed hypotheses.
\end{geometricstatement}
\begin{proof} The displayed dictionary transports the proof of NN-A012. Any additional realization, effectiveness or uniformity hypothesis is retained explicitly.\end{proof}


\section{NN-A013: geometric formulation and dictionary}
\begin{geometricstatement}[NN-A013]
The residue sum is a global trace of the infinitesimal shifted deformation on isolated finite pieces of the analytic zero divisor. At a multiple zero it uses the full meromorphic principal part, and thus its local algebra length, rather than an undefined individual velocity. The contour estimates establish asymptotic polynomial trace identities along good heights; these signed identities do not recover the Boolean sign distribution.
\end{geometricstatement}
\begin{proof} The displayed dictionary transports the proof of NN-A013. Any additional realization, effectiveness or uniformity hypothesis is retained explicitly.\end{proof}


\section{NN-A014: geometric formulation and dictionary}
\begin{geometricstatement}[NN-A014]
For each finite rank, the positive Toeplitz specialization determines a point of the positive flag chart and positive Schur tests. The explicit bound proves that this finite point remains positive on a negative-time interval even though the infinite PF specialization has already left its positive cone. The necessary rank/window escapes at least as $\sqrt{\log(1/|t|)/\log\log(1/|t|)}$. This is a quantified failure of bounded finite detection, not a derived complexity-class separation or a theorem of arithmetic purity.
\end{geometricstatement}
\begin{proof} The displayed dictionary transports the proof of NN-A014. Any additional realization, effectiveness or uniformity hypothesis is retained explicitly.\end{proof}


\section{NN-A015: geometric formulation and dictionary}
\begin{geometricstatement}[NN-A015]
The heat path gives actual Thoma points $(0,\beta(t),0)$ and actual $U(\infty)$ characters. The proved $\ell^1$ limit is $\beta_j^*=8/[\pi^2(2j-1)^2]$, whose generating function is $\cosh\sqrt{2w}$. Every fixed normalized Toeplitz flag point therefore converges to a positive $\mathbb Q$-rational point with coefficients $2^k/(2k)!$. Rietsch Theorem~2.1, equations~(2.1)--(2.3), identifies its quantum parameters. Their unnormalized versions satisfy $q_i(t)\sim1024\widetilde q_i^*/(\log t)^4$, hence approach the $q=0$ classical-cohomology specialization. This is a concrete analytic trajectory in the stated representation and flag objects, not a new pure motive, integral basis, or GCT separation.
\end{geometricstatement}
\begin{proof} The displayed dictionary transports the proof of NN-A015. Any additional realization, effectiveness or uniformity hypothesis is retained explicitly.\end{proof}


\section{NN-A016: geometric formulation and dictionary}
\begin{geometricstatement}[NN-A016]
The underlying objects are the band-limited sine projection and its codimension-one and codimension-two Palm subspaces. On each finite interval their compressions are positive trace-class contractions; the Fredholm determinant is the generating function of a Bernoulli spectral sum. This exact exterior-algebra data yields a local density limit and the periodic phase law. A two-point correlation function alone does not retain that full count determinant, and no arithmetic divisor comparison is asserted.
\end{geometricstatement}
\begin{proof} The displayed dictionary transports the proof of NN-A016. Any additional realization, effectiveness or uniformity hypothesis is retained explicitly.\end{proof}


\section{NN-A017: geometric formulation and dictionary}
\begin{geometricstatement}[NN-A017]
The empty-interval probability is the Fredholm exterior-algebra determinant of the sine projection restricted to an interval. Taylor expansion of its Fourier-feature map produces an explicit finite polynomial space with rational moment Gram matrix $G$ and determinant $\det(I-sDGD^*G)$. A Hilbert--Schmidt remainder controls trace norm and then the determinant. This gives an actual finite analytic certificate for the model probability; it is not a finite-field point count or an arithmetic motive.
\end{geometricstatement}
\begin{proof} The displayed dictionary transports the proof of NN-A017. Any additional realization, effectiveness or uniformity hypothesis is retained explicitly.\end{proof}


\section{NN-A018: geometric formulation and dictionary}
\begin{geometricstatement}[NN-A018]
The central analytic Cartier divisor has local length $m$; after the ramified base change $t=u^m$, its normalization has sections $s_j=1/2+c_ju+(C_K/m)c_j^2u^2+O(u^3)$. The coefficient $C_K$ retains the completion unit and shifted Euler value, which are invisible in the reduced support. The positive real parameter ray gives a canonical multiset of sides, not a single Boolean label on the unramified support or a colouring invariant under complex monodromy. In Martinet\textquotesingle s $Q_8$ field, the nontrivial integral $\mathbb Z[Q_8]$-module class gives an odd Artin central order by Fr\"ohlich\textquotesingle s theorem; the separately proved second-jet formula gives exactly two excess right branches. This is an explicit representation-to-functional-equation-to-analytic-divisor deduction, with no asserted pure motive or complexity separation.
\end{geometricstatement}
\begin{proof} The displayed dictionary transports the proof of NN-A018. Any additional realization, effectiveness or uniformity hypothesis is retained explicitly.\end{proof}


\section{NN-A019: geometric formulation and dictionary}
\begin{geometricstatement}[NN-A019]
The first-order generating function is $e^{\epsilon w}\prod_j(1+(1-\epsilon)\beta_j^*w)$, the genuine Thoma point $(0,(1-\epsilon)\beta^*,\epsilon)$. Thus every fixed Schur or character observable has this boundary-direction expansion even though the actual finite-time heat points have $\gamma=0$. The proved failure of an $\ell^1$ derivative records the compensating mass at indices tending to infinity. Exterior-power operations and their additive filtration identity are SGA~6, Expos\'e~V, Example~2.2(b) and Lemma~2.2.1; the representation-ring case is explicitly included there. Finite-dimensional traces give the stated unitary characters. This dictionary does not identify the analytic probability weights with integral vector-bundle multiplicities, a pure motive, or a GCT positive basis.
\end{geometricstatement}
\begin{proof} The displayed dictionary transports the proof of NN-A019. Any additional realization, effectiveness or uniformity hypothesis is retained explicitly.\end{proof}


\section{NN-A020: geometric formulation and dictionary}
\begin{geometricstatement}[NN-A020]
For a $k$ by $k$ elementary Toeplitz minor, the exact dual Jacobi--Trudi partition has $\lambda_1\le k$, irrespective of its degree or matrix indices. The two-variable sector calculation proves positivity of all such skew Schur tests, and their products preserve it by Cauchy--Binet. The exterior-power formalism is SGA~6, Expos\'e~V, Example~2.2(b) and Lemma~2.2.1, including its representation-ring example. Uniform angular confinement and compact convergence of the analytic divisor retain every fixed-width positive cone for a negative-time interval, including through multiple endpoint germs. The full infinite cone nevertheless fails at each negative time. Effective cone-complement certificates require no multiplicity identification. This is a theorem about these explicit analytic and exterior-power tests, not an arithmetic purity realization or a general complexity lower bound.
\end{geometricstatement}
\begin{proof} The displayed dictionary transports the proof of NN-A020. Any additional realization, effectiveness or uniformity hypothesis is retained explicitly.\end{proof}


\section{NN-A021: geometric formulation and dictionary}
\begin{geometricstatement}[NN-A021]
At each fixed $d$ the expansion $\det\Psi_t(U)=\sum_{\ell(\lambda)\le d}c_{\lambda,d}(t)s_\lambda(U)$ has strictly positive coefficients and $\sum_\lambda c_{\lambda,d}(t)\dim V_\lambda=1$ for a negative-time interval. These are actual probability measures on polynomial irreducible $U(d)$-representations of unbounded degree. The exact dictionary is positivity on $U(d)$ if and only if the coefficient sequence of $\Psi_t$ is $PF_d$ when $\Psi_t(0)>0$. It follows from the Weyl alternant, the ordinary Jacobi--Trudi formula, and nonnegative restriction multiplicities; their symmetry forces the required row bound. The exterior-power formalism is SGA~6, Expos\'e~V, Example~2.2(b) and Lemma~2.2.1, with the representation-ring case stated there. The compatible positive cone for all ranks fails at every negative time, so the least detecting rank escapes. These analytic real weights are not asserted integral bundle multiplicities or a GCT positive basis.
\end{geometricstatement}
\begin{proof} The displayed dictionary transports the proof of NN-A021. Any additional realization, effectiveness or uniformity hypothesis is retained explicitly.\end{proof}


\section{NN-A022: geometric formulation and dictionary}
\begin{geometricstatement}[NN-A022]
Let $m_x$ be the length of the completed local zero algebra at the positive real endpoint $x$. After $t=-u^2/2$, its leading normalization directions are $i h_{m_x,j}$, with $h_{m_x,j}$ the roots of $\operatorname{He}_{m_x}$. The proved uniform radial-tail bound promotes these local Artin data to the exact global limit $\lim_{a\downarrow0}\alpha(-a)/\sqrt a=\max_x\sqrt2 h_{m_x}/x$. Vanishing of this limit is equivalent to reducedness of every endpoint zero germ. The explicit intervals preserve every skew Schur test of width at most $r$ and every polynomial irreducible at unitary rank $d$, irrespective of degree. The exterior-power and exact-sequence dictionary is SGA~6, Expos\'e~V, Example~2.2(b) and Lemma~2.2.1; NN-A021 proves the exact $PF_d$ versus $U(d)$ criterion. These are analytic-divisor and representation-cone theorems with explicit certificate-size lower bounds, not a pure motive, integral canonical basis, first-nonreal-height formula, or general complexity separation.
\end{geometricstatement}
\begin{proof} The displayed dictionary transports the proof of NN-A022. Any additional realization, effectiveness or uniformity hypothesis is retained explicitly.\end{proof}


\section{NN-A023: geometric formulation and dictionary}
\begin{geometricstatement}[NN-A023]
The two selected zeros define a rank-two analytic finite algebra in the moving coordinate $z=b-2aL+\delta w$. Its discriminant has a unique simple zero at $a_c=\delta^2(1/8+\kappa/32+O(\kappa^2))$, and its normalization separates after a quadratic parameter extension. The complementary divisor contributes the trace $S=\sum(b-\xi)^{-2}$ and logarithmic drift $L$, retained rather than discarded. For this actual finite local algebra, SGA~1, Expos\'e~I, section~4, Proposition~10 identifies the etale locus by invertibility of its trace discriminant; section~3, Proposition~1 gives the vanishing-differentials criterion. The full infinite complementary divisor is not thereby a finite Weil object. The theorem reduces actual pair dynamics to the explicit scalar isolation parameter, with no assertion of arithmetic pair frequency.
\end{geometricstatement}
\begin{proof} The displayed dictionary transports the proof of NN-A023. Any additional realization, effectiveness or uniformity hypothesis is retained explicitly.\end{proof}


\section{NN-A024: geometric formulation and dictionary}
\begin{geometricstatement}[NN-A024]
For each CUE sample put $q=e^{-a/4}$ and $w=e^{iz}$. The polynomial $P_N(q,w)=\sum_{k=0}^N c_k q^{(2k-N)^2}w^k$ has leading and constant coefficients invertible over $\mathbb C[q,q^{-1}]$. Hence $\operatorname{Spec}\mathbb C[q,q^{-1},w,w^{-1}]/(P_N)\to\mathbb G_{m,q}$ is finite free of degree $N$. Its trace discriminant is invertible exactly on the etale locus by SGA~1, Expos\'e~I, section~4, Proposition~10, with flat/unramified Definition~1 and section~3, Proposition~1. At the first collision, with probability tending to one, exactly the smallest-gap rank-two component has simple ramification. The real circle is the fixed locus of $w\mapsto1/\overline w$ on the positive real parameter ray. This is an actual complex algebraic family and discriminant law, not an arithmetic Frobenius object or a canonical positive basis.
\end{geometricstatement}
\begin{proof} The displayed dictionary transports the proof of NN-A024. Any additional realization, effectiveness or uniformity hypothesis is retained explicitly.\end{proof}


\section{NN-A025: geometric formulation and dictionary}
\begin{geometricstatement}[NN-A025]
The actual algebra \[\mathbb C[q,q^{-1},w,w^{-1}]/\left(\sum_{k=0}^Nc_kq^{(2k-N)^2}w^k\right)\] is finite free of rank $N$ over $\mathbb C[q,q^{-1}]$. On a common small tubular analytic neighborhood of each bounded extreme-time interval, its tracked factors have ranks one or two. Each rank-two factor has at most one simple discriminant zero, with quadratic local normalization; each rank-one factor stays etale. The exact finite-locally-free trace-discriminant criterion is SGA~1, Expos\'e~I, section~4, Proposition~10, and the vanishing-differentials criterion is section~3, Proposition~1. After the stated scaling, the discriminant points and their circular positions have Poisson intensity $(3/2)\sqrt v\,dv\,db/(2\pi)$. On the common neighborhood, cluster labels are preserved and local monodromies exchange disjoint pairs, hence commute. This is an explicit random complex algebraic family, not an arithmetic Frobenius theorem, pure motive or positive integral basis.
\end{geometricstatement}
\begin{proof} The displayed dictionary transports the proof of NN-A025. Any additional realization, effectiveness or uniformity hypothesis is retained explicitly.\end{proof}


\section{NN-A026: geometric formulation and dictionary}
\begin{geometricstatement}[NN-A026]
The centered alphabet $x_j=\exp(i\vartheta(m+1-2j))$ is an actual diagonal unitary matrix. Its Weyl characters satisfy $s_\lambda(x)=\prod_{i<j}\sin(\vartheta(\lambda_i-\lambda_j+j-i))/\sin(\vartheta(j-i))\ge0$ on the stated rectangle of partition widths. Skew restriction multiplicities and the positive $\cosh\sqrt z$ specialization give strict finite-order Toeplitz tests. SGA~6, Expos\'e~V, Example~2.2(b) and Lemma~2.2.1 give the exact representation-ring exterior-power and filtration dictionary. After normalization at $1$, these are actual positive mixtures of every polynomial irreducible at each rank $d\le2n$, while the full compatible cone fails. The support retains the fixed nonreal pair. No integral multiplicity or pure-motive realization is asserted, and the family is not the heat family.
\end{geometricstatement}
\begin{proof} The displayed dictionary transports the proof of NN-A026. Any additional realization, effectiveness or uniformity hypothesis is retained explicitly.\end{proof}


\section{NN-A027: geometric formulation and dictionary}
\begin{geometricstatement}[NN-A027]
Under the displayed dominant-double hypothesis, an endpoint local Artin algebra of length two controls the maximal scaled angular displacement. Its quadratic normalization creates the conjugate pair $be^{\pm i\theta}$, and the complementary divisor has a uniform convergent product. The rectangular Schur trace has leading term $|D|\sin((n+1)\theta+\arg D)/\sin\theta$, with exponentially smaller remainder after the positive factor $(Pb)^n$ is removed. Dual Jacobi--Trudi makes this an actual order-$n$ minor in window $0,\ldots,k+n-1$. For the translated unitary alphabet it is the negative coefficient of $\det^k$ on $U(n)$. SGA~6, Expos\'e~V, Example~2.2(b), Lemma~2.2.1 give the exterior-power/representation dictionary; SGA~1, Expos\'e~I, section~4, Proposition~10 gives the finite rank-two discriminant criterion. The exact asymptotics apply to these certificate families under the additional multiplicity hypothesis, without implying that hypothesis or an integral motivic realization.
\end{geometricstatement}
\begin{proof} The displayed dictionary transports the proof of NN-A027. Any additional realization, effectiveness or uniformity hypothesis is retained explicitly.\end{proof}


\section{NN-A028: geometric formulation and dictionary}
\begin{geometricstatement}[NN-A028]
The first nonreduced endpoint algebra has length $m$ and leading heat normalization governed by $\operatorname{He}_m$. Its nonzero traceless Hermitian diagonal defines the genuine compact full-flag orbit $U(m)/U(1)^m$. The orbital Fourier transform must be negative somewhere by the explicitly proved Gaussian identity. The exact Weyl--HCIZ formula turns this into negative Schur values for integral highest weights of bounded height and width $O(a^{-1/2})$; prefix and infinite-tail factors are controlled with exponentially small errors. Dual Jacobi--Trudi gives the corresponding finite coefficient minor and the actual unitary irreducible coefficient. SGA~6, Expos\'e~V, Example~2.2(b), Lemma~2.2.1 give the representation-ring exterior-power formalism. No integrality of the Hermite orbit symplectic class is asserted. Reducedness of the endpoint analytic divisor is equivalent to faster growth of these concrete finite tests; no arithmetic purity or general complexity separation follows.
\end{geometricstatement}
\begin{proof} The displayed dictionary transports the proof of NN-A028. Any additional realization, effectiveness or uniformity hypothesis is retained explicitly.\end{proof}


\section{NN-A029: geometric formulation and dictionary}
\begin{geometricstatement}[NN-A029]
For each local Artin length $m$, the Hermite heat unfolding determines the actual full-flag coadjoint orbit of $D_m$. Its Haar orbital Fourier transform has the exact positivity radius $\mathfrak r_m=\inf\{\operatorname{diam}\operatorname{spec}X:\phi_m(X)<0\}$. The Helmholtz equation makes positivity strict below this radius; the exact Weyl--HCIZ identity consequently controls every integral partition in the corresponding scaled-width compact set. Negative orbital values give actual negative Schur coefficients after the proved coalescent-prefix and infinite-tail isolation. The minima $\min_x x\mathfrak r_{m_x}/(2\sqrt2)$ and $\min_x x(1-\beta_x)\mathfrak r_{m_x}/(2\sqrt2)$ are the exact scaled boundaries of the stated Toeplitz and unitary positive cones. SGA~6, Expos\'e~V, Example~2.2(b), Lemma~2.2.1 provide the representation-ring exterior-power formalism. No integrality of the orbit symplectic class, arithmetic Frobenius object or complexity separation is inferred.
\end{geometricstatement}
\begin{proof} The displayed dictionary transports the proof of NN-A029. Any additional realization, effectiveness or uniformity hypothesis is retained explicitly.\end{proof}


\section{NN-A030: geometric formulation and dictionary}
\begin{geometricstatement}[NN-A030]
The explicit polynomial $P_m(z)=i^m\operatorname{He}_m(z/i)$ has distinct roots, so $\mathbb Q[z]/(P_m)$ is finite etale by the finite-free discriminant criterion in SGA~1, Expos\'e~I, section~4, Proposition~10. Its companion matrix $B_m$ gives the algebraic connection $d-B_m\,dx$ on the trivial rank-$m$ bundle over $\mathbb A^1$. The analytic horizontal sections identify with $\ker P_m(D)$. Evaluation on $\coprod_\ell\operatorname{Spec}\mathbb C[x]/(x-t_\ell)^{\nu_\ell}$ is the actual Hermite jet map; its determinant is the confluent HCIZ numerator up to explicit nonzero factors. A zero of $W_{m,k}(x)$ means failure for two fat points of lengths $m-1-k$ and $k+1$. The exact orbital boundary is therefore a finite collection of one-variable jet-determinant degeneracies. The connection and finite etale spectral algebra are concrete; no arithmetic Frobenius structure, purity, or polynomial-time algorithm is inferred.
\end{geometricstatement}
\begin{proof} The displayed dictionary transports the proof of NN-A030. Any additional realization, effectiveness or uniformity hypothesis is retained explicitly.\end{proof}


\section{NN-A031: geometric formulation and dictionary}
\begin{geometricstatement}[NN-A031]
With the Haar Hermitian norm, the CUE projection space is the restriction of $H^0(\mathbb P^1,\mathcal O(N-1))$ to the unit circle. At two distinct marked points, polynomial division gives the exact evaluation sequence \[0\to(z-z_x)(z-z_y)\mathbb C[z]_{\le N-3}\to\mathbb C[z]_{\le N-1}\to\mathbb C^2\to0.\] The two-Palm projection is onto this vanishing subspace. Its exterior determinant controls the odd marked Fourier moments; the rank-two perturbation and Hilbert--Schmidt commutator bound give the quantitative empirical concentration. SGA~6, Expos\'e~V, Example~2.2(b), Lemma~2.2.1 provide the exterior-power formalism. These are concrete complex evaluation spaces with Haar probability, without an asserted arithmetic comparison, Frobenius purity, or integral multiplicity model.
\end{geometricstatement}
\begin{proof} The displayed dictionary transports the proof of NN-A031. Any additional realization, effectiveness or uniformity hypothesis is retained explicitly.\end{proof}


\section{NN-A032: geometric formulation and dictionary}
\begin{geometricstatement}[NN-A032]
The quadratic character is the actual one-dimensional Artin representation of the quadratic finite etale cover of $\operatorname{Spec}\mathbb Z[1/q]$. Its sign-$+1$ functional equation gives an even completed analytic germ at $s=1/2$. If its order is $m>0$, the local Cartier divisor has Artin algebra $\mathbb C[z]/(z^m)$; after the ramified base change $t=u^m$, normalization gives $m$ analytic sections $z=c_ju+(C_\chi/m)c_j^2u^2+O(u^3)$. Positivity of the paired Hadamard factors under A makes $C_\chi>0$ with the exact parity-dependent constants. The two vertical first jets, when present, therefore have negative second real normal component. This is a proved statement about a completed analytic divisor and its specified deformation. SGA~1, Expos\'e~I, section~3, Proposition~1 and section~4, Proposition~10 supply the differentials and discriminant dictionary for the finite local model. No arithmetic family realizing the shifted denominators, pure motive, or actual central zero is asserted.
\end{geometricstatement}
\begin{proof} The displayed dictionary transports the proof of NN-A032. Any additional realization, effectiveness or uniformity hypothesis is retained explicitly.\end{proof}


\section{NN-A033: geometric formulation and dictionary}
\begin{geometricstatement}[NN-A033]
The actual symplectic degree-two representation of Martinet\textquotesingle s $Q_8$-extension has central analytic length one by the proved bound and its negative root number. Its multiplicity two in the regular representation produces the local Dedekind algebra $\mathbb C[z]/(z^2)$. After $t=u^2$, the specified shifted divisor has sections $z_\pm=\pm iA_Nu+d_Nu^2+O(u^3)$, with common $d_N>0$. Their discriminant is $-4A_N^2t+O(t^2)$; hence the finite free cover is simply ramified at zero and etale off zero on a small disc, with transposition monodromy, by SGA~1, Expos\'e~I, section~4, Proposition~10. Reality and the involution $u\mapsto-u$ make the real normal coordinate $d_Nt+O(t^2)$. The unconditional rational test groups the analytic zero divisor into conjugate orbits; it does not require Frobenius positivity or a motivic realization of the shifted series. Length and reduced support alone do not record the completion-unit contribution to this second normal jet.
\end{geometricstatement}
\begin{proof} The displayed dictionary transports the proof of NN-A033. Any additional realization, effectiveness or uniformity hypothesis is retained explicitly.\end{proof}


\section{NN-A034: geometric formulation and dictionary}
\begin{geometricstatement}[NN-A034]
The explicitly defined compositum gives a finite etale $Q_8\times C_3$-cover of $\operatorname{Spec}\mathbb Z[1/(5\cdot41\cdot751)]$. The external tensor decomposition of its regular representation gives the exact Artin factorization, including ramified Euler factors before restriction. Its old symplectic factor contributes central length two, and the conjugate cubic degree-two twists contribute $4k$; the eight complex degree-one factors are nonzero by the proved finite certificate. The local Dedekind algebra is $\mathbb C[z]/(z^{2+4k})$. After $t=u^{2+4k}$, normalization gives its analytic sections; the two vertical first jets have negative second real normal coefficient because $C_K>0$. Thus exactly two more branches lie to the left. SGA~6, Expos\'e~V, Example~2.2(b), Lemma~2.2.1 provide the representation-ring formalism, and SGA~1, Expos\'e~I, section~4, Proposition~10 the finite-cover/discriminant dictionary. No arithmetic family realizing the analytic norm shift is asserted.
\end{geometricstatement}
\begin{proof} The displayed dictionary transports the proof of NN-A034. Any additional realization, effectiveness or uniformity hypothesis is retained explicitly.\end{proof}


\section{NN-A035: geometric formulation and dictionary}
\begin{geometricstatement}[NN-A035]
At a critical-line analytic Cartier zero of length $m$, completed reflection puts the first reduced logarithmic coefficient in $i\mathbb R$. In the genus-one factorization, this is exactly the assertion $\Re a=0$ for its exponential factor; reflection pairs $\alpha$ with $-\bar\alpha$, cancelling the real linear constants. The global divisor then gives the strictly positive real logarithmic derivative at displacement one. After $t=u^m$, the local algebra $\mathbb C[z]/(z^m)$ normalizes into $m$ sections. For $m\ge2$, the specified norm shift contributes no $t^2$ term at the second Puiseux order, and the exact gamma correction makes every vertical first direction acquire a strictly negative second normal component in the Dirichlet case and the displayed Dedekind range. SGA~1, Expos\'e~I, section~3, Proposition~1 and section~4, Proposition~10 give the finite differentials and discriminant dictionary. The simple case has an additional shift term; no arithmetic deformation, purity assertion or simple-transversality claim is inferred.
\end{geometricstatement}
\begin{proof} The displayed dictionary transports the proof of NN-A035. Any additional realization, effectiveness or uniformity hypothesis is retained explicitly.\end{proof}


\section{NN-A036: geometric formulation and dictionary}
\begin{geometricstatement}[NN-A036]
The analytic zero divisor is counted with its actual local lengths. At a reduced point, $d\zeta_t=\zeta\prime(\rho)\,ds-\rho\zeta(\rho+1)\,dt$, so its section has tangent $\partial_t+v_\rho\partial_s$. The observable $r_\rho=|\Re v_\rho|/|v_\rho|$ is its normalized real normal component relative to the critical line. The theorem proves a positive lower counting-density margin on this reduced locus while retaining every multiplicity in the ambient denominator. SGA~1, Expos\'e~I, section~3, Proposition~1 gives the local differentials dictionary. The characters $x\mapsto e^{i\gamma\log x}$ belong to the compact dual of $\mathbb Q_{>0}$, a product of prime circles. Landau\textquotesingle s fixed-frequency formula and absolute Dirichlet convergence identify the exterior ratio\textquotesingle s Haar constant term as $\sum\mu(n)n^{-3}=1/\zeta(3)$. The deterministic cosine inequality combines this with the single Hardy-phase mean. No Frobenius equidistribution for zeros, arithmetic shifted family, sign-balance law or complexity separation is asserted.
\end{geometricstatement}
\begin{proof} The displayed dictionary transports the proof of NN-A036. Any additional realization, effectiveness or uniformity hypothesis is retained explicitly.\end{proof}


\section{NN-A037: geometric formulation and dictionary}
\begin{geometricstatement}[NN-A037]
A primitive character is its actual rank-one Artin representation on the finite etale cyclotomic cover of $\operatorname{Spec}\mathbb Z[1/q]$. Its root-number half-phase makes the completed critical-line germ Real, also for complex characters. At a reduced zero, the specified norm shift has $dL_t=L\prime(\rho,\chi)\,ds-\rho L(\rho+1,\chi)\,dt$; hence its tangent section has normalized real normal component $r_{\rho,\chi}=|\Re v_{\rho,\chi}|/|v_{\rho,\chi}|$. The ambient divisor is counted with its actual local lengths. SGA~1, Expos\'e~I, section~3, Proposition~1 gives this differentials dictionary. The proved mixed-mode formula annihilates the first Hardy character against every fixed character of the compact dual of $\mathbb Q_{>0}$; absolute Euler convergence then gives the exact second moment $1/2$. The restricted, proved Rudnick--Sarnak correlation theorem bounds the nonreduced occurrence mass, producing the stated reduced-locus proportions. This is neither full product-Haar independence nor Frobenius equidistribution for zeros, and it does not construct an arithmetic family of the shifted divisors or a complexity separation. For $\zeta$, the classical $19/27$ reduced-occurrence proportion sharpens the same second-moment deduction to $11/54$ of all occurrences and $11/38$ of the reduced locus; this changes the mass bound, not the geometry or the unresolved orientation sign.
\end{geometricstatement}
\begin{proof} The displayed dictionary transports the proof of NN-A037. Any additional realization, effectiveness or uniformity hypothesis is retained explicitly.\end{proof}


\section{NN-A038: geometric formulation and dictionary}
\begin{geometricstatement}[NN-A038]
The distinct Hermite spectrum determines the actual compact orbit $U(m)/U(1)^m$, identified with the full complex flag variety of $\mathbb C^m$. Its specified Haar probability has Fourier transform $\phi_m$. The midpoint rescaling in the proof gives positivity whenever the spectral diameter is at most $2\varepsilon_0m/h_m$. Through the exact Weyl--HCIZ formula this controls every highest-weight partition in the corresponding scaled-width region; dual Jacobi--Trudi gives the actual Toeplitz minors and unitary irreducible coefficients. SGA~6, Expos\'e~V, Example~2.2(b), Lemma~2.2.1 supply the representation-ring Schur and exterior-power formalism. Thus the local Artin length $m$ contributes a factor of order $\sqrt m$ to the exact detection constants. No integral polarization of the nonintegral Hermite orbit, arithmetic cohomology realization, or complexity separation is inferred.
\end{geometricstatement}
\begin{proof} The displayed dictionary transports the proof of NN-A038. Any additional realization, effectiveness or uniformity hypothesis is retained explicitly.\end{proof}


\section{NN-A039: geometric formulation and dictionary}
\begin{geometricstatement}[NN-A039]
For the specified analytic norm shift, a length-$m$ local divisor is a finite locally free algebra over the parameter disc. After $t=u^m$, normalization gives its Puiseux sections, while their symmetric sum is the trace of the coordinate and descends to $t$. Its derivative is the full residue of $sL(s+1,\chi)/L(s,\chi)$, including nonreduced points. SGA~1, Expos\'e~I, section~3, Proposition~1 and section~4, Proposition~10 give the finite differentials and discriminant dictionary. The Hardy phase probe vanishes identically on vertical double clusters; their descended trace has strictly negative real derivative. Combining this with the exact mean probe and the small real global trace forces both normal orientations in every sufficiently high fixed multiplicative window. The Artin character is an actual rank-one arithmetic representation, but no arithmetic family realizing the denominator shift, Frobenius equidistribution, or complexity separation is asserted.
\end{geometricstatement}
\begin{proof} The displayed dictionary transports the proof of NN-A039. Any additional realization, effectiveness or uniformity hypothesis is retained explicitly.\end{proof}


\section{NN-A040: geometric formulation and dictionary}
\begin{geometricstatement}[NN-A040]
The finite locally free analytic divisor and its ramified normalized sections are those of NN-A039. Their symmetric coordinate trace descends to the original parameter, while the real normal-jet loci record the two orientations. SGA~1, Expos\'e~I, section~3, Proposition~1 and section~4, Proposition~10 provide the exact differentials and discriminant dictionary. The new estimate quantitatively pairs the Hardy character with finitely many characters of the compact dual of $\mathbb Q_{>0}$ through an absolute Euler truncation and a fully uniform stationary-phase estimate. It forces each signed jet locus to meet every stated shrinking-relative height window. No arithmetic family for the denominator shift, Frobenius equidistribution of zeros, positive full-jet Hilbert model, or complexity separation is inferred. In the sharper refinement, the finite character expansion is the actual Fourier decomposition on cyclotomic Galois quotients $(\mathbb Z/qb\mathbb Z)^*$. Parseval controls all auxiliary prime errors, and the exact induced Gauss sums make the character phases cancel. The local Euler factors yield the conductor-only constant $c_q=\frac5{4\pi}\prod_{p\mid q}\frac{p(p+1)}{p^2+1}$. This quantifies recurrence at square-root height scale without supplying a Frobenius law for the zero labels.
\end{geometricstatement}
\begin{proof} The displayed dictionary transports the proof of NN-A040. Any additional realization, effectiveness or uniformity hypothesis is retained explicitly.\end{proof}


\section{NN-A041: geometric formulation and dictionary}
\begin{geometricstatement}[NN-A041]
For each analytic Weierstrass cluster of the specified denominator shift, the derivative of the finite locally free coordinate trace is the full residue of $sL(s+1,\chi)/L(s,\chi)$. It is defined on the length-$m$ algebra, including its nilpotents, and descends from the sum of the ramified sections after $t=u^m$. Trace is additive over disjoint contour cells. Removing $J$ cells that meet a selected real normal-jet class leaves traces of one weak sign; their total is bounded above by the outer trace and the $J$ removed-cell traces. The positive phase probe therefore forces at least $T^{1-\varepsilon}\log T$ cells to meet each class in a dyadic window, up to the stated constants. SGA~1, Expos\'e~I, section~3, Proposition~1 and section~4, Proposition~10 give the finite-algebra, differentials and discriminant dictionary; the contour identity itself is proved in NN-A005 and NN-A039. The arithmetic rank-one Artin representation defines the original $L$-function, but no algebraic family, Frobenius counting law or complexity separation is asserted for its analytically moving zeros.
\end{geometricstatement}
\begin{proof} The displayed dictionary transports the proof of NN-A041. Any additional realization, effectiveness or uniformity hypothesis is retained explicitly.\end{proof}


\section{NN-A042: geometric formulation and dictionary}
\begin{geometricstatement}[NN-A042]
The real analytic Cartier divisor of $g=Z_\chi/\cos\psi_\chi$ is the difference of the Hardy and corrected-phase divisors. At a simple tangency, two local length-one factors cancel. At a vertical double, the length-two Hardy factor loses one phase factor and leaves a simple positively oriented crossing, with orientation fixed by the negative real leading discriminant coefficient. Absence of one specified normal-jet class makes all uncanceled zeros have a common crossing orientation, hence at most one more zero than pole on an interval. The exact phase-count identity and the proved positive shifted-moment tail therefore force positive counting proportions for the stated analytic normal-jet loci. SGA~1, Expos\'e~I, section~3, Proposition~1 and section~4, Proposition~10 supply the finite differentials and discriminant dictionary. This is a real ordered-divisor assertion with a specified phase, not an order on arbitrary complex schemes, a Frobenius counting law, an arithmetic family realizing the denominator shift, or a complexity separation. A single positive lower constant works for every fixed primitive character because the two microscopic moment errors have absolute leading bounds; only the eventual height threshold depends on the character. This is not uniform equidistribution in a growing conductor family. In the zeta-only signed-cancellation refinement, evaluation of $\Re\{\rho\zeta(\rho+1)\overline{\zeta^\prime(\rho)}\}/\gamma$ at the underlying point vanishes for every nonreduced zero center. On a reduced center it is the normal displacement multiplied by the positive derivative norm $|\zeta^\prime(\rho)|^2/\gamma$. Complex conjugation makes this a specified real analytic pairing, not a regular algebraic trace or a positive arithmetic measure. Its proved signed mean still needs absolute-mass estimates to count ordinary simple signs. The square-mollifier gate multiplies this normal statistic by the nonnegative Hermitian norm $|B_T(\rho)|^2$, preserving its strict sign wherever nonzero. A proved bilinear mean controls its signed mass; only if the separately stated square-mollifier moment ratio exceeds $1/2$ does the phase-discard argument give positive ordinary sign proportions. This norm remains real analytic rather than algebraic, and the strict moment inequality is not asserted as an input already established. NN-A047 subsequently proves the strict square-mollifier inequality and realizes this criterion for zeta under RH alone, with both ordinary sign proportions at least $10^{-14}$. The general fixed-character ordinary endpoint and the half-label law remain distinct. NN-A048 subsequently proves the corresponding actual moment and signed-weight interfaces for every fixed primitive character under A, closing that ordinary positive-proportion extension. The half-label law remains open.
\end{geometricstatement}
\begin{proof} The displayed dictionary transports the proof of NN-A042. Any additional realization, effectiveness or uniformity hypothesis is retained explicitly.\end{proof}


\section{NN-A043: geometric formulation and dictionary}
\begin{geometricstatement}[NN-A043]
The integers $a_K(p)$ count degree-one closed points of $\operatorname{Spec}\mathcal O_K$ above $p$. Their nonnegativity and $a_K(p)\le a_K(p)^2\le n a_K(p)$, together with the prime-ideal theorem, give the analytic variance lower bound without a conjectural Frobenius law for zeros. Each isolated length-$m$ zero of the specified norm shift gives an actual finite locally free analytic algebra over the parameter disc. After $\tau=u^m$, its normalization gives the Puiseux sections and real normal-jet classes. The real Cartier-divisor quotient $Z_K/\cos\psi_K$ cancels one factor at a vertical double and preserves the remaining positive crossing. Its crossing count, the explicit formula with both completion poles, and the proved positive shifted-moment tail force the stated proportions. SGA~1, Expos\'e~I, section~3, Proposition~1 and section~4, Proposition~10 give the finite differentials and discriminant dictionary. M remains a conditional analytic input. No arithmetic denominator-shift family, pure motive, order on arbitrary complex schemes, Frobenius distribution of zeros, or complexity separation is asserted.
\end{geometricstatement}
\begin{proof} The displayed dictionary transports the proof of NN-A043. Any additional realization, effectiveness or uniformity hypothesis is retained explicitly.\end{proof}


\section{NN-A044: geometric formulation and dictionary}
\begin{geometricstatement}[NN-A044]
At a reduced analytic zero, the specified shifted divisor projects isomorphically to the parameter disc and defines a section $s(\tau)=\rho+i\beta\tau+a_2\tau^2+\cdots$. Its first normal class relative to $\Re s=1/2$ vanishes. The proved Euler-sector constraint and reflected completed germ give \[\Re a_2=\beta d_\rho^{\,*}(\beta_\rho^*-\beta),\qquad d_\rho^{\,*}>0,\quad\beta_\rho^*>0.\] Consequently every downward tangent $\beta<0$ has a strictly negative second normal class. Upward tangents have the stated threshold, and equality retains a higher formal neighborhood. The reduced local finite algebra has vanishing relative differentials, as in SGA~1, Expos\'e~I, section~3, Proposition~1. The chosen positive real ray supplies the orientation; these are real analytic normal jets, not an intrinsic order on complex schemes. No arithmetic denominator-shift family, purity, tangent-zero existence, ordinary-sign density or complexity separation is asserted. For a quadratic number field, the degree-two finite map $\operatorname{Spec}\mathcal O_K\to\operatorname{Spec}\mathbb Z$ gives $\sum_{\mathfrak p\mid p}e_{\mathfrak p}f_{\mathfrak p}=2$, including ramification. Milne, Theorem~3.34, equation~(10), printed pp.~59--60, is the exact degree locator. At most two local Euler factors per rational prime give the proved sector, so the same ordered second-jet conclusion holds at the specified critical simple centers above height $24$. This concerns that analytic norm shift; no global zero-location premise or arithmetic realization of the shift is needed or inferred.
\end{geometricstatement}
\begin{proof} The displayed dictionary transports the proof of NN-A044. Any additional realization, effectiveness or uniformity hypothesis is retained explicitly.\end{proof}


\section{NN-A045: geometric formulation and dictionary}
\begin{geometricstatement}[NN-A045]
At each selected isolated zero, the local analytic zero algebra is reduced and its shifted divisor defines a unique parameter-disc section. The finite relative-differentials criterion is SGA~1, Expos\'e~I, section~3, Proposition~1. Its real normal tangent is $\Re v_\rho$, and the statistic \[K_\rho=\frac{|\zeta^\prime(\rho)|^2}{\gamma}\Re v_\rho\] is its pairing with a positive derivative norm. Evaluation gives zero at every nonreduced center. Blaschke factorization retains the complementary divisor in a nonnegative logarithmic distance potential. Separation of the selected reduced centers bounds the sum of those potentials even when the complement has arbitrary multiplicities. The resulting absolute mass and the proved signed cancellation force almost linear counts in both real normal-orientation loci. Complex conjugation and the chosen real ray are essential analytic data; this does not define an algebraic order, pure motive, Frobenius counting law or complexity separation. This derivative method retains a subpower loss; NN-A047 separately proves positive ordinary zeta proportions by square-mollifier moments. Half balance remains unresolved.
\end{geometricstatement}
\begin{proof} The displayed dictionary transports the proof of NN-A045. Any additional realization, effectiveness or uniformity hypothesis is retained explicitly.\end{proof}


\section{NN-A046: geometric formulation and dictionary}
\begin{geometricstatement}[NN-A046]
The factor $L^\prime(\rho,\chi)$ annihilates the statistic at every nonreduced zero center. On the reduced analytic locus, the specified denominator-shift divisor defines a unique local section, and its real normal tangent is paired with the positive derivative norm $|L^\prime(\rho,\chi)|^2/\gamma$. The finite relative-differentials criterion is SGA~1, Expos\'e~I, section~3, Proposition~1. The proof uses the actual finite character decomposition on $(\mathbb Z/q\mathbb Z)^*$: after primitive Gauss pairing, the principal component of the conjugated residue statistic is real; GRH and the fully proved Perron estimate control the other components. Thus concrete rank-one arithmetic characters control a real analytic normal statistic and the stated almost linear reduced-center counts. The argument does not construct an arithmetic denominator-shift scheme, a Frobenius distribution of zero labels, positive counting proportions, a pure motive or a complexity separation. For $q=1$ the finite nonprincipal part is absent, so the complete result requires RH alone. The separate NN-A047 square-mollifier deduction closes positive ordinary sign proportions for the zeta case under RH alone; the corresponding fixed-character moment interface is not supplied by this contour theorem. NN-A048 supplies the separate fixed-character square-mollifier calculation and closes the positive-proportion endpoint; it uses additional auxiliary family GRHs and does not alter the weaker input scope of this contour theorem.
\end{geometricstatement}
\begin{proof} The displayed dictionary transports the proof of NN-A046. Any additional realization, effectiveness or uniformity hypothesis is retained explicitly.\end{proof}


\section{NN-A047: geometric formulation and dictionary}
\begin{geometricstatement}[NN-A047]
On the reduced analytic zero divisor, the specified denominator-shift family has a unique local section with normal tangent $\Re v_\rho$ relative to the critical real line. The finite relative-differentials criterion is SGA~1, Expos\'e~I, section~3, Proposition~1. The real statistic \[K_\rho=\frac{|\zeta^\prime(\rho)|^2}{\gamma}\Re v_\rho\] vanishes at every nonreduced center. Multiplication by the Hermitian norm $|B_T(\rho)|^2$ preserves its strict sign. The proved moment concentration inequality and vanishing signed mean yield \[\liminf_{T\to\infty}\frac{\#\{\rho\ {\rm reduced}:T<\Im\rho\le2T,\ \pm\Re v_\rho>0\}}{N(2T)-N(T)}\ge10^{-14}.\] Thus both oriented reduced loci have positive counting proportions under RH alone, even with arbitrary multiplicities in the complementary divisor. The exact rational Bernstein certificate verifies the finite kernel inequality; the analytic localization and fixed-parameter smoothing supply its arithmetic interpretation. The Hermitian norm and real orientation are specified analytic data, not an order on arbitrary complex schemes, an arithmetic denominator-shift family, a Frobenius half-label law, a pure motive or a complexity separation.
\end{geometricstatement}
\begin{proof} The displayed dictionary transports the proof of NN-A047. Any additional realization, effectiveness or uniformity hypothesis is retained explicitly.\end{proof}


\section{NN-A048: geometric formulation and dictionary}
\begin{geometricstatement}[NN-A048]
At a reduced analytic zero, the specified denominator-shift divisor has a unique local section; the finite relative-differentials criterion is SGA~1, Expos\'e~I, section~3, Proposition~1. Its real normal tangent is paired with the positive derivative norm and then with the Hermitian square-mollifier norm: \[K_{\rho,\chi}|B_{\chi,T}(\rho)|^2=\frac{|L^\prime(\rho,\chi)|^2}{\gamma}\Re v_{\rho,\chi}\,|B_{\chi,T}(\rho)|^2.\] The statistic vanishes at all nonreduced centers. The actual finite Fourier decomposition on $(\mathbb Z/(qk)\mathbb Z)^*$ makes the principal Gauss scalar real, while character orthogonality controls each whole coefficient polynomial before summation. The proved moments and signed cancellation yield \[\liminf_{T\to\infty}\frac{\#\{\rho\ {\rm simple}:T<\Im\rho\le2T,\ \pm\Re v_{\rho,\chi}>0\}}{N_\chi(2T)-N_\chi(T)}\ge10^{-14}.\] Both oriented reduced loci therefore have positive counting proportions for every fixed primitive character under A. The norm and real orientation are specified analytic data; no algebraic ordering, arithmetic deformation scheme, Frobenius half-label law, pure motive or complexity separation is inferred. Auxiliary moduli grow in the proof, so A supplies genuinely additional family GRHs even though the counted character is fixed.
\end{geometricstatement}
\begin{proof} The displayed dictionary transports the proof of NN-A048. Any additional realization, effectiveness or uniformity hypothesis is retained explicitly.\end{proof}


\section{NN-A049: geometric formulation and dictionary}
\begin{geometricstatement}[NN-A049]
Under A, there is one absolute $c>0$ such that each fixed primitive character has two sets of at least $10^{-24}N_{\chi,T}$ reduced analytic zero centers. Their actual divisor sections are disjoint degree-one analytic covers of $|	au|le c/T$, with nonvanishing relative $s$ derivative throughout. Their real normal orientations remain strict for positive $	au$ on this disk. The same sections remain simple and disjoint on the larger disk
\[
 |	au|le \frac{\sqrt{\log\log T}}{100T},
\]
without an enlarged orientation assertion. The derivative section and the actual square-mollifier section define the five weighted analytic counting measures in the proof. Their marked-prime Gaussian moment factorization and the complete complementary zero-kernel estimate imply, for each strict normal orientation and every fixed finite interval $I$,
\[
 \liminf_{T\to\infty}\frac{\#\{\rho:\ \pm\Re v_\rho>0,\ 
 (\log|\Re v_\rho|-\log(T/\log T))/\sqrt{(\log\log T)/2}\in I\}}
 {N_{\chi,T}}\ge10^{-16}\mathbb P(G\in I).
\]
The same lower component holds for the full tangent norm and the normalized derivative norm. These are actual real analytic Hermitian norms and ordinary counts on the selected reduced locus, not a globally algebraic order or a Frobenius probability. SGA1, Expos\'e I, section3, Proposition1, \url{https://grothendiecksga.com/read/sga1/fr/I-3.html}, supplies the reduced differential dictionary; the analytic section is proved by explicit Rouch\'e and derivative bounds. The earlier $10^{-16}$ sublogarithmic disk retains its RH-alone zeta scope. No full CLT, half-label law, heat identification, arithmetic algebraization, or growing-conductor threshold is inferred.
The same degree-one analytic covers extend throughout the anisotropic region
\[
 (1+T|\tau|)e^{12T|\Im\tau|}\le c_1e^{\sqrt{(\log\log T)/2}}.
\]
For each fixed $A_0>0$, a smaller positive fraction depending on $A_0$ gives the corresponding region with right side $c_{A_0}e^{2A_0\sqrt{(\log\log T)/2}}$. The direct difference estimate and a fixed-disk zero count of one prove reducedness and unique patching throughout these regions. Real orientation is an additional condition and remains guaranteed only on the stated constant-over-height interval. No common arithmetic scheme over varying conductors is inferred.

There exist absolute $\kappa,c,T_0>0$ such that for every $T\ge T_0$ and every primitive character with $q\le T^\kappa$, the same analytic statement holds uniformly with two reduced sets of cardinality at least $10^{-30}N_{\chi,T}$. Each gives disjoint degree-one analytic covers of the uniform $c/T$ strict-side disk and of
\[
 (1+T|\tau|)e^{12T|\Im\tau|}
       \le c_1e^{\sqrt{(\log\log T)/2}}
\]
for simplicity and separation. For each fixed $A_0>0$, a uniform fraction $10^{-31}\mathbb P(G\in[A_0+1,A_0+2])$ has the corresponding region with right side $c_{A_0}e^{2A_0\sqrt{(\log\log T)/2}}$. The normalized logarithmic section-norm counting measures have the same uniform lower Gaussian component, with coefficient $10^{-30}$. Uniformity refers to this height-dependent finite collection of analytic zero divisors. It asserts no common finite-type arithmetic family over the conductor parameter, no Frobenius probability, and no full limiting measure. The actual derivative and fixed-disk zero-count criteria give the sections; the SGA1 reduced-differential dictionary above remains unchanged.

\end{geometricstatement}
\begin{proof} The displayed dictionary transports the proof of NN-A049. Any additional realization, effectiveness or uniformity hypothesis is retained explicitly.\end{proof}


\section{NN-A050: geometric formulation and dictionary}
\begin{geometricstatement}[NN-A050]
Let $\mathcal Z=\{(t,z):H_t(z)=0\}$ be the analytic zero divisor. For every fixed $0<\eta<2/3$ and sufficiently large $T$, at least \[(2/3-\eta)\bigl(N_H(2T)-N_H(T)\bigr)\] reduced centers over $t=0$ extend to pairwise disjoint degree-one sections over the common disc $|t|<c_\eta/(\log T)^2$. The local relative differentials vanish because $\partial_zH_t\ne0$, the analytic counterpart of SGA~1, Expos\'e~I, section~3, Proposition~1. Conjugation fixes each section over the real parameter interval, so the corresponding fiber points stay real and their local discriminants do not vanish. The proof retains every complementary multiplicity in the inverse-square interaction divisor and controls its average by restricted pair correlation. This is an analytic unramified-survival theorem, not an algebraization of the heat family, a Frobenius law, a global discriminant-exhaustion statement or a complexity separation. Under the PF/character dictionary, these real surviving factors remain positive elementary factors; the full negative-time character still fails positivity because other zeros can be nonreal.
\end{geometricstatement}
\begin{proof} The displayed dictionary transports the proof of NN-A050. Any additional realization, effectiveness or uniformity hypothesis is retained explicitly.\end{proof}


\section{NN-A051: geometric formulation and dictionary}
\begin{geometricstatement}[NN-A051]
The holomorphic zero locus $\mathcal Z=\{(t,z):H_t(z)=0\}$ is an analytic effective Cartier divisor. At the constructed points $(-a,z_a)$ and $(-a,\overline{z_a})$, the inequality $\partial_zH_{-a}\ne0$ gives two distinct local unramified sections over the parameter line, the analytic form of SGA~1, Expos\'e~I, section~3, Proposition~1. Their positive real coordinates satisfy \[\Re z_a\le\exp\!\left(\exp(C/\sqrt a)\right),\qquad |\Im z_a|\ge\tfrac12.\] This quantifies a non-real-supported fiber without asserting that the heat divisor algebraizes over a number field. Under RH, the PF/character correspondence of NN-A007--A008 interprets the nonreal zeros as failure of full total positivity at $-a$. The height estimate alone supplies no numerical upper witnessing Toeplitz order or unitary rank, because the rest of the zero alphabet enters those tests. Neither a Frobenius distribution nor a complexity separation is inferred. For a primitive character, the root-number-normalized completion has a Hermitian Fourier kernel, so conjugation fixes the real parameter fiber even when the kernel is not even. Its constructed pair gives the same reduced analytic divisor conclusion. For a number field $K$, the coefficients $b_K(m)$ count effective ideal divisors on $\operatorname{Spec}\mathcal O_K$; their positivity and the positive pole residue yield the analytic profile $\kappa_K E(w)/\sqrt a$. This connects actual arithmetic ideal counts to an unconditional non-real-supported heat fiber, without asserting an algebraization of the analytic family or discharging M in any other Dedekind theorem. Constants may depend on the fixed arithmetic object. Under RH, if the least multiple positive endpoint is $x_*$ with multiplicity $m$, its special analytic fiber is $\mathbb C[w]/(w^m)$. The base change $a=q^2$ followed by normalization separates the $m$ Hermite branches; the deck involution has $\lfloor m/2\rfloor$ transpositions and one fixed branch when $m$ is odd. Every nonreal branch has real coordinate $x_*-2\ell_*a+O(a^2)$. Cone compactness identifies this cluster as the first nonreal-height cluster. If the endpoint divisor is reduced everywhere, that first height instead tends to infinity. This conditional characterization does not decide the arithmetic reducedness of the endpoint divisor.
\end{geometricstatement}
\begin{proof} The displayed dictionary transports the proof of NN-A051. Any additional realization, effectiveness or uniformity hypothesis is retained explicitly.\end{proof}


\section{NN-C001: geometric formulation and dictionary}
\begin{geometricstatement}[NN-C001]
In the Boolean coordinate algebra $B_n=\mathbb F_2[x]/(x_i^2-x_i)$, evaluation gives $B_n\cong\prod_{x\in\mathbb F_2^n}\mathbb F_2$. An ideal in this product is determined by the coordinates on which it vanishes. Hence the original and compressed equations do generate the same ideal after passage to $B_n$, and define the same finite reduced, finite etale zero-dimensional scheme. The resulting isomorphism and its count are computed without listing its possibly $2^n$ components. The trace dictionary is SGA7 II XXII Corollary 3.2; the relative finite trace is SGA4 XVII 6.2.3. There is no new reductive-group action or representation multiplicity attached by this compression.
\end{geometricstatement}
\begin{proof} The displayed dictionary transports the proof of NN-C001. Any additional realization, effectiveness or uniformity hypothesis is retained explicitly.\end{proof}


\section{NN-C002: geometric formulation and dictionary}
\begin{geometricstatement}[NN-C002]
The field-coordinate construction is the elementary affine coordinate description of restriction of scalars for the finite separable extension $K/\mathbb F_2$, restricted to rational points. To obtain the finite counting scheme one also imposes $x_i^{2^r}-x_i=0$, or equivalently all bit relations $x_{ij}^2-x_{ij}=0$. Its coordinate algebra is a product indexed by the actual field-valued solutions. Counting is rank/trace; conditional sampling is repeated restriction to clopen subsets. No purity beyond the trivial weight-zero Frobenius representation on these rational components is required. SGA4 XVII 6.2.3 supplies the finite trace formalism; no claim of a canonical GCT representation is made.
\end{geometricstatement}
\begin{proof} The displayed dictionary transports the proof of NN-C002. Any additional realization, effectiveness or uniformity hypothesis is retained explicitly.\end{proof}


\section{NN-C003: geometric formulation and dictionary}
\begin{geometricstatement}[NN-C003]
For the finite split scheme $X=\operatorname{Spec}(\mathbb Q^{q^n})$, the product of local factors is an element $W\in\Gamma(X,\mathcal O_X)$, and its sum over assignments is the trace of multiplication by $W$ in that finite etale algebra. In the idempotent basis of this product algebra, multiplication by $W$ is diagonal with entries $W(x)$, so its trace is exactly the displayed sum. EGA IV4, printed p.111,18.2.1--18.2.2 is the original finite-algebra trace construction; equation18.2.1.2 states base-change compatibility. SGA4 XVII 6.2.3 provides the relative-dimension-zero trace formalism. This is a weighted endomorphism trace, not automatically a point count or a pure Frobenius trace: arbitrary rational or negative weights need not be eigenvalues of a geometric Frobenius with the desired weight. The natural basis is indexed by assignments; there is no promised short basis for the whole exponential-dimensional algebra.
\end{geometricstatement}
\begin{proof} The displayed dictionary transports the proof of NN-C003. Any additional realization, effectiveness or uniformity hypothesis is retained explicitly.\end{proof}


\section{NN-C004: geometric formulation and dictionary}
\begin{geometricstatement}[NN-C004]
The contraction is composition of evaluation/coevaluation pairings in finite-dimensional $K$-vector spaces, and can also be written as the finite-algebra trace of a diagonal weighted assignment operator as in NN-C003. The diagonal matrix calculation in NN-C003 supplies that algebraic trace dictionary directly; SGA4 XVII 6.2.3 is the verified finite-trace comparison. The temporary $V$ is a coefficient-encoding device; it is not a Frobenius coordinate, deformation of a motive, or canonical-basis parameter. Fixed algebraicity ensures exact arithmetic but does not imply purity or a representation-theoretic obstruction.
\end{geometricstatement}
\begin{proof} The displayed dictionary transports the proof of NN-C004. Any additional realization, effectiveness or uniformity hypothesis is retained explicitly.\end{proof}


\section{NN-C005: geometric formulation and dictionary}
\begin{geometricstatement}[NN-C005]
The Boolean quotient is the coordinate algebra of a finite split etale scheme. The coefficient functional above is its sum-of-evaluations trace over $\mathbb F_2$, expressed in the squarefree monomial basis. This basis has $2^n$ elements. Its finite scheme has weight-zero Frobenius but the trace modulo two is still $\oplus\mathrm P$-complete from a succinct equation presentation. SGA7 II XXII Corollary 3.2 gives exactly the coherent trace congruence, with properness satisfied for this finite scheme. It gives no polynomial algorithm for constructing an explicit cohomology basis.
\end{geometricstatement}
\begin{proof} The displayed dictionary transports the proof of NN-C005. Any additional realization, effectiveness or uniformity hypothesis is retained explicitly.\end{proof}


\section{NN-C006: geometric formulation and dictionary}
\begin{geometricstatement}[NN-C006]
The dictionary is literal: solution set to finite etale scheme, count to rank, parity to coherent Frobenius trace, and affine restriction to proper object minus boundary. On the representation side $H^0$ is a direct sum of trivial Frobenius representations indexed by solutions. Weight zero controls eigenvalue magnitude but not the cost of extracting the number of summands from equations. SGA7 II XXII 3.1--3.2 and SGA5 III 6.13 are the exact locators; no Grothendieck archive document is invented.
\end{geometricstatement}
\begin{proof} The displayed dictionary transports the proof of NN-C006. Any additional realization, effectiveness or uniformity hypothesis is retained explicitly.\end{proof}


\section{NN-C007: geometric formulation and dictionary}
\begin{geometricstatement}[NN-C007]
A restriction cuts a finite Boolean solution scheme by additional equations, equivalently selects a clopen collection of its rational components. Odd isolation asks that one selected finite scheme have odd degree; unique isolation asks for degree at most one. The finite trace of SGA4 XVII 6.2.3 records these degrees. A weight function's unique minimum is a combinatorial selection on this set, not a consequence of Frobenius purity. In the representation dictionary, direct sums encode alternative restrictions and tensor products encode simultaneous witness choices; the displayed parity OR uses their mod-two rank arithmetic.
\end{geometricstatement}
\begin{proof} The displayed dictionary transports the proof of NN-C007. Any additional realization, effectiveness or uniformity hypothesis is retained explicitly.\end{proof}


\section{NN-C008: geometric formulation and dictionary}
\begin{geometricstatement}[NN-C008]
Unsatisfiability is emptiness of the finite Boolean solution scheme. A computation transcript certifies this from its equations; it is not automatically an effective Nullstellensatz identity or a short free resolution of its coordinate algebra. In the finite etale model, emptiness means rank zero and zero $H^0$, while SGA7 II XXII 3.2 only identifies a residue of that rank. E computes the integer rank through its own algorithm, not through a claimed small cohomological presentation.
\end{geometricstatement}
\begin{proof} The displayed dictionary transports the proof of NN-C008. Any additional realization, effectiveness or uniformity hypothesis is retained explicitly.\end{proof}


\section{NN-C009: geometric formulation and dictionary}
\begin{geometricstatement}[NN-C009]
The quotient $B_n$ is the coordinate algebra of a finite split scheme; its idempotent decomposition makes semantic ideal membership transparent after all components are listed. The telescoping identity is an ideal-theoretic statement on that finite scheme. Producing polynomial-ring correction terms is a presentation problem, analogous to replacing an abstract zero object by an explicit exact complex. Rank or Frobenius-trace identities in SGA7 XXII do not control those presentation sizes. The translation is algebraic; no unproved categorical purity theorem is used.
\end{geometricstatement}
\begin{proof} The displayed dictionary transports the proof of NN-C009. Any additional realization, effectiveness or uniformity hypothesis is retained explicitly.\end{proof}


\section{NN-C010: geometric formulation and dictionary}
\begin{geometricstatement}[NN-C010]
The constant-replacement equations define an affine scheme over $\mathbb Z$. A closed point of its characteristic-zero fiber supplies $\operatorname{Spec}K$; reductions with residue degree one give rational points over finite fields. Its finite permutation representation has Dedekind zeta function $\zeta_K$. The relevant trace counts prime-ideal splitting, not satisfying assignments. The elementary matrix trace of multiplication supplies the finite-algebra language, and Burgisser Section 4.6 supplies the exact arithmetic estimate. A covers abelian character decompositions; the general field need not admit a decomposition into Dirichlet characters. This is the same type of analytic field gate as Koiran, but the output here is nonuniform circuit simulation.
\end{geometricstatement}
\begin{proof} The displayed dictionary transports the proof of NN-C010. Any additional realization, effectiveness or uniformity hypothesis is retained explicitly.\end{proof}


\section{NN-C011: geometric formulation and dictionary}
\begin{geometricstatement}[NN-C011]
An inversion fiber is a rational-point fiber of the polynomial map $F_n$. Collisions are rational points of its fiber product with itself, with the diagonal removed; the first-differing-bit decomposition is a disjoint clopen stratification of this finite Boolean scheme. Counts are degrees of those finite etale pieces, with trace formalism SGA4 XVII 6.2.3. This construction supplies an algorithm for the succinct fibers, not a theorem that GRH or purity makes arbitrary fibers easy.
\end{geometricstatement}
\begin{proof} The displayed dictionary transports the proof of NN-C011. Any additional realization, effectiveness or uniformity hypothesis is retained explicitly.\end{proof}


\section{NN-C012: geometric formulation and dictionary}
\begin{geometricstatement}[NN-C012]
Formal identity asks whether a regular function is zero in an affine coordinate ring. Boolean testing asks whether its restriction to the finite reduced subscheme cut out by $x_i^2-x_i$ is zero. The restriction map has a nonzero kernel, as the example shows. A trace on the finite subscheme forgets still more data. Finite-algebra matrix traces and the SGA7 XXII trace statement therefore do not identify formal PIT with Boolean zero testing. On the representation side, one must retain the polynomial family and its arithmetic circuit model; the finite evaluation representation is only a quotient.
\end{geometricstatement}
\begin{proof} The displayed dictionary transports the proof of NN-C012. Any additional realization, effectiveness or uniformity hypothesis is retained explicitly.\end{proof}


\section{NN-C013: geometric formulation and dictionary}
\begin{geometricstatement}[NN-C013]
On the finite Boolean solution scheme, the energy partitions the rational components into clopen subsets $X_j$. The enumerator is $\sum_j\operatorname{rank}H^0(X_j,\mathcal O)T^j$, equivalently the character of the graded vector space with one basis vector per solution in its energy degree. The finite-algebra trace is immediate in the component idempotent basis; SGA4 XVII 6.2.3 supplies the verified finite-trace formalism. The variable $T$ records this chosen grading; it does not turn $T^j$ into a Tate twist or a geometric Frobenius eigenvalue without a separate realization. Thus the result produces an effective graded rank computation but no plethysm multiplicity formula or GCT obstruction.
\end{geometricstatement}
\begin{proof} The displayed dictionary transports the proof of NN-C013. Any additional realization, effectiveness or uniformity hypothesis is retained explicitly.\end{proof}


\section{NN-C014: geometric formulation and dictionary}
\begin{geometricstatement}[NN-C014]
Each singleton system defines one rational reduced point, a proper finite etale scheme of degree one. A universal set of evaluation points must meet every such closed subscheme of the Boolean scheme, forcing every component to occur. In its coordinate algebra $\mathbb F_2^{2^n}$, these are the primitive coordinate idempotents. The sketch counterexample changes the ideal and doubles the selected clopen locus. Finite trace and degree are as in SGA4 XVII 6.2.3; weight-zero purity does not remove the need to retain each possible component in a universal evaluation set.
\end{geometricstatement}
\begin{proof} The displayed dictionary transports the proof of NN-C014. Any additional realization, effectiveness or uniformity hypothesis is retained explicitly.\end{proof}


\section{NN-C015: geometric formulation and dictionary}
\begin{geometricstatement}[NN-C015]
The verifier product gives a finite Boolean scheme together with its acceptance idempotent; summing the idempotent is the finite-algebra trace. Over $\mathbb Z$, the supplied integer presentation and formal-degree bound make all intermediate coefficients effectively bounded. The trace dictionary follows directly by using the component idempotent basis, as in NN-C003; compare SGA4 XVII 6.2.3. By contrast NN-C010 starts with arbitrary complex constants, whose replacement by finite-field data uses number-field points and Dedekind zeta estimates. These are different effective presentations of related rank/count constructions, not different purity statements.
\end{geometricstatement}
\begin{proof} The displayed dictionary transports the proof of NN-C015. Any additional realization, effectiveness or uniformity hypothesis is retained explicitly.\end{proof}


\section{NN-C016: geometric formulation and dictionary}
\begin{geometricstatement}[NN-C016]
The partition polynomial is the trace of multiplication by $\prod_f\Phi_f$ on the finite free algebra $\mathbb Q[y][x]/(x_i^2-x_i)$ over the parameter scheme $\operatorname{Spec}\mathbb Q[y]$. In the orthogonal idempotent basis indexed by bit assignments, this multiplication matrix is diagonal with the corresponding factor products on the diagonal. Its trace is their sum. Applying a coefficient-ring homomorphism entrywise commutes with summing the diagonal, proving the claimed base-change compatibility directly. The original EGA IV4, printed p.111,18.2.1--18.2.2 records this trace, with base change in equation18.2.1.2; compare the relative-dimension-zero formalism in SGA4 XVII 6.2.3. Specializing formal weights commutes with the trace and with the constructed arithmetic circuit. This is an actual family-level symbolic computation, unlike a single numerical rank. It still does not provide orbit-closure multiplicity obstructions, border lower bounds, or a finite Weil realization of nonstandard canonical bases.
\end{geometricstatement}
\begin{proof} The displayed dictionary transports the proof of NN-C016. Any additional realization, effectiveness or uniformity hypothesis is retained explicitly.\end{proof}


\section{NN-C017: geometric formulation and dictionary}
\begin{geometricstatement}[NN-C017]
After adjoining $x_i^{|K|}-x_i$, the original zero set is a finite reduced split scheme over $K$ whose rank is the count. The basis expansion and valid-code equations give a bijection of rational components with a finite Boolean scheme over $\mathbb F_p$; they are an effective coordinate presentation, not an identification of unrestricted affine zero schemes. The $p$-adic order in $h(Y)$ is a divisibility filtration on integer-valued indicators. The formal nilpotent $V$ records this filtration and is not a Frobenius coordinate or a Tate twist. On the resulting finite scheme, the component-idempotent basis makes counting the trace of the identity; SGA7 II XXII Corollary 3.2 supplies its mod-$p$ coherent trace expression and SGA4 XVII 6.2.3 the finite-trace formalism. The algorithm computes this succinctly presented rank without constructing the exponential basis and supplies no new representation-theoretic positivity statement.
\end{geometricstatement}
\begin{proof} The displayed dictionary transports the proof of NN-C017. Any additional realization, effectiveness or uniformity hypothesis is retained explicitly.\end{proof}


\section{NN-C018: geometric formulation and dictionary}
\begin{geometricstatement}[NN-C018]
A polynomial solution over a finite nonreduced ring is an $R$-point of the affine equation scheme, not in general a geometric point of a reduced scheme. The primary decomposition is a decomposition of $\operatorname{Spec}R$ into clopen components; its functor of points turns the equation solution set into the displayed Cartesian product. Additive coordinates and binomial digits then present this finite set by a finite reduced Boolean scheme over $\mathbb F_p$. That presentation preserves the set and its rank, not the nilpotent structure, tangent spaces or deformation theory of the original equation scheme. Counting is the trace of the identity on the product algebra indexed by these points, and weighted counting is the diagonal multiplication trace; these assertions follow directly in the idempotent basis, with SGA4 XVII6.2.3 as the finite-trace dictionary. SGA7 II XXII Corollary3.2 applies to the resulting finite scheme for its mod-$p$ trace, not automatically to the original nonreduced ring-valued functor. No canonical-basis or orbit-closure positivity is inferred.
\end{geometricstatement}
\begin{proof} The displayed dictionary transports the proof of NN-C018. Any additional realization, effectiveness or uniformity hypothesis is retained explicitly.\end{proof}


\section{NN-C019: geometric formulation and dictionary}
\begin{geometricstatement}[NN-C019]
The finite witness scheme has a diagonal grading operator with one basis vector of degree $w\cdot x$ for every witness. Its character is $\sum_{x\in\mathcal F}T^{w\cdot x}$, and the least-degree graded piece has rank one. Equivalently, $w$ exposes a single vertex of the convex hull of the witness bit vectors and selects one monomial in the initial form of the positive generating polynomial. The count/character interpretation follows directly from the component-idempotent basis and the finite-trace formalism of SGA4 XVII6.2.3. These are finite-set and torus-grading statements; they do not construct a GCT canonical basis or make a signed structure-constant formula positive. A numerical signed specialization forgets the grading and can reintroduce cancellation.
\end{geometricstatement}
\begin{proof} The displayed dictionary transports the proof of NN-C019. Any additional realization, effectiveness or uniformity hypothesis is retained explicitly.\end{proof}


\section{NN-C020: geometric formulation and dictionary}
\begin{geometricstatement}[NN-C020]
The coefficient object is a finite locally free affine scheme $\operatorname{Spec}B_n\to\operatorname{Spec}\mathbb Z$ with an explicit polynomial-rank trivialization of its pushforward algebra. The algorithm performs base change to $\operatorname{Spec}(\mathbb Z/M\mathbb Z)$ and evaluates in the full fiber algebra, including all of its components and nilpotents. It does not seek a rational point of that fiber. The scalar-output condition makes the resulting section descend to the base ring, where the magnitude bound identifies its unique integer lift. EGA IV4, printed p.111,18.2.1--18.2.2 gives the finite-algebra trace dictionary and base-change formula; the computational proof itself is the explicit multiplication table above. A weight or purity assertion without this small presentation does not supply the simulation, and no GCT orbit-closure obstruction is constructed.
\end{geometricstatement}
\begin{proof} The displayed dictionary transports the proof of NN-C020. Any additional realization, effectiveness or uniformity hypothesis is retained explicitly.\end{proof}


\section{NN-C021: geometric formulation and dictionary}
\begin{geometricstatement}[NN-C021]
For every hard contradiction, the Boolean-coordinate equation scheme is empty; its finite algebra has rank and Frobenius trace zero. E computes this vanishing efficiently relative to the former exhaustive exponent. Resolution width and polynomial-calculus degree/monomial size concern the cost of deriving the unit ideal from the chosen generators in specified syntactic systems. Thus the same zero scheme can have a short computation transcript and a long ideal-derivation certificate. The finite trace dictionary is SGA7 II XXII Corollary3.2, but that identity places no bound on the chosen presentation's derivation complexity. This is an effective distinction between a geometric zero object and an explicit proof of emptiness in a fixed calculus, not a new failure of purity or a representation-theoretic multiplicity statement.
\end{geometricstatement}
\begin{proof} The displayed dictionary transports the proof of NN-C021. Any additional realization, effectiveness or uniformity hypothesis is retained explicitly.\end{proof}


\section{NN-C022: geometric formulation and dictionary}
\begin{geometricstatement}[NN-C022]
Let $X_n=\operatorname{Spec}\mathbb Z[x_1,\ldots,x_n]/(x_j^2-x_j)$, the disjoint union of $2^n$ copies of $\operatorname{Spec}\mathbb Z$. Each Boolean output table has a unique integer multilinear interpolation which is an idempotent in this algebra, so $F$ defines a morphism $X_n\to X_m$. Over a Boolean target section $y$, its fiber is the finite split scheme with coordinate algebra $\mathbb Z^{|F^{-1}(y)|}$. The local clauses give a concise presentation of this same fiber. Its finite-locally-free rank, or trace of the identity, is the inverse count; reduction to $\mathbb F_2$ gives only the parity of that rank. E computes the rank without expanding the exponential product algebra. This is the explicit multiplication-matrix trace of EGA IV4,18.2.1--18.2.2, and the finite trace formalism of SGA4 XVII6.2.3. The image distinguisher compares support in the finite target with its uniform counting measure; the image has at most the source's number of components. No zeta-function estimate, cohomological purity theorem, or orbit-closure representation is needed, and A or M does not alter this parameter calculation.
\end{geometricstatement}
\begin{proof} The displayed dictionary transports the proof of NN-C022. Any additional realization, effectiveness or uniformity hypothesis is retained explicitly.\end{proof}


\section{NN-C023: geometric formulation and dictionary}
\begin{geometricstatement}[NN-C023]
The construction is already a scheme-theoretic theorem: parity-3-SAT is encoded in the absolute-Frobenius endomorphism of an explicitly trivialized rank-one coherent cohomology space of a projective hypersurface. Its sheaf-theoretic exact sequence and Cech basis prove the dictionary directly. Katz's original pp.35--36 locate the corresponding smooth Hasse--Witt formula, and SGA7 II XXII3.1--3.2 locate the zeta congruence and trace-point-count identity for arbitrary proper schemes. These reducible hypersurfaces are not asserted to furnish pure l-adic realizations of canonical bases, and their one-dimensional coherent Frobenius space is not a GCT representation multiplicity. The theorem identifies an effective arithmetic operation inside the geometric formalism, rather than importing Boolean hardness merely by analogy.
\end{geometricstatement}
\begin{proof} The displayed dictionary transports the proof of NN-C023. Any additional realization, effectiveness or uniformity hypothesis is retained explicitly.\end{proof}


\section{NN-C024: geometric formulation and dictionary}
\begin{geometricstatement}[NN-C024]
Let $U=\operatorname{Spec}\mathbb F_2[u,\Delta(u)^{-1}]$. The pencil defines a smooth projective hypersurface family over $U$, with a canonically based coherent line $R^{r-2}\pi_*\mathcal O$; its linearized relative Frobenius is multiplication by $uc$. The open smooth family is established by an explicit homogeneous Jacobian minor, which is the local differential criterion of SGA1 II4.10--4.12. The finite-field trace and full zeta congruence are exactly SGA7 II XXII3.1--3.2. Every good closed fiber has pure l-adic cohomology by Deligne, Weil I1.6; purity constrains the complex absolute values of Frobenius eigenvalues and does not compute this coherent characteristic-two invariant. The construction supplies genuine finite Weil objects for these counting gadgets, not canonical realizations of GCT's nonstandard basis or its structure constants. It sharpens an effectiveness gate while leaving that separate functorial/canonical realization problem intact.
\end{geometricstatement}
\begin{proof} The displayed dictionary transports the proof of NN-C024. Any additional realization, effectiveness or uniformity hypothesis is retained explicitly.\end{proof}


\section{NN-C025: geometric formulation and dictionary}
\begin{geometricstatement}[NN-C025]
For every fixed-width CNF $\Phi$ on $n\geq1$ Boolean variables there is an explicitly constructed polynomial $f_\Phi\in\mathbb F_2[x_1,\ldots,x_n]$, represented by a polynomial-size arithmetic circuit and having odd degree, such that
\[
 U_\Phi=\operatorname{Spec}\mathbb F_2[x,t]/(t^2+t-f_\Phi(x))
\]
is smooth and geometrically integral of dimension $n$, finite etale of degree two over $\mathbb A^n$, and satisfies $\#U_\Phi(\mathbb F_2)=2\#\operatorname{SAT}(\Phi)$. Its graph presentation is a polynomial-size complete intersection of quadratic equations with a displayed triangular Jacobian minor equal to one. In particular its count is always zero modulo two, whereas its count modulo four is twice the Boolean parity. The proof above constructs $f_\Phi=1+S_\Phi+x_1^d+x_1$ with $d$ odd and larger than the verifier degree; the odd pole at infinity proves that $f_\Phi$ is not $h^2+h$ over the algebraic closure. Thus this geometric integrality claim has an explicit certificate argument.

The map is an Artin--Schreier torsor under the constant group $\mathbb Z/2\mathbb Z$. Over a satisfying rational input its fiber is $\operatorname{Spec}(\mathbb F_2\times\mathbb F_2)$; over a nonsatisfying one it is $\operatorname{Spec}\mathbb F_4$. These fibers have the same geometric degree, but respectively two and zero rational points. The multiplication trace of the unit in either rank-two algebra is zero in characteristic two; its algebraic trace locator is EGA IV4,18.2.1--18.2.2, printedp111. SGA4 XVII6.2.3 is the corresponding sheaf-theoretic finite trace, a distinct formalism. The whole affine point count is the alternating compact-support etale trace, by Deligne Weil I(1.5.1); it is not the trace on infinite-dimensional ordinary affine coherent $H^0$. Its integer residue modulo four is a separate operation, not a map from an odd-adic coefficient field to $\mathbb Z/4\mathbb Z$. Exact point counting on these circuit-presented smooth objects is counting-complete by division by two, while E evaluates the constructed objects through the original $n$-variable formula. No canonical GCT multiplicity or positive polytope formula is identified.
\end{geometricstatement}
\begin{proof} The displayed dictionary transports the proof of NN-C025. Any additional realization, effectiveness or uniformity hypothesis is retained explicitly.\end{proof}


\section{NN-C026: geometric formulation and dictionary}
\begin{geometricstatement}[NN-C026]
For each fixed-width CNF $\Phi$ on $n\geq1$ variables, the proof explicitly constructs a finite field $K=\mathbb F_q$, $q=2^b$ with $b$ polynomial in the input size, and a monic odd-degree circuit-presented polynomial $f_\Phi(X)\in K[X]$. Put $d=\deg f_\Phi=2g+1$ and $w=g+1$. Its weighted projective completion
\[
 C_\Phi:\quad Y^2+YZ^w=Z^{d+1}f_\Phi(X/Z)
       \subset\mathbb P(1,w,1)
\]
is a smooth geometrically integral projective curve of genus $g$, with one rational point at infinity and
\[
 \#C_\Phi(K)=1+2\#\operatorname{SAT}(\Phi).
\]
The polynomial and the integers $d,g$ have polynomial-size circuit/binary descriptions; $d$ and $g$ themselves may be exponential. The two affine charts in the proof establish smoothness, completion and the Cech basis $e_j=[x^{-j}y]$, $1\leq j\leq g$, for $H^1(C_\Phi,\mathcal O)$. Absolute Frobenius sends $e_j$ to $e_{2j}$, with out-of-range indices zero; $q$-Frobenius sends it to $e_{qj}$. These coherent operators are nilpotent, independently of the encoded satisfying count.

For every prime $\ell\ne2$, the actual $H^1_{\rm et}(C_{\overline K},\mathbb Q_\ell)$ is pure of weight one and has Frobenius trace $q-2\#\operatorname{SAT}(\Phi)$. SGA7 XXII3.1 explains the common full zeta reduction modulo two; Deligne Weil I1.5.1,1.6 and2.3--2.5 supply the integral trace and purity interfaces. This proves a precise contrast between the characteristic-two coherent realization and the next binary digit of the integer etale trace. The divided count $(\#C(K)-1)/2$ is a counting-complete function on the succinct class, and its parity is parity-counting-complete. E evaluates the constructed instances from their original Boolean witnesses; it does not give a polynomial algorithm on arbitrary high-degree circuit-presented curves. No plethysm/Kronecker multiplicity identification or canonical-basis construction is claimed.
\end{geometricstatement}
\begin{proof} The displayed dictionary transports the proof of NN-C026. Any additional realization, effectiveness or uniformity hypothesis is retained explicitly.\end{proof}


\section{NN-C027: geometric formulation and dictionary}
\begin{geometricstatement}[NN-C027]
Let $f_1,\ldots,f_h\in\mathbb Z[x_1,\ldots,x_v]$ be the sparse quadratic Boolean and verifier equations in the proof, so that $v,h\geq1$, $f_1=x_1^2-x_1$, and $Y=V(f_1,\ldots,f_h)\simeq\coprod_N\operatorname{Spec}\mathbb Z$. Then the sparse cubic arithmetic hypersurface
\[
 X=V\left(t^2-t-\sum_{i=1}^h y_i f_i(x)\right)
             \subset\mathbb A^{v+h+1}_\mathbb Z
\]
is smooth of relative dimension $v+h$ with geometrically integral fibers. Its class is $[X]=\mathbb L^{v+h}+N\mathbb L^h$. For every field $K$ and every invertible prime $\ell$, its actual geometric compact-support groups, with their Galois actions, are
\[
 H_c^{2h}(X_{\overline K},\mathbb Q_\ell)=\mathbb Q_\ell(-h)^N,
 \qquad
 H_c^{2(v+h)}(X_{\overline K},\mathbb Q_\ell)=\mathbb Q_\ell(-v-h),
\]
and all other groups vanish. When $K$ is finite the displayed identifications are Frobenius-equivariant; they do not require finite $K$ to hold as Galois representations. Its arithmetic zeta is
$\zeta_X(s)=\zeta(s-v-h)\zeta(s-h)^N$.
These are separate proved assertions: affine-bundle stratification proves the scissor identity, while !-base change, the global trace and the involution $t\mapsto1-t$ prove the actual sheaf decomposition. Their exact interfaces are SGA4 XVII5.2.6/5.2.8 and XVIII2.8.1/2.9. The relative Jacobian identity proves smoothness over the full integer base, using SGA1 II4.10(iii)/4.12.

Over $B=\mathbb Z[1/2]$, replace the equation by $t^2+1=\sum_i y_i f_i(x)$. The resulting $X^-$ is also smooth with geometrically integral fibers. Its lower compact-support group is $(\chi_{-4}\otimes\mathbb Q_\ell(-h))^N$, its upper group is the same Tate line, and
\[
 \zeta_{X^-}(s)=\zeta^{(2)}(s-v-h)L(s-h,\chi_{-4})^N.
\]
Here the sign is the actual anti-invariant summand of the regular permutation representation $1\oplus\chi_{-4}$ of $\operatorname{Spec}\mathbb Z[i,1/2]\to\operatorname{Spec}B$. The first construction has a separate algebraic involution; the second realizes a quadratic Galois sign. Neither asserts an injective realization of the full split-zero semiring.

The proof constructs these data from a Boolean formula with $N$ satisfying assignments, with $v,h$ and the cubic description polynomial in that formula's size. Consequently extracting the lower cohomology multiplicity $N$ from this succinct geometric input is counting-complete, its parity is parity-counting-complete, and its nonvanishing is NP-complete. The corresponding shifted nonreal nontrivial zeta-zero orders are $N$ times a fixed nonzero analytic multiplicity; this identity is unconditional, while A locates the shifted critical lines. E evaluates $N$ with its original-witness bound. Thus actual purity, a concrete sign representation, smooth integral geometry and A-controlled analytic zero locations all coexist with this precise multiplicity-extraction problem. No unconditional time lower bound or effective GCT canonical-object identification is inferred.
\end{geometricstatement}
\begin{proof} The displayed dictionary transports the proof of NN-C027. Any additional realization, effectiveness or uniformity hypothesis is retained explicitly.\end{proof}


\section{NN-C028: geometric formulation and dictionary}
\begin{geometricstatement}[NN-C028]
The algebraic-geometric assertion is that, for the explicitly padded Boolean systems $Y_j\simeq\coprod_{N_j}\operatorname{Spec}\mathbb Z$, the two cubic schemes
\[
 X_j=V\left(t^2-t-\sum_{i=1}^H y_i f_{j,i}(x)\right)
          \subset\mathbb A^{V+H+1}_\mathbb Z
\]
are smooth with geometrically integral fibers, and their actual lower compact-support realizations are $\mathbb Q_\ell(-H)^{N_j}$. In the category of rational l-adic Frobenius representations, an injection between these realizations exists if and only if $N_1\leq N_2$. This equivalence holds for every prime with identical geometric parameters, and for the anti-invariant quadratic-character realizations at every odd prime. The scheme construction, relative trace and projectors are those proved in NN-C027, with exact SGA4 XVII5.2.6/5.2.8 and XVIII2.8.1/2.9 interfaces. Parsimonious padding proves that the geometric dictionary preserves the relevant counts and common Tate twist.

Thus even on this actual pure Tate subcategory, a succinct defining-equation presentation can make the existence of a subrepresentation PP-complete. The explicit matrices, once supplied, would be scalar and their rank comparison elementary; obtaining their multiplicities from the equations is the hard operation. Their zeta factors have A-controlled shifted critical lines, but the orders at a fixed nonreal lower-factor zero give precisely the same comparison. This is a theorem about linear injections of cohomology realizations, not about scheme morphisms, algebraic correspondences or orbit-closure containment. No fullness statement for a realization functor is assumed.
\end{geometricstatement}
\begin{proof} The displayed dictionary transports the proof of NN-C028. Any additional realization, effectiveness or uniformity hypothesis is retained explicitly.\end{proof}


\section{NN-C029: geometric formulation and dictionary}
\begin{geometricstatement}[NN-C029]
The algebraic-geometric theorem is constructive in the following precise representation. For the split Boolean coordinate algebra
\[
 B=\mathbb Z[x_1,\ldots,x_v]/(x_j^2-x_j)
                 \simeq\prod_{b\in\{0,1\}^v}\mathbb Z,
\]
let $Y=V(f_1,\ldots,f_h)$ be the closed subscheme cut out by the local verifier equations. If $Y$ is empty, E produces rational arithmetic circuits representing integral polynomials $a_i$ of degree at most $v$ whose ordinary ideal combination is one in $\mathbb Z[x]$. The first-violation construction partitions the component idempotents of $B$; the constant-scope division above lifts that partition to the ambient affine coordinate ring. It thereby constructs an integral polynomial isomorphism
\[
 V\left(t^2-t-\sum_i y_i f_i(x)\right)
        \longrightarrow\mathbb A^{v+h}_{\mathbb Z},
 \quad (x,y,t)\longmapsto(x,y+a(2t-t^2)),
\]
with the explicitly proved inverse and degree bounds. Its circuit description has the displayed $S(v)$ bound. Moreover that description is defined over the explicit open arithmetic base $\operatorname{Spec}\mathbb Z[1/B(v)!]$: the Hasse-jet Chinese-remainder isomorphism proves the interpolation denominators are invertible there. Thus the same effective presentation of the ideal identity and affine trivialization base-changes to every field of characteristic $p>B(v)$. The integral morphisms exist at all primes; only their claimed arithmetic-circuit presentation has this restricted base.

The symbolic preliminary multipliers are finite-algebra traces of products of local weights, so C016's trace construction and base change apply, with exact classical locator EGA IV4, printedp111,18.2.1--18.2.2 and equation18.2.1.2. The passage from quotient identities to ordinary ideal multipliers is proved here by local interpolation, without invoking an unproved effectivity theorem. In the sparse-CNF case, the classical first-violation partition already gives the ordinary identity without the Boolean equations. These are integral geometric statements with localized arithmetic-circuit presentations; no efficient circuit presentation at every prime, motivic equivalence theorem, or GCT representation-theoretic lower bound is inferred.
\end{geometricstatement}
\begin{proof} The displayed dictionary transports the proof of NN-C029. Any additional realization, effectiveness or uniformity hypothesis is retained explicitly.\end{proof}


\section{NN-C030: geometric formulation and dictionary}
\begin{geometricstatement}[NN-C030]
Let $P\in\mathbb F_2[z_{e,0},z_{e,1}]$ be nonzero of degree two in each of $M$ pairs, and put $r=2M+1$, $F=z_0P$. In the Fermat pencil $H+uF$, every singular geometric fiber occurs at a parameter in the one finite field $\mathbb F_{2^{r\operatorname{ord}_r(2)}}$. The proof is the pairwise Euler identity, normalization at $z_0=1$, and the resulting root-of-unity and power equations above. A parameter of absolute degree greater than that bounded field degree therefore gives a smooth projective fiber; a quadratic Artin--Schreier tower supplies such a parameter with degree $O(r^2)$ and a polynomial-length presentation.

For the incidence-rail polynomials, the proper smooth fiber has $H^0(\mathcal O)=\mathbb F_{2^b}$ and one other nonzero coherent group, a canonically based line in degree $r-2$. Absolute Frobenius has semilinear coefficient $\theta c$, and the linearized field Frobenius has scalar $c$, the original parity count. The complete Cech proof is NN-C023 and the classical smooth formula is Katz2.3.7.17, p36. SGA1 II4.10(iii)/4.12 is the smoothness interface, SGA7 II XXII3.1--3.2 gives the proper zeta congruence, and Deligne Weil I1.6 gives purity of the actual l-adic realizations. These provide an exact deterministic finite-field geometric dictionary for the hard coefficient. A purity or existence theorem does not evaluate it, and the construction supplies no canonical identification with GCT multiplicities.
\end{geometricstatement}
\begin{proof} The displayed dictionary transports the proof of NN-C030. Any additional realization, effectiveness or uniformity hypothesis is retained explicitly.\end{proof}


\section{NN-C031: geometric formulation and dictionary}
\begin{geometricstatement}[NN-C031]
For every odd $k\geq3$, every nonzero pairwise-degree-two $P$ over $\mathbb F_2$ in $M=2^{k-1}-1$ pairs defines the smooth projective hypersurface
\[
 V\left(z_0^{2^k-1}+\sum_e(z_{e,0}^{2^k-1}+z_{e,1}^{2^k-1})
                        +\omega z_0P\right)/\mathbb F_4.
\]
The singular-parameter field from NN-C030 has odd degree $k(2^k-1)$ and therefore excludes the fixed quadratic parameter $\omega$. Its top coherent line is explicitly trivialized by the boundary/Cech construction. Absolute Frobenius has semilinear coefficient $\omega c$; $4$-power Frobenius has scalar $c$. The lower coherent groups are only the degree-zero field, and the proper zeta congruence is $(1-cT)/(1-T)$ modulo two. Exact classical interfaces are Katz2.3.7.17, SGA1 II4.10(iii)/4.12, SGA7 II XXII3.1--3.2 and Deligne Weil I1.6.

The forced-unit-clause padding turns any source formula into this geometric class while preserving its count exactly. Thus the field, smoothness and coherent basis are all supplied effectively, and the sole unknown scalar is parity-counting complete. This is an actual smooth projective finite-field construction, with pure l-adic cohomology, and a precise computational operation inside it. It does not assert a canonical representation-theoretic realization of permanent/determinant multiplicities.
\end{geometricstatement}
\begin{proof} The displayed dictionary transports the proof of NN-C031. Any additional realization, effectiveness or uniformity hypothesis is retained explicitly.\end{proof}


\section{NN-C032: geometric formulation and dictionary}
\begin{geometricstatement}[NN-C032]
Let $X=V(H+\omega F)/\mathbb F_4$ be a hypersurface of NN-C031. The displayed Hermitian-matrix equations define a projective $\mathbb F_2$ scheme $W$ whose base change is the Segre product $X\times X^\sigma$ and whose descent datum is factor exchange. Thus $W$ represents the Weil restriction. It is smooth and geometrically integral, with coherent cohomology only in degrees $0,d,2d$, of ranks $1,2,1$. In the explicit descended bases $1$, $(e+e',\omega e+\omega^2e')$, and $e\cup e'$, its Frobenius operators are $1$, $c\left(\begin{smallmatrix}0&1\\1&0\end{smallmatrix}\right)$, and $c$. Its exact rational-point bijection with $X(\mathbb F_4)$ and its proper zeta congruence are proved above.

The scheme-theoretic interfaces are the explicit Segre chart proof with Stacks27.13.6, restriction-of-scalars sheaf/base-change Lemmas97.11.1--97.11.2, smoothness descent Lemma35.23.29, coherent flat base change Lemma30.5.2 and coherent Kunneth Lemma33.29.1. SGA7 II XXII3.1--3.2 supplies the coherent zeta congruence and Deligne Weil I1.6 supplies purity. These are exact locators, with no original Grothendieck archive cote claimed. The formula for the descended middle matrix is proved by semilinear arithmetic and cup-product compatibility, not guessed from point counts. The underlying count reduction and E evaluator are unchanged; the new construction removes the fixed-field obstruction while preserving an explicit small coherent realization.
\end{geometricstatement}
\begin{proof} The displayed dictionary transports the proof of NN-C032. Any additional realization, effectiveness or uniformity hypothesis is retained explicitly.\end{proof}


\section{NN-C033: geometric formulation and dictionary}
\begin{geometricstatement}[NN-C033]
For every prime $p$, let $B=\mathbb F_p[y]$ and let
\[
 \mathscr A=B[x_1,\ldots,x_n]/(x_i^2-x_i),\qquad
       a=\prod_f\Phi_f(x).
\]
This is a finite free split algebra of rank $2^n$, with the Boolean orthogonal idempotents as a basis. Its multiplication trace $\operatorname{Tr}_{\mathscr A/B}(a)$ is the universal partition polynomial, because multiplication is diagonal in that basis. The theorem constructs a prime-field arithmetic circuit for this relative trace in size $S(n)$, naturally under arbitrary substitutions of the parameters. The auxiliary scalar extension is finite separable and free; the proof computes the trace after that extension and recovers its original scalar coordinate by an explicit basis. This descent is an identity of regular functions on the full parameter scheme, not pointwise interpolation on its rational points.

For an unsatisfiable local verifier, the same construction supplies an ordinary ideal identity in $\mathbb F_p[x]$ and polynomial maps identifying its associated smooth cubic with affine space over every extension field of characteristic $p$. The maps and inverses are explicitly represented by subexponential circuits and have the degree bounds from NN-C029. Thus the characteristic-two geometric trivialization gate is also closed for this family. The exact classical finite-locally-free trace locator is EGA IV4,18.2.1--18.2.2, printed p111, especially base-change equation18.2.1.2. The integral-matrix identity, change of interpolation nodes, integer-fiber separation and coordinate descent are proved here; no additional Grothendieck archive locator is asserted. This family-level effective trace theorem supplies neither a canonical GCT multiplicity realization nor general polynomial-time geometric recognition.
\end{geometricstatement}
\begin{proof} The displayed dictionary transports the proof of NN-C033. Any additional realization, effectiveness or uniformity hypothesis is retained explicitly.\end{proof}


\section{NN-C034: geometric formulation and dictionary}
\begin{geometricstatement}[NN-C034]
Let $B=\mathbb Z[y]$ and
\[
 \mathscr A=B[x_1,\ldots,x_n]/(x_i^2-x_i),\qquad
                       a=\prod_f\Phi_f(x).
\]
The finite locally free multiplication trace $\operatorname{Tr}_{\mathscr A/B}(a)$ is the universal partition polynomial. The theorem supplies an integral arithmetic circuit for this regular function on $\operatorname{Spec}B$ in size $S(n)$, with a constant-free presentation of polynomial formal degree. The auxiliary extension $B\to A\otimes_{\mathbb Z}B$ is finite free and faithfully flat. It contains the required distinct unit-separated interpolation sections; these differences remain units after every base change, including ramified residue characteristics. Exact interpolation and coordinate projection descend the computed regular function to $B$.

Fix $C,a\ge1$.  Let $m\le Cn$ explicit polynomial lists
$\Phi_f=\sum_j y_{f,j}x^{\alpha(f,j)}$ have description
length $L$, declared scopes of size at most
$(\lceil\log_2n\rceil)^a$, and total degree budget
$\sum_f\max_j|\alpha(f,j)|\le Cn$.  Put $B=\mathbb Z[y]$ and
\[
 \mathscr A=B[x_1,\ldots,x_n]/(x_i^2-x_i),\qquad
                       u=\prod_f\Phi_f(x).
\]
The finite locally free algebra $\mathscr A$ has rank $2^n$.
Its multiplication trace
$\operatorname{Tr}_{\mathscr A/B}(u)=Z_\Phi(y)$, as is
seen in the orthogonal Boolean-idempotent basis.  This
regular function on $\operatorname{Spec}B$ has an integral
circuit of size and bit construction time
$\operatorname{poly}(n,L)2^{O_{C,a}(n(\log\log n)^2/\sqrt{\log n})}$,
with polynomial-bit integer constants and a constant-free
compilation of formal degree at most $\max(m,1)$, with $n$
the original number of Boolean coordinates.  For every map
$B\to R$ the same circuit computes the base-changed
multiplication trace over $R$; coordinate descent from
the finite free cyclotomic extension proves the identity
before any specialization.  The exact classical trace
and base-change interface is EGA IV4, 18.2.1--18.2.2,
printed p.~111, including equation 18.2.1.2.  This is an
effective theorem for these specified regular trace
functions, not a claimed polynomial-time algorithm for
general finite algebra traces or GCT multiplicities.

For a support $\Omega\subseteq[r]\times[r]$ satisfying
$|\Omega|\le Cr(\lceil\log_2r\rceil)^a$, put
$B=\mathbb Z[a_{ij}:(i,j)\in\Omega]$ and
$\mathscr A=B[x_1,\ldots,x_r]/(x_j^2-x_j)$.  The regular
function
\[
 \operatorname{Tr}_{\mathscr A/B}\left(
     (-1)^r\prod_j(1-2x_j)
                 \prod_i\sum_{j:(i,j)\in\Omega}a_{ij}x_j\right)
                              =\operatorname{per}(A)
\]
has an integral constant-free circuit of size
$\operatorname{poly}(r)2^{O_{C,a}(r(\log\log r)^2/\sqrt{\log r})}$
and formal degree at most $2r$, valid after every base
change.  Ryser's identity identifies this signed trace,
while the heavy-row decomposition constructs it using
smaller trace functions with polylogarithmic row scopes.

For the split torus with characters indexed by $\Omega$,
let $\mathcal V_\Omega$ have one basis vector of character
$\prod_i a_{i,\sigma(i)}$ per allowed perfect matching.
Its effective character is the same permanent polynomial.
The restriction partition above gives the actual direct-sum
decomposition
\[
 \mathcal V_\Omega\simeq
  \bigoplus_{\pi}
    \chi_\pi\otimes\mathcal V_{\Omega_\pi},
 \qquad \chi_\pi=\prod_{i\in H}a_{i,\pi(i)},
\]
where $\Omega_\pi$ is the remaining minor support and its
characters are included in the original torus.  Applying
the integral minor circuits to these summands gives the
stated character circuit; positivity of the decomposition
does not remove the virtual differences inside the minor
circuits.  The precise classical trace and character-ring
interfaces remain EGA IV4, 18.2.1--18.2.2, printed p.~111,
and the SGA3/SGA6 locators already stated in NN-S019.

For the unsatisfiable local-verifier scheme, the same circuits give an ordinary global ideal identity $\sum_i a_if_i=1$ over $\mathbb Z$ and the explicit mutually inverse maps identifying its smooth counting cubic with affine space. The maps have the actual degree bounds from NN-C029, integral coefficient circuits of size $S(v)$, and polynomial formal degree after removing all large constant leaves. Thus these are effective integral geometric trivializations, valid after arbitrary base change, with no denominator-prime exclusion. On the empty Boolean equation schemes of NN-C021, these circuit ideal certificates over $\mathbb F_2$ have size $2^{o(n)}$, while deriving their unit ideal by ordinary monomial polynomial calculus requires size $2^{\Omega(n)}$. This compares specified presentations of the unit-ideal proof, not different underlying zero schemes or different cohomological degrees. The exact classical multiplication-trace and base-change locator is EGA IV4,18.2.1--18.2.2, printed p111, equation18.2.1.2. All node, inverse, flat scalar-extension and polynomial-identity calculations are proved above; no unverified Grothendieck archive cote or GCT canonical-basis functor is invoked.
\end{geometricstatement}
\begin{proof} The displayed dictionary transports the proof of NN-C034. Any additional realization, effectiveness or uniformity hypothesis is retained explicitly.\end{proof}


\section{NN-C035: geometric formulation and dictionary}
\begin{geometricstatement}[NN-C035]
Let $Y_F$ be the finite etale subscheme of $\operatorname{Spec}(\mathbb F_2[x_1,\ldots,x_n]/(x_i^2-x_i))$ cut out by the clause-violation equations. Put $B=\mathbb F_2[z_{i,b}]$ and $\mathscr A_F=\mathcal O(Y_F)\otimes_{\mathbb F_2}B$. It is finite locally free over $B$, with one rank-one summand for every satisfying assignment. For
\[
                           w=\prod_i\bigl((1-x_i)z_{i,0}+x_i z_{i,1}\bigr),
\]
the multiplication trace $\operatorname{Tr}_{\mathscr A_F/B}(w)$ is $A_F(z)$ in the component-idempotent basis. For fixed clause width, this regular function has a deterministically constructible constant-free circuit of size $(n+m)^{O_k(1)}2^{O_k(n(\log\log n)^2/\sqrt{\log n})}$ and formal degree $O_k(n)$. This statement remains true after every base change to a characteristic-two coefficient ring. The same trace theorem holds for the explicitly listed weight factors with polylogarithmic scopes and linear total input degree stated above.

The compressed equations define the same closed subscheme inside the Boolean scheme, because its product coordinate algebra has ideals determined by their vanishing components. The sparsification leaves form a disjoint union decomposition of $Y_F$, and multiplication trace adds over that decomposition. The nonzero two-rail coefficient records an actual component; trace at all weights equal to one is only its rank modulo two. Every point of $Y_F$ is $\mathbb F_2$-rational, so Frobenius on $H^0_{\mathrm{et}}(Y_F\times\overline{\mathbb F}_2,\mathbb Q_\ell)$, for $\ell\ne2$, is the identity. This weight-zero purity does not compute the symbolic trace. The exact classical multiplication-trace and base-change locator is EGA IV4,18.2.1--18.2.2, printed p111, equation18.2.1.2. No reductive-group multiplicity or canonical GCT object is asserted.
\end{geometricstatement}
\begin{proof} The displayed dictionary transports the proof of NN-C035. Any additional realization, effectiveness or uniformity hypothesis is retained explicitly.\end{proof}


\section{NN-C036: geometric formulation and dictionary}
\begin{geometricstatement}[NN-C036]
Let $F$ be a fixed-width CNF on $n$ declared Boolean coordinates. Its integral Boolean scheme is
\[
 Y_F=\operatorname{Spec}\!\left(
 \mathbb Z[x_1,\ldots,x_n]/(x_i^2-x_i,\ 1-c_C(x):C\in F)\right).
\]
Here $c_C=1-\prod_{\ell\in C}(1-\ell(x))$ is the integral clause-satisfaction indicator. The scheme is the disjoint union of one $\operatorname{Spec}\mathbb Z$ for each satisfying assignment. Put $B=\mathbb Z[z_{i,b}]$, $\mathscr A_F=\mathcal O(Y_F)\otimes_{\mathbb Z}B$, and
\[
                         w=\prod_i((1-x_i)z_{i,0}+x_i z_{i,1}).
\]
Then the finite locally free multiplication trace $\operatorname{Tr}_{\mathscr A_F/B}(w)$ equals the integral assignment polynomial $A_F$. For fixed width $k$, this regular trace function has a deterministically constructible integral constant-free circuit of size $(n+m)^{O_k(1)}2^{O_k(n(\log\log n)^{1+1/k}/(\log n)^{1/(2k)})}$ and polynomial formal degree. The same circuit computes the trace after every base change, including all characteristics and nonreduced coefficient rings. Explicit weight factors with polylogarithmic scopes and linear total input degree have the same bound with their description length and fixed parameters included.

Disjoint sparsification is a decomposition of $Y_F$ into clopen finite etale subschemes over $\mathbb Z$, and the trace adds over it. The construction retains the actual integer component multiplicities, rather than only their residues. Equivalently, the torus with one character $z_{i,b}$ per literal rail has an effective module with one basis vector of character $\prod_i z_{i,x_i}$ for each satisfying assignment; its character is computed by the integral circuit using virtual differences. For a graph, this yields the effective independent-subset character after setting the zero rails to one. The exact trace/base-change source is EGA IV4,18.2.1--18.2.2, printed p111, equation18.2.1.2; the split-torus character and Grothendieck-ring locators are those in NN-S019. This is an effective statement for these presented finite schemes and characters, not a theorem about general GCT multiplicities.

For an unsatisfiable normalized width-$k$ CNF with clause
polynomials $v_1,\ldots,v_m$, the closed subscheme
$V(v_1,\ldots,v_m)\subseteq\mathbb A^n_{\mathbb Z}$ is empty,
and its unit ideal has the displayed explicit integral
circuit certificate of size
$(n+m)^{O_k(1)}2^{O_k(n(\log\log n)^{1+1/k}/(\log n)^{1/(2k)})}$.
No Boolean equations occur in that ideal.  The multipliers
have degree at most $n-|S_j|$ and integral constant-free
circuits of polynomial formal degree.

For the affine hypersurface
\[
 X_F=\{t^2-t=\sum_{j=1}^m y_jv_j(x)\}
                          \subseteq\mathbb A^{n+m+1}_{\mathbb Z},
\]
the same certificate gives the explicit isomorphism
$X_F\simeq\mathbb A^{n+m}_{\mathbb Z}$ by the unit-ideal
coordinate formulas already proved in NN-C029:
\[
 z_j=y_j+h_j(x)(2t-t^2),\qquad
 t=\sum_jv_j(x)z_j,\qquad
 y_j=z_j+h_j(x)(t^2-2t).
\]
The first formula is the map from $X_F$ to coordinates
$(x,z)$, and the last two define its inverse.  The equality
$\sum_jv_jh_j=1$ makes $\sum_jv_jz_j=t$ on $X_F$ and
$\sum_jv_jy_j=t^2-t$ on the inverse image; substitution
then gives both identity composites.  These calculations
use no division and hold over every base ring.  The forward
map has degree at most $n+2$ and the inverse has degree at
most $n+2k+2$, since $\deg(\sum_jv_jz_j)\le k+1$.
Sharing the common sum and multiplier circuits gives both
maps the same mildly subexponential original-$n$ circuit
bound and polynomial formal degree.  This is an effective
integral trivialization for the stated empty clause systems,
not a general recognition theorem for affine space.
\end{geometricstatement}
\begin{proof} The displayed dictionary transports the proof of NN-C036. Any additional realization, effectiveness or uniformity hypothesis is retained explicitly.\end{proof}


\section{NN-C037: geometric formulation and dictionary}
\begin{geometricstatement}[NN-C037]
For any field $K$, let $\varphi:\mathcal O_{\mathbb P^1_K}(-R)^w\to\mathcal O_{\mathbb P^1_K}^{w}$ be the vector-bundle map obtained by homogenizing a $w\times w$ polynomial matrix $A(T)$ of degree at most $R$, with $A(0)=I_w$. Its top exterior power is explicitly the determinant section
\[
                   s=\det(\varphi)\in
                     H^0(\mathbb P^1_K,\mathcal O(wR)).
\]
It is nonzero at $[1:0]$. If $\varphi$ has rank at most $k_t$ at the distinct rational points $[1:\alpha_t]$, its effective Cartier zero divisor contains
\[
                        \sum_t(w-k_t)[1:\alpha_t],
                         \qquad \sum_t(w-k_t)\le wR.
\]
This is a standalone geometric form of the obstruction. On the affine chart the local parameter is $T-\alpha_t$; constant row and column reduction makes $w-k_t$ rows divisible by this parameter, proving the claimed local divisor multiplicity. The homogeneous determinant has total degree $wR$; after base change to an algebraic closure, its affine zero multiplicities plus its possible multiplicity at infinity sum to $wR$. Thus the inequality is proved explicitly, without assuming an unverified determinant-functor formula.

For an E factor of scope $I$, restriction to its $w$ coordinate columns gives this map, with $k_t\le\min(w,a_t+c_t)$ because its image at the node factors through the assigned coordinate group and the other groups' measurement quotients. Equation~\eqref{eq:C037-scope-bound} is exactly the resulting divisor-degree inequality. For a rationally coefficient-dependent universal linear factor, apply the same theorem over the generic coefficient field $K(u)$ to the Jacobian matrix proved above. It expresses a genuine obstruction to these families of low-degree vector-bundle degenerations, not a statement that the permanent itself is a determinant.

The exact Grothendieck source interface is EGA~IV4, 21.1.4, printed p.~256, for the divisor of an invertible-sheaf meromorphic section; 21.1.7 and equation~21.1.7.1, p.~257, for effectivity of a regular section; and Proposition~21.2.11 with its proof, p.~262, for the line-section/divisor correspondence. The present local multiplicity and degree calculation supplies the additional quantitative assertion needed here.

Finally, the dense symbolic permanent is still the finite locally free multiplication trace of the displayed Ryser weight in
$\mathbb Z[a_{ij}][x_1,\ldots,x_n]/(x_j^2-x_j)$, up to the scalar $(-1)^n$. EGA~IV4,18.2.1--18.2.2, p.~111 and equation~18.2.1.2 give its trace and base-change interface. The obstruction says that this compactly presented regular trace function cannot acquire NN-C034's useful separated layout merely from its linear factor degrees. It imposes no obstruction on the existence of that finite etale algebra or its trace, and no lower bound on all algorithms for the regular function.\paragraph{Geometric form of the Boolean-summary obstruction.}
Over the generic coefficient field $F$, homogenize the
$I$-column Jacobian minor to obtain
\[
 \varphi:\mathcal O_{\mathbb P^1_F}(-R)^w
               \longrightarrow\mathcal O_{\mathbb P^1_F}^{w},
 \qquad
       s=\det\varphi\in H^0(\mathbb P^1_F,\mathcal O(wR)).
\]
It is nonzero at zero. The partition of the Boolean cube
by the formal values of each deformed linear form
forces the fiber of this Jacobian map at $\alpha_t$ to
have rank at most $\lfloor\log_2N_t\rfloor+\delta_t$.
For moving rational nodes, the additional rank-one map
$v_tq_t$ is exactly the chain-rule correction. Thus the
effective zero divisor of $s$ contains
\[
 \sum_t
 \max(0,w-\lfloor\log_2N_t\rfloor-\delta_t)[1:\alpha_t],
\]
whose degree is bounded by $wR$. A summary with at most
$M_t$ states on a Zariski-dense set of coefficient
specializations bounds $N_t$ by $M_t$, regardless of
coefficient adaptivity of that summary. This is the precise
geometric restatement of equation~\eqref{eq:C037-summary-bound}.
The line-section/divisor conventions and effectivity are
EGA IV4, 21.1.4, 21.1.7 and Proposition21.2.11, printed
pp.~256,257,262; the quantitative multiplicities and the
Boolean fiber count are proved above. The statement concerns
this family of Jacobian degenerations over coefficient space.
It does not equate a finite summary partition with an
algebraic quotient scheme or assert any general obstruction
to existence of the Ryser trace function.


\paragraph{Geometric meaning of finite branching.}
Affine coefficient space is integral. Finitely many input
sets covering a nonempty open cannot all have proper
closed closures, as the polynomial-product proof shows.
On a dense leaf the numerical identities therefore define
rational identities over the generic coefficient field.
Its scope Jacobian minor yields the same map
$\mathcal O_{\mathbb P^1}(-R)^w\to\mathcal O_{\mathbb P^1}^w$,
and the effective zero divisor of its determinant must
contain the rank-defect divisor supplied by the generic
value count. Its degree is at most $wR$, giving the displayed
inequality. The exact divisor/section interface remains
EGA IV4,21.1.4,21.1.7 and Proposition21.2.11, printed
pp.~256,257,262. No quotient scheme is attached to a
discrete test or branch set, and no algebraic structure
of those sets was assumed.

\end{geometricstatement}
\begin{proof} The displayed dictionary transports the proof of NN-C037. Any additional realization, effectiveness or uniformity hypothesis is retained explicitly.\end{proof}


\section{NN-C038: geometric formulation and dictionary}
\begin{geometricstatement}[NN-C038]
Fix $C\ge1$, $m\le Cn$ explicit polynomial factors with independent coefficients, scope bound $w$ and degree budget $\sum_f\deg\Phi_f\le Cn$, with full list description length $L$. Put $B_0=\mathbb Z[y]$ and
\[
 \mathscr A=B_0[x_1,\ldots,x_n]/(x_i^2-x_i),\qquad
                         u=\prod_f\Phi_f(x).
\]
This finite locally free algebra has rank $2^n$, and its multiplication trace is the regular function
\[
               \operatorname{Tr}_{\mathscr A/B_0}(u)
                 =\sum_{x\in\{0,1\}^n}\prod_f\Phi_f(x)
                                     \quad\text{on }\operatorname{Spec}B_0.
\]
Define
\[
 \begin{aligned}
 \lambda&=\lceil\log_2n\rceil,\\
 B&=\lceil\log_2(8(C+1)(w+1)\lambda)\rceil,\\
 z&=\max(2,\lambda/B),\\
 \theta&=\min(1,\lceil\log_2z\rceil^{3/2}/\sqrt z).
 \end{aligned}
\]
Then
this regular trace function has an integral constant-free circuit of size and deterministic construction time
$\operatorname{poly}_C(n,L)2^{O_C(n\theta)}$ and formal coefficient degree at most $\max(m,1)$. The identity is established before base change and therefore holds over every commutative ring. It is subexponential in the original number of Boolean coordinates whenever the scopes satisfy $w=n^{o(1)}$ and the explicit input length is subexponential. The scheme still has relative dimension zero; no claim of a small explicit idempotent basis is made.

For a bipartite support graph with $r$ vertices in each part and $S\le r^{1+o(1)}$ edges, let a split torus have one specified character for each allowed edge. The actual free torus module
\[
                   \mathcal V_\Omega
                    =\bigoplus_{\pi\ {\rm allowed\ bijection}}
                            \bigotimes_{i=1}^r L_{i,\pi(i)}
\]
has the symbolic permanent as its character. Its class has an integral expression with shared arithmetic subexpressions, allowing virtual differences, of size $2^{o(r)}$ and formal degree at most $2r$, with the full uniform bound given by $B_{\rm per}$ and $\theta_{\rm per}$ in the statement. The heavy-row expansion is an actual direct sum decomposition into selected-edge characters tensored with the corresponding minor modules. Ryser's formula also identifies its character with the multiplication trace of the stated signed weight on the finite Boolean algebra in the original $r$ coordinates.

The exact classical trace and base-change interface is EGA IV4,18.2.1--18.2.2, printed p111, including equation18.2.1.2. Split-torus gradings are SGA3, ExposeI,4.7.2--4.7.3, and direct sum/tensor product in the representation Grothendieck ring are SGA6, ExposeIV,(1.5.2) and Definition2.2, as verified for NN-S019. These statements concern compactly presented regular trace functions and specified effective torus characters; they do not claim arbitrary geometric pushforwards or general GCT multiplicities are polynomial-time computable.
\end{geometricstatement}
\begin{proof} The displayed dictionary transports the proof of NN-C038. Any additional realization, effectiveness or uniformity hypothesis is retained explicitly.\end{proof}


\section{NN-C039: geometric formulation and dictionary}
\begin{geometricstatement}[NN-C039]
For the compactly presented finite locally free Boolean
algebra
\[
 \mathscr A=\mathbb Z[y][x_1,\ldots,x_n]/(x_i^2-x_i),
                \qquad u=\prod_f\Phi_f(x),
\]
the coefficientwise monomial reduction above leaves $u$
unchanged. Its multiplication trace therefore has the
integral constant-free circuit bound
\[
        \operatorname{poly}_C(n,L)\,
           2^{O_C(n\min(1,(1+d/n)\theta(n,w)))}
\]
and formal coefficient degree at most $\max(m,1)$, where
$d$ is the sum of the maximum support sizes of the reduced
factor monomials and $\theta(n,w)$ is the explicitly
defined uniform parameter of NN-C038. This is an identity
of regular trace functions over $\mathbb Z[y]$ and survives
every commutative-ring base change, even when the original
binary exponents are very large. The trace and base-change
interface is EGA IV4,18.2.1--18.2.2, printed p.~111,
including equation18.2.1.2.

For full local tables, let the split torus have one
specified character $y_{f,a}$ for each local assignment.
The actual free module
\[
             \bigoplus_{x\in\{0,1\}^n}
                          \bigotimes_f L_{f,x|I_f}
\]
has that table partition polynomial as its character,
with multiplicities retained if different assignments
have the same character. For $m\le Cn$ and the growing
arity range above, its effective class has an integral
expression with virtual differences and shared arithmetic
subexpressions of size $2^{o(n)}$ and formal degree at
most $\max(m,1)$. For the CNF satisfaction specialization
with its added $n$ unary weights, the character belongs to
the finite etale satisfying-assignment scheme; its circuit
has the same time bound and formal degree at most $m+n$,
with $m$ the clause count. The exact torus and Grothendieck-ring dictionary
is SGA3, ExposeI,4.7.2--4.7.3 and SGA6, ExposeIV,(1.5.2),
Definition2.2. These specified circuit bounds do not assert
polynomial construction of a basis or of general
representation-theoretic multiplicities.

\end{geometricstatement}
\begin{proof} The displayed dictionary transports the proof of NN-C039. Any additional realization, effectiveness or uniformity hypothesis is retained explicitly.\end{proof}


\section{NN-C040: geometric formulation and dictionary}
\begin{geometricstatement}[NN-C040]
Let $F=K(u_1,\ldots,u_w)$ and let
\[
             Y=\operatorname{Spec}
             F[x_1,\ldots,x_w]/(x_j^2-x_j:1\le j\le w)
\]
be the finite etale Boolean scheme of rank $N=2^w$. For each $j$, the coefficient derivative $P_j(T)$ of the assigned family $L(T)\in\Gamma(Y,\mathcal O_Y)[T]$ is an arbitrary section of that algebra with $P_j(0)=x_j$ and $T$-degree at most $R$. There is no claim that the $P_j(T)$ remain idempotent or define an algebra endomorphism of $Y$.

On $\mathbb P^1_F$, squarefree multiplication of these sections and evaluation on the $N$ points of $Y$ define a morphism of vector bundles
\[
 \varphi:\bigoplus_{S\subseteq[w]}
                  \mathcal O_{\mathbb P^1}(-R|S|)
                 \longrightarrow\mathcal O_{\mathbb P^1}^{N}.
\]
Its affine-coordinate matrix is $V(T)^{\mathsf T}$, with column $S$ of degree at most $R|S|$; homogenization defines the indicated global map. At $T=0$ it is invertible. Its determinant is therefore a nonzero global section of
$\mathcal O_{\mathbb P^1}(Rw2^{w-1})$ and defines an effective divisor of that degree. At a rational node $\alpha_t$, coefficient differentiation confines the initial evaluation points of each formal-value class to one point if $\nabla\alpha_t=0$, or to one affine line otherwise. The explicit Hasse expansion in equation~\eqref{eq:C040-Hasse} bounds its order-$k$ coefficient-vector space by $H_w(k)$ or $A_w(k)$, respectively. Consequently this determinant divisor contains
\[
 \sum_t\left(\sum_{k=0}^{w-1}
            \max(0,2^w-M_tD_t(k))\right)[\alpha_t],
\]
with the stronger class-size multiplicities proved above. Comparison of effective-divisor degrees proves equation~\eqref{eq:C040-exact}, and subexponential value multiplicities force $R/Q\ge1-o(1)$.

The section/divisor interface is EGA IV4,21.1.4 p256,21.1.7 and equation21.1.7.1 p257, and Proposition21.2.11 with full proof p262, verified from the original scan. The bundle, its determinant degree and all local multiplicities have been constructed and proved here. The value partition is used only to constrain jets of sections; it is not asserted to define a geometric quotient, an algebra endomorphism, or a new representation multiplicity. This is an explicit geometric obstruction for low-degree families of regular functions on a finite etale scheme, with an honest limitation to that interpolation interface.

After the finite separable change of coefficient field $F\subset F'$,
replace $Y$ and $\mathbb P^1_F$ by $Y_{F'}$ and $\mathbb P^1_{F'}$.
The determinant section again belongs to
$\mathcal O_{\mathbb P^1_{F'}}(Rw2^{w-1})$, is nonzero at zero, and
has exactly the proved lower bounds on its divisor multiplicities at
the $F'$-rational interpolation nodes. Thus finite separable coefficient
extensions preserve this effective-divisor obstruction, including the
$R/Q\ge1-o(1)$ consequence. The extension of coefficient derivations is
the stated classical separability theorem, and the divisor interface
remains EGA IV4,21.1.4,21.1.7 and Proposition21.2.11 at the original
pages cited in NN-C040. This is a base change of the explicitly
constructed bundle calculation; it does not assert an efficient
representation of algebraic coefficients or a general circuit lower
bound.


The rational coefficient-derivative sections on
$Y_{E_0}\times\mathbb P^1_{E_0}$ give a meromorphic squarefree
evaluation matrix. Clearing its Boolean-component column $x$ by
$B_x^{2w}$ constructs a global bundle map
\[
 \mathcal O_{\mathbb P^1}(-2wR)^{2^w}
       \longrightarrow \mathcal O_{\mathbb P^1}^{2^w}.
\]
Its determinant is a nonzero section of
$\mathcal O_{\mathbb P^1}(2w2^wR)$ and has precisely the proved
lower bounds on its orders at the nodes, because every clearing
factor is a unit there. Its effective divisor therefore yields
\eqref{eq:C040-rational}. With a common denominator, multiplication
by $B(T)/B(0)$ reduces the value partitions to the earlier polynomial
bundle statement of degree $a w2^{w-1}$. The same verified EGA IV4
section/divisor locators apply. Thus allowing controlled poles does
not remove this effective-divisor obstruction; neither efficient
construction of the exponentially large matrix nor a geometric
quotient of the summary classes is asserted.

\end{geometricstatement}
\begin{proof} The displayed dictionary transports the proof of NN-C040. Any additional realization, effectiveness or uniformity hypothesis is retained explicitly.\end{proof}


\section{NN-C041: geometric formulation and dictionary}
\begin{geometricstatement}[NN-C041]
Over the universal coefficient ring
$R=\mathbb Z[S_{ij},U_{ia},V_{aj}]$, the explicitly constructed
algebra $B_R$ is finite free of rank $H^2$. Its spectrum is a
finite locally free, generally nonreduced thickening of the
coefficient base. Let
\[
 Y=\operatorname{Spec}
 B_R[x_1,\ldots,x_n]/(x_j^2-x_j:1\le j\le n).
\]
This Boolean scheme is finite etale of rank $2^n$ over
$\operatorname{Spec}B_R$. The sum inside
\eqref{eq:C041-ryser} is its multiplication trace of the
specified product section. EGA IV4,18.2.1--18.2.2, printed
p.~111, is the already checked finite-free trace and base-change
interface. The new integer algebra and the section have been
given explicitly, including all of their ranks and costs.

Give $t^\alpha z^{[\beta]}$ the torus character
$\alpha-\beta\in\mathbb Z^d$. The multiplication preserves
this grading. Under SGA3, Expose I,4.7.2--4.7.3, it is the
corresponding diagonal-torus action on $B_R$. The functional
$\ell$ projects to the weight-zero free summand and then
sums its basis coordinates. It is an explicitly specified
linear contraction, not the trace of the nonreduced algebra
$B_R$ and not an asserted ring homomorphism. In particular
nilpotent moment coordinates are retained rather than erased
by an ordinary algebra trace. The final universal polynomial
is the character of the usual free permutation-indexed
module with each entry expanded as $S_{ij}+\sum_aU_{ia}V_{aj}$;
the signed circuit computes that positive final character.
No canonical GCT object or effective extraction of arbitrary
representation multiplicities is being supplied.
\end{geometricstatement}
\begin{proof} The displayed dictionary transports the proof of NN-C041. Any additional realization, effectiveness or uniformity hypothesis is retained explicitly.\end{proof}


\section{NN-C042: geometric formulation and dictionary}
\begin{geometricstatement}[NN-C042]
Let
\[
 X_n=\operatorname{Spec}
       \mathbb Z[x_1,\ldots,x_n]/(x_i^2-x_i:1\le i\le n).
\]
A local Boolean map gives a morphism $X_n\to X_m$ by integer multilinear idempotents, and a Boolean target section $y$ has finite split fiber $Y_y$ of rank $|f^{-1}(y)|$. For the universal two-rail section over the coefficient base $\mathbb Z[z]$,
\[
            w=\prod_i((1-x_i)z_{i,0}+x_i z_{i,1}),
\]
the regular function
$\operatorname{Tr}_{(\mathcal O(Y_y)\otimes_{\mathbb Z}\mathbb Z[z])/\mathbb Z[z]}(w)$ is its full preimage polynomial. Under either the original exact parameter condition~\eqref{eq:C042-gate} or the stronger Catalan-threshold condition~\eqref{eq:C042-Catalan-gate}, this function has an integral constant-free circuit of size~\eqref{eq:C042-time} and formal degree at most $2n$, compatible with every commutative-ring base change. In particular the conclusion holds for $k=o(\log\log n/\log\log\log n)$, or in the explicit range~\eqref{eq:C042-Catalan-range}, with polynomially many output coordinates.

E's disjoint sparsification is a disjoint-union decomposition of this Boolean fiber on the original $n$ coordinates. Grouping clauses changes only the presentation of its acceptance idempotent; it does not change the scheme or introduce new witness coordinates. Trace adds over the disjoint components and is then computed by NN-C039's integral construction. The rank can be recovered exactly from the computed integer trace, while reduction modulo two retains only its parity. The source-interface locators are EGA IV4,18.2.1--18.2.2, printed p111, with trace base change in equation18.2.1.2. The image test concerns the support of a morphism between finite split schemes; no purity, zeta, or arbitrary representation-multiplicity extraction claim is made. This is an effective finite-scheme theorem for the specified growing local presentations, and not a polynomial-time consequence of cohomological purity.


The finite split fibers, acceptance idempotents and two-rail trace
functions of NN-C042 are unchanged. The new threshold tree gives a
finer effective presentation by disjoint Boolean subschemes on the
same original coordinates. Under~\eqref{eq:C042-Catalan-gate}, and
in particular in range~\eqref{eq:C042-Catalan-range}, their integral
trace functions have the same uniform constant-free circuit and
formal-degree bounds. Exact rank extraction and image support of
$X_n\to X_m$ therefore have the stronger growing-locality scope.
The exact trace/base-change interface remains EGA IV4,18.2.1--18.2.2,
p111, equation18.2.1.2. This alters the effective presentation bound,
not the weight-zero purity of the finite fibers, and supplies no
polynomial-time or general GCT multiplicity algorithm.

\end{geometricstatement}
\begin{proof} The displayed dictionary transports the proof of NN-C042. Any additional realization, effectiveness or uniformity hypothesis is retained explicitly.\end{proof}


\section{NN-C043: geometric formulation and dictionary}
\begin{geometricstatement}[NN-C043]
Let $b_1,\ldots,b_r$ be the independent radicands constructed in the proof and let
\[
 R=\mathbb Z[(2D\prod_j b_j)^{-1}],\qquad
 X=\operatorname{Spec}
     R[u_1,\ldots,u_r]/(u_j^2-b_j:1\le j\le r).
\]
Then $X\to\operatorname{Spec}R$ is a finite etale Kummer cover of rank $2^r$, with the independent sign action of $(\mathbb Z/2\mathbb Z)^r$. Each quadratic relation has invertible derivative $2u_j$ because $2$ and $b_j$ are units. The exact standard-etale criterion is Stacks Definition10.144.1 and Lemma10.144.2(1)--(2), tags00UB and00UC; the generic field and character group are given by Milne Theorem5.30 and Example5.31(b). The monomial basis is the corresponding sign-character decomposition over the generic fiber.

There is a deterministically constructible prime $p=\operatorname{poly}(L)$, and for every requested precision $e$ a finite etale rank-two algebra
\[
 A_e=(\mathbb Z/p^e\mathbb Z)[T]/(T^2-\eta)
\]
with an explicit $\mathbb Z/p^e\mathbb Z$-morphism
\[
       \operatorname{Spec}A_e
              \longrightarrow X\times_R\operatorname{Spec}(\mathbb Z/p^e\mathbb Z).
\]
Thus a local point of residue degree at most two suffices even though the global cover has exponential degree. No degree-one point or completely split prime is required. The Hensel construction explicitly lifts that point through every indicated nilpotent thickening.

For any encoded arithmetic section whose values on the Boolean input components descend to integers $g(x)$ with $|g(x)|\le2^w$, pulling back along this point gives $g(x)\cdot1$ in $A_{w+2}$. Its normalized finite-algebra trace is
\[
              \frac12\operatorname{Tr}_{A_{w+2}/(\mathbb Z/p^{w+2}\mathbb Z)}
                           (g(x)\cdot1)
                    =g(x)\pmod{p^{w+2}}.
\]
The inverse of two exists, and the trace is literally the trace of the two-by-two multiplication matrix. EGA IV4,18.2.1--18.2.2,p111, with base change in equation18.2.1.2, supplies the exact classical trace interface. The magnitude bound recovers the unique integer lift in polynomial time. This is an effective local-specialization and descent theorem for the given Kummer presentations, not a purity-to-multiplicity assertion or a new GCT obstruction.

Invert $2\ell(\ell-1)D\prod_a b_a$ in $\mathbb Z$, giving $R$.
The corresponding scheme $X=\operatorname{Spec}\mathcal O_\ell
\times_{\operatorname{Spec}\mathbb Z[D^{-1}]}\operatorname{Spec}R$
is finite etale of rank $(\ell-1)\ell^r$ over $R$.
Indeed the derivatives of the radical relations are units.
For the cyclotomic relation, $Z-1$ is a unit because
$\Phi_\ell(1)=\ell$ is a unit, and differentiation of
$(Z-1)\Phi_\ell(Z)=Z^\ell-1$ gives the same derivative criterion.
After the cyclotomic base change it is a $\mu_\ell^r$-torsor.
This is the standard-etale criterion in Stacks
Definition 10.144.1 and Lemma 10.144.2, tags 00UB and 00UC.

The construction gives a point over the finite etale rank-$d$
algebra $A_e$ on every indicated reduction of $X$ modulo $p^e$.
Its lifted scalar section has normalized trace
\[
       d^{-1}\operatorname{Tr}_{A_e/(\mathbb Z/p^e\mathbb Z)}
                          (g(x)\cdot1)=g(x)\pmod{p^e};
\]
$d$ is invertible by the chosen avoidance factor in $W$.
The finite-algebra trace and base-change dictionary is
EGA IV$_4$, 18.2.1--18.2.2, printed p.~111, especially
equation 18.2.1.2, as already checked in NN-C020 and NN-C043.
The global cover has exponential rank, while this explicit
local coefficient algebra has the fixed rank $\ell(\ell-1)$.

This extends a proved coefficient-presentation class relevant
to the bounded counting-output transfer. It does not construct
such presentations for arbitrary complex arithmetic circuits.
The prime $\ell$ is fixed throughout; neither a growing prime
nor mixed unrelated radical degrees nor unrestricted nested
radicals is covered by the cost and injectivity proof above.
The integer-output promise and the exclusion of computed
division remain exactly as in NN-C043. No new positivity or
occurrence obstruction follows from this local trace identity.

With the fixed field presentation and atoms above, put
$R=\mathbb Z[(2MDC R_F\prod_jc_j)^{-1}]$ and base change
$\mathcal O$ to this scalar ring. Over the fixed etale coefficient
algebra $R[Z]/(F)$,
the equations $u_j^{m_j}=c_j$ define a finite etale Kummer cover of
rank $\prod_jm_j$; its generic character group is
$\prod_j\mathbb Z/m_j\mathbb Z$. Over $R$ its total rank is
$b\prod_jm_j$. The monic presentation and Kummer degree equality
identify its selected generic field without expanding this rank.
The standard-etale criterion is Stacks Definition10.144.1 and
Lemma10.144.2, tags00UB and00UC. Kummer duality and its degree formula
are Milne Theorem5.30 and Remark5.32, pp75--76; the only additional
field source is the full primitive-element Theorem5.1/Remark5.2,
pp61--62, together with Proposition2.12, pp29--30.

For an effectively chosen prime $p$ of polynomial numerical size, the
proof constructs compatible points of this cover over the finite etale
rank-$D$ rings $A_e$, with $D$ depending only on $M$. The points preserve
all relations over $R$ among the selected integral coordinate values.
The algebraic section representing a bounded integer Boolean output
pulls back to $g(x)1$, and since $p\nmid D$ its normalized trace is
\[
           D^{-1}\operatorname{Tr}_{A_e/(\mathbb Z/p^e\mathbb Z)}
                         (g(x)1)=g(x)\pmod{p^e}.
\]
The exact finite-algebra trace and base-change interface is EGA IV4,
18.2.1--18.2.2,p111, equation18.2.1.2. The bounded-output lift recovers
the integer in polynomial time. The new conclusion is uniform local
specialization for the stated bounded flat-radical presentations;
neither a small global cover nor a purity-to-complexity implication
is asserted.

Localize the scalar base further by $2M\prod_{\ell\in S}\ell\prod_jc_j$
and the denominator $C$, obtaining a ring $R_0\subset\mathbb Q$.
The cyclotomic algebra $R_0[Z]/(\Phi_{8M})$ is finite etale, and over
it the monic ordinary and exceptional Kummer equations above define
a finite etale cover. To verify the first assertion, differentiate
$Z^{8M}-1$: its derivative is a unit since $8M$ and $Z$ are units,
and use the cyclotomic product identity to see that the derivative
of $\Phi_{8M}$ is a unit in its own quotient. For each Kummer equation
the derivative is a unit because the degree and its constant term
are units. This is the standard-etale criterion in Stacks
Definition10.144.1 and Lemma10.144.2, tags00UB and00UC.

The generic cover has the actual faithful field component established
in (C043.G4); its possibly exponential degree is not expanded. The proof
constructs compatible points over the finite etale rank-$R$ algebras
$A_e$, where $R=\varphi(K)/\varphi(p-1)$ is polynomial in $M+L$ and
$p\nmid R$. No degree-one point or completely split rational prime
is required. The coordinate algebra can be a product of local factors;
the idempotents in (C043.G8) impose the required quadratic relations
without factoring that product.

For a section with bounded integer Boolean values, pullback is the
scalar $g(x)1$. Its normalized trace is
\[
 R^{-1}\operatorname{Tr}_{A_e/(\mathbb Z/p^e\mathbb Z)}(g(x)1)
                         =g(x)\pmod{p^e}.
\]
The exact classical trace and base-change interface is EGA IV4,
18.2.1--18.2.2, printed p111, equation18.2.1.2. Integer recovery is
the proved effective step. Thus this is a uniform local-specialization
theorem for explicitly described growing flat-radical presentations.
It is not an inference of efficient multiplicity extraction from purity.
The only number-field inputs beyond the already written local construction
are the classical cyclotomic and Kummer results located above.
No Riemann-hypothesis input is used.

The growing-lcm theorem supplies a faithful finite cyclotomic--Kummer
coefficient order $\mathcal O$ whose generic degree may be exponential.
After inverting $W_0$, it is the finite etale cover specified there.
For each Boolean input the circuit describes a rational section on
the selected generic component, with finitely many required inverses.
The norm estimate bounds the number of rational primes at which one
of the nonzero integral fraction components can fail to be a unit.
It does so without constructing the corresponding divisor or norm
product. In the explicitly sampled interval, more than half the
primes avoid this locus.

For any retained prime the algorithm constructs, in expected
polynomial time, a finite etale local algebra
\[
 A_e=(\mathbb Z/p^e\mathbb Z)[T]/(G),
 \qquad \operatorname{rank}A_e=D=M\varphi(8M),
\]
and a compatible point of the coefficient cover. The cyclotomic and
Kummer root relations are imposed together. If all required
denominators pull back to units, the point extends to the specified
localization carrying the rational section, and that trial is
accepted. Every accepted pullback of a promised integer section is
exactly the same scalar integer modulo $p^e$; the acceptance rule is
therefore a local-unit test with zero output error, not a test of
whether the sampled point represents a preferred complex embedding.

Since $p>2M$, no prime factor of $M\varphi(8M)$ equals $p$, and $D$
is a unit modulo $p^e$. Thus the scalar can also be read as
\[
       D^{-1}\operatorname{Tr}_{A_e/(\mathbb Z/p^e\mathbb Z)}(g(x)1)
                            =g(x)\pmod{p^e}.
\]
The trace construction and base change are EGA IV4,18.2.1--18.2.2,
p111, equation18.2.1.2. Standard etaleness is Stacks
Definition10.144.1 and Lemma10.144.2, tags00UB and00UC. The additional
algorithmic classical inputs are precisely Rabin's certified finite-
field and root algorithms at the locators above, AKS Theorems4.1
and5.1, and Rosser--Schoenfeld Corollary3,equation(3.8).
This is a uniform expected-polynomial-time local-specialization
theorem for the stated rational sections. It supplies neither a
deterministic general root-selection procedure nor a
purity-to-multiplicity principle.

On each actual Boolean input component the register execution
specifies finitely many rational coefficient sections and tests
whether those sections vanish. The faithful cyclotomic--Kummer
order from NN-C043 and the norm estimate give an explicit sampling
interval in which fewer than $1/64$ of the rational primes can
annihilate any of the nonzero integral fraction components on that
path. For each retained prime a compatible point over a finite
etale rank-$M\varphi(8M)$ local algebra is constructible in expected
polynomial time.

At a good prime every encountered rational section is defined,
and a nonzero section pulls back to a unit. Hence the point preserves
the itinerary through the specified zero loci and principal opens.
At a bad prime that itinerary can change. The algorithm explicitly
caps both the control path and the random construction time, giving
the proved bounded-error computational statement. A correct
integer-valued final pullback is recovered by its scalar coordinate
or by the normalized finite-algebra trace from EGA IV4,
18.2.1--18.2.2,p111, equation18.2.1.2. The classical coefficient,
norm, prime-count, primality and root-finding interfaces are exactly
the previously audited ones in NN-C043 and NN-C044. The statement
concerns specified zero-test paths and does not assert order
preservation or efficient multiplicity extraction from purity.
\end{geometricstatement}
\begin{proof} The displayed dictionary transports the proof of NN-C043. Any additional realization, effectiveness or uniformity hypothesis is retained explicitly.\end{proof}


\section{NN-C044: geometric formulation and dictionary}
\begin{geometricstatement}[NN-C044]
Let $f_i,a_i,R_i,D,W$ be as in the proof, put $R=\mathbb Z[1/W]$, and define the finite cover
\[
 Y=\operatorname{Spec}
       R[T_1,\ldots,T_r]/\bigl(f_i(T_i)/a_i:1\le i\le r\bigr).
\]
It is finite etale of rank $\prod_i d_i$: the equations are monic, and inversion of the resultants makes each derivative a unit in its corresponding factor algebra. This is precisely the standard-etale criterion of Stacks Definition 10.144.1 and Lemma 10.144.2, tags00UB and00UC. The chosen actual tuple $(\alpha_i)$ selects the generic field component $K$ of $Y_{\mathbb Q}$; no algorithm for computing this component or its degree is claimed.

There is a prime $p$ of polynomial numerical size, with $p\nmid W$, and a degree-$D$ finite etale algebra
\[
       A_e=(\mathbb Z/p^e\mathbb Z)[T]/(G)
\]
for every $e\ge1$, together with compatible morphisms
\[
       \operatorname{Spec}A_e\longrightarrow
           Y\times_R\operatorname{Spec}(\mathbb Z/p^e\mathbb Z).
\]
These points preserve every relation over $R$ among the coordinates of the selected generic component: they arise from a single embedding $K\hookrightarrow A_\infty[1/p]$ whose integral constant values lie in $A_\infty$. The local rank is $D$, even when the generic component has much larger degree. The existence uses complete splitting of the residue polynomials in $\mathbb F_{p^D}$ and Hensel lifting; it does not assert a degree-one point, a completely split prime over $\mathbb Q$, or a canonical selection of these local points.

For an arithmetic section whose restriction to every Boolean input component in the selected generic field is an integer $g(x)$ with $|g(x)|\le2^w$, pullback to $A_{w+2}$ is the scalar $g(x)1$. Because $p\nmid D$, normalized finite-algebra trace gives
\[
 D^{-1}\operatorname{Tr}_{A_{w+2}/(\mathbb Z/p^{w+2}\mathbb Z)}(g(x)1)
                  =g(x)\pmod{p^{w+2}}.
\]
The trace is the trace of the multiplication matrix in the displayed free basis. Its exact classical construction and base change are EGA IV4,18.2.1--18.2.2,p111, equation18.2.1.2. The output bound recovers $g(x)$ uniquely. If $D$ and the remaining descriptions are polynomially bounded, the field presentation and compatible local point are polynomial advice and each pullback is polynomial-time computable. This is a finite-etale specialization theorem with an explicit nonuniform computational conclusion; it supplies no purity-to-multiplicity implication and no new representation-positivity criterion.

With the notation of the proof, put $R=\mathbb Z[1/W]$ and
\[
 \mathcal O=R[\gamma_1,\ldots,\gamma_r]\subset K,
 \qquad \gamma_i=a_i\alpha_i.
\]
This is a finite coefficient order over $R$. Since each $a_i$ is a unit
of $R$, it contains the original constants. Every fraction component
belongs to $\mathcal O$. If it is nonzero, its minimal-polynomial
constant term has all prime divisors among those of $W$, so the inverse
formula in the proof belongs to $\mathcal O$. Thus all the specified
rational gate sections, on every Boolean input component, are regular
sections of $\operatorname{Spec}\mathcal O$ after this one localization.
The order and localization are used as existence objects; their explicit
descriptions may have exponential size.

There exists a prime $p$ of polynomial binary length outside $W$, and
compatible finite etale rank-$D$ algebras
\[
                  A_e=(\mathbb Z/p^e\mathbb Z)[T]/(G)
\]
with morphisms
\[
 \operatorname{Spec}A_e\longrightarrow
       \operatorname{Spec}\mathcal O
                   \times_R\operatorname{Spec}(\mathbb Z/p^e\mathbb Z)
 \qquad(e\ge1).
\]
They preserve every relation over $R$ among the selected constants and
every required inverse. In particular these are points of the finite
etale ambient cover defined by the separate leaf polynomials, satisfying
the actual generic-component relations. The map is from the localized
coefficient order, not from the whole characteristic-zero field.
The standard-etale interface is Stacks Definition10.144.1 and
Lemma10.144.2, tags00UB and00UC; the existence of the coherent generic
embedding is Milne Proposition2.12.

If a rational gate section has the bounded integer value $g(x)$ on the
selected generic component, its pullback is $g(x)1$. Since $p\nmid D$,
\[
 D^{-1}\operatorname{Tr}_{A_{w+2}/(\mathbb Z/p^{w+2}\mathbb Z)}
                      (g(x)1)=g(x)\pmod{p^{w+2}}.
\]
The exact classical trace and base-change locators are EGA IV4,
18.2.1--18.2.2,p111, equation18.2.1.2. Scalar recovery and every
rational-section evaluation are polynomial-time operations given the
polynomial advice describing the local algebra and its coherent point.
This extends the finite-etale bounded-output dictionary to specified
rational sections whose denominators are everywhere defined on the
Boolean components. It makes no efficient-construction claim about the
global localization, no purity-to-multiplicity inference, and no new
representation-positivity assertion.

The finitely many actual executions produce rational functions in the
selected generic coefficient component. Invert the one integer $W$
above and let
\[
 R=\mathbb Z[1/W],\qquad
 \mathcal O=R[a_1\alpha_1,\ldots,a_r\alpha_r]\subset K.
\]
Every value occurring on an actual Boolean execution belongs to
$\mathcal O$ and is either zero or a unit, by the same minimal-polynomial
inverse formula. Thus each encountered zero locus is either the whole
selected component or empty after this localization, and every
corresponding principal open is preserved by the local point.

There are compatible morphisms from the finite etale rank-$D$ local
schemes $\operatorname{Spec}A_e$ to
$\operatorname{Spec}\mathcal O$, over the indicated reduction of $R$,
which preserve all these open-or-closed decisions simultaneously. Pullback
therefore preserves the input-specific itinerary through the finite
sequence of zero-locus and principal-open tests. The number of Boolean
inputs and the description of the localization may be exponential,
while the local algebra and its coherent point require only polynomial
advice. A bounded integer-valued final section is recovered by its scalar
coordinate, equivalently by the normalized finite-algebra trace from
EGA IV4,18.2.1--18.2.2,p111, equation18.2.1.2. The classical integrality,
norm, Hensel, and standard-etale interfaces are exactly those already
listed in NN-C044; this is their stated computational consequence for
specified zero-test paths, not a statement about ordered real sections.
\end{geometricstatement}
\begin{proof} The displayed dictionary transports the proof of NN-C044. Any additional realization, effectiveness or uniformity hypothesis is retained explicitly.\end{proof}


\section{NN-C045: geometric formulation and dictionary}
\begin{geometricstatement}[NN-C045]
Let $K$ be the actual generic coefficient field of the specified
nested-radical constant graph. Form the finite order generated by
its integral fraction components, and localize at the integer $W$
from the proof; write this order as $\mathcal O$ over
$R=\mathbb Z[1/W]$. All relevant denominator components and all
nonzero tested components are units in this order, by the monic
minimal-polynomial inverse formula. Its generic fiber is the
selected actual field $K$, rather than a product of all formally
possible radical branches. The coefficient divisions and all
sections encountered on the finite set of Boolean input paths are
regular on $\operatorname{Spec}\mathcal O$. A nonzero root equation
is a Kummer equation $U^d-a$ with $d$ and $a$ invertible; a zero
radicand is the selected zero section. No independence of the root
classes or full product degree is asserted.

There is a prime $p$ of polynomial binary length, $p\nmid W$, and
a tower of finite etale complete local $\mathbb Z_p$-algebras
\[
 A_0\subset A_1\subset\cdots\subset A_h,
 \qquad \operatorname{rank}_{\mathbb Z_p} A_j
                       =\varphi(M)M^j,
 \qquad A_j/pA_j=\mathbb F_{p^{\varphi(M)M^j}}.
\]
The coherent field embedding restricts to a point
$\operatorname{Spec}A_h\to\operatorname{Spec}\mathcal O$.
Its truncations give compatible points over $\mathbb Z/p^e\mathbb Z$
with polynomial numerical rank $D$. At each nonzero radical step
all possible transported roots lie in the next local layer, so the
chosen generic component and every relation over $R$ among its
coordinates are preserved. The construction does not require
choosing unrelated roots for different equations.

This point carries each actual execution itinerary through zero
loci and principal opens to the same itinerary in the special
fiber: a zero section pulls back to zero, and every nonzero
encountered section pulls back to a unit. If the final section is
the bounded integer scalar $g(x)$, then, because $p\nmid D$,
\[
       D^{-1}\operatorname{Tr}_{A_h/p^eA_h\,/\,\mathbb Z/p^e\mathbb Z}
                           (g(x)\cdot1)
                         =g(x)\pmod{p^e}.
\]
This is the multiplication-matrix trace and its base-change
compatibility in EGA IV4,18.2.1--18.2.2,p111,
equation18.2.1.2. The magnitude bound recovers the integer at
$e=w+2$. Thus the theorem is a bounded-rank local-specialization
statement for actual nested coefficient components and rational
execution sections. Its effectivity is polynomial advice;
it does not claim uniform extraction of a generic component,
complex order preservation, or multiplicity computation from
purity. The classical embedding, integrality, finite-field,
Hensel, norm and etale interfaces are stated with their exact
locators in the proof; no Riemann-hypothesis input is used.

For the actual generic coefficient field of this polynomial-root
constant presentation, localize its finite integral order at the
norm-avoidance integer $W$. All coefficient divisions and all
actual program sections are then regular. At every root gate the
monic squarefree replacement $Q$ has invertible resultant
$\operatorname{Res}(Q,Q')$. Its root cover is finite etale, and the
chosen actual generic component is retained through one coherent
embedding; no independent product-of-degrees assertion is made.

There exists a finite etale complete local tower with residue
ranks $1,\Lambda,\ldots,\Lambda^h$ and a compatible point of that
actual component over its top algebra. Every polynomial of degree
at most $b$ that occurs in a layer has all its reduced roots in
the next residue extension, and every simple root lifts. This is
the geometric reason the local rank depends on root depth and
$\Lambda$, despite an exponentially large possible global field
degree and nonsolvable relative Galois groups.

The point preserves every relation over the localized coefficient
ring and the zero-locus/principal-open itinerary of each actual
Boolean execution. At precision $p^e$, normalized finite-algebra
trace of its scalar output is that integer modulo $p^e$, since
$p\nmid D$. The exact trace/base-change anchor is EGA IV4,
18.2.1--18.2.2,p111,equation18.2.1.2. The conclusion is polynomial
advice for this explicit tower presentation; it does not identify
unknown generic components uniformly or infer numerical
multiplicities from purity.
\end{geometricstatement}
\begin{proof} The displayed dictionary transports the proof of NN-C045. Any additional realization, effectiveness or uniformity hypothesis is retained explicitly.\end{proof}


\section{NN-C046: geometric formulation and dictionary}
\begin{geometricstatement}[NN-C046]
The constant presentation is a succession of monic finite locally
free covers
\[
 \operatorname{Spec}\mathcal O_j\longrightarrow
                         \operatorname{Spec}\mathcal O_{j-1}.
\]
The full-degree promise makes each generic fibre integral and makes
the monic algebra the actual coefficient order. On the principal
open where the degrees and nonzero radicands are units, each root
cover is standard etale: its derivative is $d_jY_j^{d_j-1}$.
This is Stacks Definition10.144.1 and Lemma10.144.2, tags00UB and00UC.
Fraction gates are rational sections on their stated principal
opens. The constant-depth filtration prescribes the residue-field
tower degrees $D_j$; the cyclic finite-field calculation (C046.8)
splits all fibres at the next level, with no choice of a preferred
geometric point required.

The norm estimate gives a uniformly samplable positive density of
primes at which all sections encountered on the actual Boolean
component remain defined and all nonzero ones pull back to units.
The monic relations make every constructed local point a point of
the actual order, so the integer section (C046.11) is preserved at
every accepted point. This is the geometric content of zero-error
evaluation. Exact-zero branches specify an itinerary through zero
loci and principal opens; good primes preserve that itinerary,
whereas bad primes may change it, accounting for the bounded-error
branch theorem.

The rank-$D$ residue algebra is finite etale. Its scalar output may
also be recovered by normalized trace: all prime divisors of
$D=\varphi(M)M^h$ are at most $M$, so $p>2M$ makes $D$ a unit.
The trace interface is EGA IV4,18.2.1--18.2.2,printedp111,
equation18.2.1.2. This statement concerns effectively presented
coefficient covers, not extraction of GCT multiplicities from
weights or purity.

Geometrically, these are full-degree monic coefficient covers.
Their generic fibres are integral. Localizing away from the monic
discriminants makes their root covers finite etale; the residual
rank filtration is $1,\Lambda,\ldots,\Lambda^h$. All their reduced
fibres split at the next prescribed level and every simple root
lifts, so every constructed point is a point of the actual order.
Integer sections and good-prime zero-locus itineraries are preserved
as above. Since every prime factor of $D$ is at most $b$, normalized
trace is available for $p>2b$ with the same EGA IV4 locator. This
strictly enlarges the uniform coefficient presentation to certain
nonsolvable extensions, while retaining the explicit full-degree
promise and leaving arbitrary complex-constant elimination open.
\end{geometricstatement}
\begin{proof} The displayed dictionary transports the proof of NN-C046. Any additional realization, effectiveness or uniformity hypothesis is retained explicitly.\end{proof}


\section{NN-C047: geometric formulation and dictionary}
\begin{geometricstatement}[NN-C047]
For any nonzero odd integer $b$ represented by the preceding
straight-line program, let
\[
 R=\mathbb Z[(2b)^{-1}],\qquad
 X_b=\operatorname{Spec}R[z]/(z^2-b^2).
\]
The base and coefficients are succinctly given by the same
arithmetic graph. If an ordinary polynomial presentation is
wanted, introduce one coordinate for each graph value and its
triangular addition or multiplication equation, and a coordinate
$u$ with $(2b)u=1$. This gives an equivalent base presentation
with polynomially many equations and bounded degrees, without
writing the binary expansion of $b$.

The two factors $z-b$ and $z+b$ are comaximal because $2b$ is a
unit. Evaluation at $b$ and $-b$ explicitly gives
\[
                  R[z]/(z^2-b^2)\simeq R\times R.
\]
An inverse sends $(r,s)$ to
$r(1+z/b)/2+s(1-z/b)/2$. Thus this is a split finite etale cover
of rank two; the derivative $2z$ is a unit. The exact standard-etale
anchor is Stacks Definition10.144.1 and Lemma10.144.2,
tags00UB and00UC. The section
\[
               e_b=\frac{1+z/b}{2}
                   =\frac{b+z}{2z}
\]
is an idempotent, with values $(1,0)$ on the two components.
Its multiplication matrix has trace one and rank one. The section
$1-e_b$ has the same trace and rank. This trace is the explicit
finite-algebra trace of EGA IV4,18.2.1--18.2.2,p111,
with base change in equation18.2.1.2.

At every prime not dividing $2b$ the fiber consists of two rational
points, with identity Frobenius on its degree-zero etale
cohomology and no positive-degree geometric cohomology. Its
count is two and its local zeta function is $(1-T)^{-2}$.
The generic two-point permutation representation is likewise
trivial. The arithmetic covers for $b$ and $-b$ are literally the
same, and their idempotents are exchanged; all the listed cover,
trace and rank data agree.

In the real fiber the principal-root point is the one with
$z=|b|>0$. Pulling back $e_b$ along that selected point gives one
when $b>0$ and zero when $b<0$, exactly (C047.1). Hence a uniform
algorithm that recovers this selected component's scalar from the
succinct presentation solves $\mathrm{PosSLP}$, despite the
split, zero-dimensional, pure Tate cover and rank-one idempotent.
The positivity selection is additional real-fiber data; finite
etale trace or purity alone does not identify it. This statement
does not claim a complexity lower bound for explicit matrices,
whose displayed two-by-two trace is immediate.

\paragraph{The actual split cubic and selected principal component.}
Use the specified complex embedding of $K$ and put
\[
 R=\mathbb Z[\eta,(3A)^{-1}],\qquad
 \zeta_3=\eta^2,\qquad
 \mathcal B=R[T]/(T^3-A^3).
 \tag{C047c.7}
\]
The coefficients and localization may be given by the original
straight-line graph, with one coordinate and triangular equation
per gate and a final inverse equation. This is a polynomial-size
bounded-degree presentation; no binary expansion of $A$ is used.
Since $A$ is a unit, setting $U=T/A$ identifies
$\mathcal B$ with $R[U]/(U^3-1)$.
The three roots $1,\zeta_3,\zeta_3^2$ are distinct and their
differences are units: their relevant product is a unit multiple
of $3$. The Chinese remainder map therefore gives
\[
                  \mathcal B\simeq R\times R\times R,
 \qquad
                  e_0=\frac{U^2+U+1}{3}\longmapsto(1,0,0).
 \tag{C047c.8}
\]
Thus this is an actual split finite etale cover of rank three,
and $e_0$ is a genuine idempotent. The derivative $3U^2$ is
a unit. The exact standard-etale locators are Stacks
Definition~10.144.1 and Lemma~10.144.2, tags 00UB and 00UC,
as in NN-C047.

Multiplication by each primitive idempotent has rank and
trace one. The trace/base-change dictionary is EGA IV$_4$
18.2.1--18.2.2, printed p.~111, particularly (18.2.1.2).
Every finite-field fiber of this split cover has three
rational points and local zeta factor $(1-T)^{-3}$;
its permutation representation over the base field is trivial.
These cover, trace and rank data do not choose a component.

At the archimedean principal-root point,
$U=\gamma/A$ is $1$ if $a>0$, and
$-\eta=\zeta_3^2$ if $a\le0$.
Hence pulling back the actual idempotent $e_0$ gives precisely
the bit in (C047c.5). The linear expression
$(U+\eta)/(1+\eta)$ gives that same bit on these two selected
components; it is not claimed to be an idempotent on the full
three-component cover.
Selecting the principal complex branch from the succinct graph
therefore carries the PosSLP decision, despite the fixed
degree-two containing field and split etale cubic presentation.
The actual selected factor has relative degree one, whereas
the displayed equation has degree three. An arbitrary point
of the displayed cover need not preserve that selected-factor
relation. No Frobenius, trace-positivity or canonical-multiplicity
argument has been asserted to recover the archimedean branch.

\paragraph{The selected real point in a tower of finite covers.}
Fix one defined executed path and its rational input point.  Successive
root instructions define finite monic algebra extensions
$R[T]/(T^d-a)$, after restricting to the principal opens on which
the denominators used by the path are invertible.  At each real root
step the semantic point is selected by $T\geq0$ when $d$ is even,
and by the unique real solution when $d$ is odd.  Thus the algorithm
selects a point in a tower of finite real fibers, with sign tests
choosing semialgebraic subsets and domain tests choosing the stated
open conditions.  The tower need not be etale at zero radicands,
and the selected-point evaluation is not the finite-algebra trace or
an average over conjugates.  The shared DAG is a succinct description
of this tower and of regular functions after the required
localizations; it is not an explicit basis of its potentially
exponential rank.  The exact classical localization dictionary is
Stacks Project, Lemmas 10.17.4--10.17.7, tags 00E2--00E5.  The
complexity assertion concerns the selected real-point functor with
its order data; no unprovided reductive representation or
cohomological multiplicity is inferred.
\end{geometricstatement}
\begin{proof} The displayed dictionary transports the proof of NN-C047. Any additional realization, effectiveness or uniformity hypothesis is retained explicitly.\end{proof}


\section{NN-C048: geometric formulation and dictionary}
\begin{geometricstatement}[NN-C048]
Let $R=\mathbb Z[1/m]$. The primitive-character cover of the
constant cyclic group of order $m$ is the finite etale scheme
\[
 \begin{gathered}
 \mathcal P_m=\operatorname{Spec}R[U]/(\Phi_m(U)),
 \qquad \operatorname{rank}_R\mathcal P_m=\varphi(m),\\
 \operatorname{Gal}(\mathbb Q(\zeta_m)/\mathbb Q)
                   \simeq(\mathbb Z/m\mathbb Z)^\times.
 \end{gathered}
\]
It has the alternative succinct presentation
\[
 R[U,V_1,\ldots,V_r]\big/
 \big(U^m-1,\ V_i(U^{m/\ell_i}-1)-1\ (1\le i\le r)\big).
                                                               \tag{C048.11}
\]
Binary powering presents these equations by circuits of total
size polynomial in $\log m$ and the supplied factorization size.
To verify the presentation, decompose $U^m-1$ by its cyclotomic
factors over $R$. They are pairwise comaximal: in every residue
characteristic prime to $m$ their roots are disjoint because
$U^m-1$ is separable, and the same holds over $\mathbb Q$; the
finite free Chinese-remainder map is therefore an isomorphism
on every fiber and hence an isomorphism. On a proper divisor
component $d<m$, some $\ell_i$ has $d\mid m/\ell_i$, so inverting
$U^{m/\ell_i}-1$ removes that component. On the primitive
component each of these functions is a unit on every fiber,
so its inversion changes nothing. This proves (C048.11).

For $p\equiv1\pmod m$ the fiber consists of $\varphi(m)$
rational points, one for each primitive character over
$\mathbb F_p$. The certified primitive-root and Hensel algorithms
construct one compatible point over $\mathbb Z/p^e\mathbb Z$
without listing the other points or expanding the rank-$\varphi(m)$
algebra. A says exactly enough about abelian Dirichlet factors
to supply the explicit progression density (C048.6)--(C048.8).
The resulting local coefficient algebra has rank one, even
when the generic cover has exponential rank in its succinct
description.

On each actual Boolean input path, the circuit gives rational
sections of this faithful cover. Its nonzero integral numerator
and denominator sections have nonzero integer norms with the
budget (C048.1). A prime avoiding their norms makes their
pullbacks units under every chosen primitive point. Every
completed straight-circuit evaluation preserves its rational
integer section exactly; good primes additionally preserve the
itinerary through the zero loci and principal opens of an
exact-zero program. The scalar is the trace of multiplication
in the rank-one local coefficient algebra. This is the literal
finite locally free algebra trace of EGA IV4,
18.2.1--18.2.2,p111, with base change in equation18.2.1.2.
The finite-etale simple-root description is also Stacks
Definition10.144.1,tag00UB, and Lemma10.144.2,tag00UC.
This computation does not build a global trace matrix of
exponential rank or extract multiplicities from purity.

\emph{Factorization-free geometric statement.} Over $R=\mathbb Z[1/m]$, replace the primitive-character cover
in (C048.11) by the explicitly presented finite etale open and
closed subscheme of $\mu_m$ obtained by inverting only
$U^{m/\ell}-1$ for the prime divisors $\ell\le256\ell_2(m)$.
The full-order component remains among its components. For
every splitting prime, a uniform point of this partial cover
lies in the primitive component with probability greater than
$255/256$, by the exact cyclic-group calculation (C048.15).
All points lift simply over $\mathbb Z/p^e\mathbb Z$; those
on the primitive component preserve the actual coefficient
order. A norm-good prime additionally preserves every principal
open and zero locus visited by the actual program.

Thus the effective finite-cover statement and (C048.16) give
the claimed bounded-error pullback algorithm without computing
the complete conductor factorization. A proper-order component
need not preserve a rational identity of the primitive generic
component, so no zero-error geometric specialization claim is
made for it. The grading-character, finite-etale and rank-one
trace locators are exactly those of NN-C048; the partial-factor
sampling mechanism is credited to the primary CIT theorem above.
\end{geometricstatement}
\begin{proof} The displayed dictionary transports the proof of NN-C048. Any additional realization, effectiveness or uniformity hypothesis is retained explicitly.\end{proof}


\section{NN-C049: geometric formulation and dictionary}
\begin{geometricstatement}[NN-C049]
Over $R=\mathbb Z[1/(2\prod_i a_i)]$, let
\[
 \mathcal X=\operatorname{Spec}
       R[Y_1,\ldots,Y_r]/(Y_i^2-a_i:1\le i\le r).
                                                               \tag{C049.9}
\]
This is the actual finite etale coefficient cover of degree
$2^r$, with geometrically split sign action $\mu_2^r$.
Its monomial basis consists of the products
$\prod_iY_i^{e_i}$ and transforms under the corresponding
one-dimensional sign characters. Its generic fiber is the
field in (C049.2), not a union of unrelated root components.
Each individual monic quadratic chart is etale because
$2Y_i$ is a unit, the exact interface of Stacks
Definition10.144.1,tag00UB, and Lemma10.144.2,tag00UC.

For every odd prime outside the stated denominators, base
change to $\mathbb F_{p^2}$ splits this cover into $2^r$
rational points. One such point is constructed in expected
polynomial time using (C049.7)--(C049.8), without writing
the cover's exponential multiplication table. Frobenius over
$\mathbb F_p$ acts on each square-root coordinate by a sign,
so its square fixes every point; residue fields of degree
one and two both embed in the target. No Frobenius density
estimate is needed for this statement.

A circuit defines an integral section $\alpha$ of this
coefficient cover. Its norm, the determinant of multiplication
by $\alpha$ on the finite free algebra, is an integer with
the bound (C049.3)--(C049.4). If it is nonzero modulo $p$,
that multiplication endomorphism is invertible on the whole
fiber. Every point's pullback of $\alpha$ is then nonzero.
A zero integral section pulls back to zero at every point.
This is the geometric content of the one-sided algorithm.
The determinant norm and its product formula are Milne's
definition and theorem cited above. Determinant commutes with
base change by its polynomial expression in matrix entries.
The separate finite locally free trace and its base-change
formalism are EGA IV4,18.2.1--18.2.2,p111,
equation18.2.1.2. This theorem concerns
succinct sections and explicit small-rank residue points;
it makes no sign, purity, or GCT multiplicity inference.
\end{geometricstatement}
\begin{proof} The displayed dictionary transports the proof of NN-C049. Any additional realization, effectiveness or uniformity hypothesis is retained explicitly.\end{proof}


\section{NN-C050: geometric formulation and dictionary}
\begin{geometricstatement}[NN-C050]
The tower (C050.4) consists of monic finite locally free covers
of degrees $d_j$, with integral generic fibers by (C050.1).
Its final affine scheme is the actual coefficient order
$\mathcal X=\operatorname{Spec}\mathcal O$, of rank
$N=\prod_jd_j$ over $\mathbb Z$. On the principal open where
degrees, radicands and the relevant denominators are units,
each stage is a finite etale Kummer cover: the derivative is
$d_jY_j^{d_j-1}$. This is the exact simple-root interface of
Stacks Definition10.144.1 and Lemma10.144.2, tags00UB and00UC.
At an individual stage the unit-radicand cover is a
$\mu_{d_j}$-torsor with monomial characters
$1,Y_j,\ldots,Y_j^{d_j-1}$. No global product action on the
whole nested tower is asserted: later radicands can depend on
earlier coordinates.

For every good prime in the progression $2\bmod m$, the power
map $z\mapsto z^{d_j}$ is a permutation of $\mathbb F_p^\times$
at every stage. Consequently the tower has one compatible
degree-one residue point over that prime, constructible by
binary exponentiation, and this point lifts uniquely through
the finite rings $\mathbb Z/p^e\mathbb Z$. The full degree-$N$
geometric fiber need not split over $\mathbb F_p$; other closed
points or geometric roots are not constructed. The statement
is an effective rank-one specialization of a possibly
exponentially large finite coefficient cover.

Norms are determinants of multiplication on its finite free
coordinate algebra. Invertibility of these determinants at
good primes makes the encountered nonzero sections invertible
on the whole fiber, and hence at the selected point. A rational
integer section represented by $A/B$ satisfies the global
order identity $A-gB=0$, so every completed local pullback
recovers $g$ at finite precision. Zero sections pull back to
zero even at bad primes, explaining the coRP identity test.
Exact-zero controls are itineraries through zero loci and
principal opens; norm-good primes preserve the actual itinerary,
giving bounded error for branching. The norm definition and
product formula are Milne FT,p82--84; determinant specialization
follows from its polynomial expression in matrix entries.
The rank-one target has ordinary trace equal to its scalar value.
The separate trace and its base-change formalism are EGA IV4,
18.2.1--18.2.2,printedp111,equation18.2.1.2. An exponential
global trace computation is not required.
\end{geometricstatement}
\begin{proof} The displayed dictionary transports the proof of NN-C050. Any additional realization, effectiveness or uniformity hypothesis is retained explicitly.\end{proof}


\section{NN-C051: geometric formulation and dictionary}
\begin{geometricstatement}[NN-C051]
Over $R=\mathbb Z[1/(D\prod_jb_j)]$, the coefficient scheme
\[
       \mathcal X=\operatorname{Spec}
          R[Y_1,\ldots,Y_r]/(Y_j^D-b_j)
\]
is finite locally free of degree $D^r$ and etale, because
$DY_j^{D-1}$ is a unit. It is a torsor under $\mu_D^r$;
the separate monomial basis gives its character grading after
splitting that group scheme. Its generic algebra is the
actual positive-radical number field, by the classical
degree lemma and (C051.3). The finite locally free trace and
its base-change dictionary are EGA IV4, 18.2.1--18.2.2, p.~111,
equation 18.2.1.2; the monic etale chart is Stacks
Definition 10.144.1, tag 00UB, and Lemma 10.144.2, tag 00UC,
with the SGA 1, Expose I, Section 3 relative-differentials
criterion. These are the same exact interfaces already
recorded in NN-C046--C049.

At every admissible progression prime, the proof constructs
an $\mathbb F_{p^2}$-point of this scheme, or already an
$\mathbb F_p$-point in the odd case, in expected polynomial
bit time from its succinct presentation. This does not say
that the entire cover splits over that field. The norm of a
section is the determinant of its multiplication map, by
Milne FT's norm definition, p.~82; determinant specialization
follows directly from its polynomial expression. If
this norm is invertible on the base, every point pulls the
section back to a unit. Hensel lifts these points to the
explicit length-$h$ unramified local algebras. Thus neither
the exponential global rank nor the number of geometric
points has to be written out by the evaluator. No purity,
Frobenius-sign distribution or canonical GCT multiplicity
is inferred from these point constructions.
\end{geometricstatement}
\begin{proof} The displayed dictionary transports the proof of NN-C051. Any additional realization, effectiveness or uniformity hypothesis is retained explicitly.\end{proof}


\section{NN-C052: geometric formulation and dictionary}
\begin{geometricstatement}[NN-C052]
For every supplied input and actual path, the real odd-root
description defines an actual order
$\mathcal O_x\subset K_x\subset\mathbb R$, finite free over
$\mathbb Z$ with rank at most $m^U$. The triangular monic
presentation maps onto this order but need not be faithful.
For every $p>m$ with $p\equiv2\pmod m$, the proof constructs
only existentially a compatible actual-field embedding into
a finite totally ramified $E/\mathbb Q_p$ with residue
$\mathbb F_p$. Whenever the explicit modular division tests
complete, the actual geometric specialization is
\[
 \mathcal O_x[\text{executed denominators}^{-1}]
       \longrightarrow\mathcal O_E
       \longrightarrow\mathbb F_p.                        \tag{C052r.15}
\]
Thus the computed point belongs to the actual selected
component, including when a displayed root equation has
repeated roots on the special fiber. No rational-point
lifting to $\mathbb Z_p$ is asserted at such a prime.
The finite cover $\operatorname{Spec}\mathcal O_E\to
\operatorname{Spec}\mathbb Z_p$ can be ramified; neither
its rank, its basis nor its nonreduced special fiber is
constructed. A zero root step is not called an etale chart.
Milne ANT 7.31, 7.38, 7.42 and 7.55 supply the exact
complete-DVR and ramification dictionary used here.

For an integral section $z\in\mathcal O_x$, its norm is
$\det(m_z)$, the determinant of its multiplication map;
this is Milne FT's norm definition, p.~82, with the
conjugate-product formula on pp.~83--84. Matrix determinant
commutes with specialization because it is an integer
polynomial in the matrix entries. At a norm-good prime
the determinant is invertible, so every encountered
nonzero section pulls back to a unit. The separate algebra
trace and its base-change formula are EGA IV4,
18.2.1--18.2.2, p.~111, equation 18.2.1.2; that locator is
used for trace, not as a definition of the determinant norm.
On the rank-one target $\mathbb F_p$, scalar multiplication
has trace equal to its scalar. If $A/B=g\in\mathbb Z$,
the actual order identity $A-gB=0$ survives (C052r.15),
and $p>2^{w+1}$ uniquely recovers the bounded integer.

For a promised family on all Boolean inputs, the orders
$\mathcal O_x$ form a finite locally free algebra over the
split Boolean base $\prod_x\mathbb Z$, with ranks allowed
to differ on its clopen components. The algorithm fixes
one input and does not construct this exponentially large
family. Exact-zero programs select an itinerary of zero
loci and principal opens, preserved at norm-good primes.
The actual selected component need not carry the entire
$\mu_d$ action of a displayed root cover, and successive
coefficient-dependent covers need not have a product torus
action. No canonical GCT multiplicity, purity assertion,
or fast global norm/trace computation follows from these
residue-point constructions.

\paragraph{Mixed flat initialization, under A.}
The initial flat presentation is the actual finite free
real order \(\mathcal O_0\) in (C052f.2), of rank
\(D^r\), with its separate monomial grading after
adjoining the relevant roots of unity. Away from
\(D\prod b_j\), its binomial charts are etale.
For every prime \(p>Q\) in the explicit reduced
class \(-1\bmod Q\), (C052f.6) constructs an
\(\mathbb F_p\)-point by binary powering. Its
simple lifts give the actual embedding
\(K_0\hookrightarrow\mathbb Q_p\). This is one
point of the initial cover, not a claim that all
its conjugates split over \(\mathbb F_p\).

The subsequent real odd-root description defines
an actual finite free order \(\mathcal O_x\) in
the real field \(K_x\). Its displayed monic cover
need not be faithful or etale at zero-root fibers.
Every completed trial is the genuine point
(C052f.8), through an unconstructed totally ramified
complete DVR. This point preserves every relation
in the actual localized order, including those
between the initial even radicals and the later
odd-root tower. The preceding real order need not
be stable under every torus action suggested by
the syntactic binomial covers.

The multiplication determinant norm and conjugate
formula are Milne FT, p.~82 and pp.~83--84;
determinant specialization is its elementary
polynomial-in-entries identity. The separate
algebra trace and its base change are EGA IV4,
18.2.1--18.2.2, p.~111, equation 18.2.1.2.
At a good prime every encountered nonzero section
pulls back to a unit. On the rank-one residue
target the trace of a scalar is that scalar,
and the actual relation \(A-gB=0\) plus the large
prime bound recovers the bounded integer.

This theorem enlarges the allowed coefficient
presentation but not the gate convention for
computed even radicals. A succinct arithmetic
radicand is not an explicitly written integer:
replacing it by its potentially enormous value
would invalidate the polynomial description bound.
The principal-root/sign obstruction NN-C047
therefore remains outside this mixed theorem.
No deterministic identity algorithm, unrestricted
complex-root selection, sign test, purity or
GCT multiplicity assertion is inferred.

\paragraph{The three actual constructible strata.}
Fix a Boolean input. Let \(\mathcal O_x\) be the actual integral
order generated by all values in the compiled division-free circuit;
these values are algebraic integers when rational leaves are represented
by the integer pairs in (C052d.2). In a mixed model the initial faithful
order and its fixed explicit denominator localization can instead be
retained. The actual relations, rather than an independently chosen
component of a formal binomial tower, define this order. Let \(b\) be
the product of the explicitly written nonzero rational-leaf denominators,
and put \(X_x=\operatorname{Spec}\mathcal O_x[b^{-1}]\). This
localization does not change the real characteristic-zero classification.
The two compiled values \(V,Z\) are global sections. The original
rational computation lives on the principal open \(D(V)\): every
\(B_j\) is a unit there because it is a product of the rational-leaf
denominators and powers or products of the denominator numerators
that divide \(V\). The original output is consequently the regular
section \(A_{\mathrm{out}}/B_{\mathrm{out}}\) on \(D(V)\).

The three constructible strata are
\[
        V(V),\qquad D(V)\cap V(A_{\mathrm{out}}),\qquad D(Z).
                                                                    \tag{C052d.5}
\]
The notation \(V(f)\) in (C052d.5) denotes a zero locus, whereas the
unparenthesized \(V\) is the compiled section. Their pullbacks to the
specified real point express exactly Undefined, Defined-zero and
Defined-nonzero. Their set-theoretic union identities are
\(V(Z)=V(V)\cup V(A_{\mathrm{out}})\) and
\(D(Z)=D(V)\cap D(A_{\mathrm{out}})\), since \(Z=VA_{\mathrm{out}}\).
These formulas include the case \(V=0\), when \(D(V)\) is empty.

The exact affine localization and closed-subscheme interfaces are
Stacks Project Lemmas 10.17.4--10.17.7, tags 00E2--00E5:
\url{https://stacks.math.columbia.edu/tag/00E2},
\url{https://stacks.math.columbia.edu/tag/00E3},
\url{https://stacks.math.columbia.edu/tag/00E4}, and
\url{https://stacks.math.columbia.edu/tag/00E5}.
Their complete short proofs specify pullback of principal opens and
the prime ideals of a localization or quotient. For the arithmetic
specializations and trace/base-change dictionary one retains exactly
the EGA IV, fourth part, 18.2.1--18.2.2, printed p.~111, interface
of NN-C052; the determinant norm remains the separately cited Milne
field-norm interface. No trace is being used to decide which of the
three strata contains the real point. The decision algorithm is the
proved finite-residue zero test applied to its two explicit sections.
\end{geometricstatement}
\begin{proof} The displayed dictionary transports the proof of NN-C052. Any additional realization, effectiveness or uniformity hypothesis is retained explicitly.\end{proof}


\section{NN-C053: geometric formulation and dictionary}
\begin{geometricstatement}[NN-C053]
The actual finite free orders of NN-C050--C052 and their
localized sections are unchanged. This theorem gives an
unconditional supply of residue characteristics admitting
their explicitly computed points of degree one or two.
For the real odd-root model, the compatible actual-field
embedding factors through the valuation ring of a finite
totally ramified extension of $\mathbb Q_p$; its reduced
residue field is still $\mathbb F_p$. The finite ramified
cover itself is a proof device and is not constructed.
For flat radicals the faithful monic cover admits the
rank-at-most-two points supplied by its degree conditions.
A zero section pulls back to zero at every completed point;
a section whose determinant norm is invertible pulls back
to a unit. The norm definition is Milne FT, p.~82, with
its product and integrality formulas on pp.~83--84.
Its specialization is the elementary polynomial formula
for a determinant. The separate trace and base-change
formalism is EGA IV$_4$,18.2.1--18.2.2,p.~111,
(18.2.1.1)--(18.2.1.2). In a degree-$f$ residue field,
scalar multiplication by $g$ has trace $fg$; here $f=1$
or $2$ and $p$ is odd. The readout above uses the scalar
coordinate directly, without computing the enormous
trace or norm of the source order.
The classical complete-DVR dictionary is the exact
Milne ANT7.31,7.38,7.42,7.55 interface recorded in NN-C052.
This changes the arithmetic supply of specializations,
not their ranks, actual components or geometric interpretation.
No finite Weil object or canonical GCT multiplicity is supplied.

\paragraph{Fixed-alphabet mixed actual-order specializations.}
For every fixed \(\Pi\), the faithful initial real
orders and the actual later orders are exactly those
of (C052f.2) and (C052f.8). The unconditional sampler
constructs degree-one residue points at primes for
which every initial flat binomial chart has a simple
point and the later actual real odd-root field embeds
through a finite totally ramified local extension.
Its residue remains \(\mathbb F_p\), including at
zero-root special fibers. No basis of the potentially
exponential order or auxiliary complete DVR is built.
The point preserves every actual localized relation;
norm-good sections pull back to units, and bounded
integer sections are recovered from the single scalar
residue. The norm and trace/base-change locators remain
the distinct Milne FT and EGA IV4 interfaces recorded
in NN-C052, respectively.

The new finite-family computation concerns classical
real-character \(L\)-values at one. It certifies a
zero-free real neighborhood for the fixed conductors
dividing \(24R\), and thereby supplies arithmetic
specializations without assuming A. It does not
construct a Weil object, a canonical GCT multiplicity,
or a uniform zero-free neighborhood over growing
prime alphabets. The constants may depend arbitrarily
but effectively on the fixed alphabet. The general
mixed theorem with a growing alphabet retains A.

\paragraph{Actual mixed orders with quadratic residue points.}
The actual initial real order is (C053q.9), finite
free of rank \(D^r\), with simple binomial charts
away from \(D\prod b_j\). For the unconditionally
sampled primes, its explicitly constructed
\(\mathbb F_{p^2}\)-point lifts to
\(K_0\hookrightarrow E_2=\mathbb Q_p(i)\).
Each later actual real component embeds through
a finite totally ramified extension of \(E_2\).
Consequently its actual integral order, localized
at completed denominators, has a genuine point
\[
             \mathcal O_x[\mathcal D^{-1}]
                           \longrightarrow\mathbb F_{p^2}
                                                               \tag{C053q.17}
\]
that preserves every actual relation, including
between the initial even radicals and later odd
ones. The displayed monic tower need not be a
faithful order at the later gates and need not
be etale at zero-root fibers; the argument uses
the actual minimal components. The two-dimensional
residue field is constructed, but the potentially
exponential-rank order and complete ramified
extensions are not.

Invertibility of the determinant norm of a section
ensures its image is a unit. The norm/conjugate
formulas are Milne FT v5.10, p.~82 and pp.~83--84;
determinant specialization is the polynomial
matrix formula. The separate trace/base-change
locator is EGA IV$_4$, 18.2.1--18.2.2, printed
p.~111, especially (18.2.1.2). A scalar integer
section has residue-field trace twice its scalar
coordinate, so the readout is exact without
forming a trace on the original enormous order.
Frobenius is the conjugation \(a+bi\mapsto a-bi\);
an integer output is Frobenius-fixed. No purity,
positivity or canonical GCT multiplicity claim
is made.

The exclusion of a later degree divisible by three
is a precise limitation of this residue construction:
for every prime \(p\ne3\), one has
\(3\mid p^2-1\), and the cube map on
\(\mathbb F_{p^2}^\times\) is not bijective.
The uniqueness step (C053q.13)--(C053q.14) would
therefore fail. This is not a lower bound for
general mixed-root algorithms. Initial flat
degrees may still contain three, because their
odd-root inversion occurs in \(\mathbb F_p\),
before the quadratic square-root stages. The
general mixed model retains the separately
stated A-dependent prime-field route; its fixed
prime-alphabet subclass is already unconditional.
\end{geometricstatement}
\begin{proof} The displayed dictionary transports the proof of NN-C053. Any additional realization, effectiveness or uniformity hypothesis is retained explicitly.\end{proof}


\section{NN-C054: geometric formulation and dictionary}
\begin{geometricstatement}[NN-C054]
For a defined complex principal-root circuit, the actual scaled roots
$\gamma=B\alpha$ and their algebraic-integer arithmetic components
generate a finite order $\mathcal O\subset K$, where
$[K:\mathbb Q]\le N$.  Inverting the finitely many actual nonzero
denominator components gives an order localization on which every
executed value is a regular function.  A displayed root step is the
finite free cover
\[
 \operatorname{Spec}R[T]/(T^d-a)\longrightarrow\operatorname{Spec}R
\]
after the required localization.  Its displayed rank is $d$, even
when the chosen generic component has smaller degree.  It is etale
where $d$ and $a$ are invertible, by the monic derivative criterion
(Stacks Project, Definition 10.144.1 and Lemma 10.144.2,
tags 00UB and 00UC); ramification at zero is permitted.

The compositum order generated by $\mathcal O$, its conjugate and $i$
has generic degree at most $2N^2$.  The displayed integral formulas
for real and imaginary parts give a norm separation of its selected
archimedean coordinates by $2^{-B_0}$.  On each actual fiber the
algorithm constructs a polynomial-size rational description of the
point selected successively by the principal argument convention,
to any absolute accuracy $2^{-P}$ whose exponent has polynomial
binary length, together with exact coordinate signs.  It follows
that the selected-point regular-function zero and coordinate-sign
problems for these succinct cover towers are Turing-equivalent to
$\mathrm{PosSLP}$ in the stated uniform model.
Their Boolean versions are also many-one equivalent by the separately
credited local adaptive-query compilation theorem.

The principal selection at the negative real cut is archimedean
order and argument data; it is not a globally algebraic section of
the power cover.  The proof does not construct a polynomial-size
basis of $\mathcal O$, replace evaluation at a selected point by a
trace over all conjugates, or infer a Frobenius/representation
multiplicity algorithm.  The norm product used for separation is
the classical number-field norm, not an invocation of a trace-form
locator.  Principal-open localization has the exact dictionary
Stacks Lemmas 10.17.4--10.17.7, tags 00E2--00E5.

\paragraph{Explicit cyclic branch action.}
For a nonzero actual radicand over a characteristic-zero field
containing $\zeta_d$, the fiber of $T^d-z$ is a torsor for the
cyclic group generated by multiplication by $\zeta_d$.
The displayed label specifies which translate of the distinguished
principal point is evaluated.  At $z=0$ all labels have the same
point, while the displayed degree-$d$ fiber can be nonreduced.
The auxiliary cyclotomic point is itself specified by the principal
root of $-1$, so its computation falls inside NN-C054 rather than
assuming an external compatible embedding.  This explicit finite
cyclic action is the representation dictionary: on the coordinate
line $T$, the label $k$ acts by the scalar $\zeta_d^k$, the
tautological one-dimensional character.  It supplies no efficient algorithm for an
unspecified generic component or an unlabeled coherent tuple of
arbitrary algebraic constants.  The root-cover and etale-away-from-zero
source interfaces are exactly those already recorded for NN-C054.
\end{geometricstatement}
\begin{proof} The displayed dictionary transports the proof of NN-C054. Any additional realization, effectiveness or uniformity hypothesis is retained explicitly.\end{proof}


\section{NN-C055: geometric formulation and dictionary}
\begin{geometricstatement}[NN-C055]
Let $\mathcal P$ be the affine scheme over $\mathbb Z$ whose coordinates
are the nested generators of (C055.1), and let
$\alpha:\mathcal P\to\operatorname{Mat}_{n,\mathbb Z}$ be the
resulting polynomial matrix morphism.  Then the regular function
$\alpha^*\operatorname{per}$ has the integral circuit (C055.3),
of formal degree $O(n^2(k+1)^2)$ and actual degree at most
$n(2h+3)$.  For separately supplied coordinates the same identity
holds on the closed polynomial compatibility locus, with no
claim of an identity on the larger space of unrelated coordinates.

For a leaf $J$, the finite etale Boolean scheme
\[
 Y_J=\operatorname{Spec}
   \mathbb Z[b_1,\ldots,b_m]/(b_1^2-b_1,\ldots,b_m^2-b_m)
       \simeq\coprod_{2^m}\operatorname{Spec}\mathbb Z
\]
has relative dimension zero and finite free rank $2^m$.
After base change to the explicit truncated algebra (C055.7),
let $f_J(t,b)$ be the product of the $3m$ factors in (C055.6),
with $u_i=R_J(i,:)x_J$ and $v_j=C_J(j,:)y_J$ already substituted.
The leaf boundary section is exactly
\[
                   P_J=[t^m]\operatorname{Tr}(m_{f_J}),
\tag{C055.10}
\]
This is the ordinary multiplication trace in a finite free algebra; its
definition and base-change law are EGA IV$_4$, Sections
18.2.1--18.2.2, printed p. 111, including (18.2.1.2).
It is not a Frobenius trace or a determinant norm.
The trace formula is valid even when the coefficient algebra
has nilpotents; the Boolean factor remains etale over that base.

At a node, the boundary section lies in the direct sum of the
balanced homogeneous modules
$\operatorname{Sym}^a(R^{r_I})\otimes
 \operatorname{Sym}^a(R^{s_I})$ for $0\le a\le|I|$, with the
variables viewed as the chosen coordinate generators.
The cross tensors give the integral divided-power contraction
(C055.5), and the transfers give functorial symmetric-power maps.
These are honest polynomial maps of finite free modules over
$\mathbb Z$; an exponential of derivatives with factorial
denominators is not required in positive characteristic.
Equivalently the permanent is the weight $(1,\ldots,1)$ coefficient
of $\prod_i\sum_j A_{ij}z_j$ for the diagonal column torus.
This character dictionary does not provide an unspecified
positive rule for GCT multiplicities or an algorithm for finding
a compatible cut tree.  The rank-state bound charges the entire
finite free boundary module, not only its geometric support.
The exact split-torus grading locator is SGA~3, Expose~I,
4.7.2--4.7.3, with the direct-sum/tensor-product Grothendieck-ring
interfaces SGA~6, Expose~IV, (1.5.2) and Definition~2.2, as
previously source-checked for NN-C038 and NN-S019.

Finally the generic family in the proof is a point of such a
parameter scheme with constant coarse width two.  Formula
(C055.9) says that its generic point is outside every presented
sparse-plus-rank locus with $s=n^{1+o(1)}$ and $r=o(n)$.
It is a transcendence-degree obstruction to that presentation,
not a claim that no small circuit computes the permanent.

\paragraph{Total-support leaf trace and matching modules.}
Over the nested-generator affine parameter scheme, the leaf
matching module is a direct sum indexed by the above heavy-row
statuses.  Each summand is the selected edge/row-monomer character
tensored with the rectangular residual matching module.  Its
character is an actual sum of matching monomials before any
arithmetic cancellations.  The residual character is computed as
the $t^{b'}$ coefficient of multiplication trace on the finite
etale Boolean scheme of rank $2^{b'}$, over the same truncated
boundary-coordinate algebra.  Thus (C055t.1) is an integral
presentation of the same regular boundary section and the same
global permanent pullback; it does not change the parameter
scheme or introduce unrelated original matrix coordinates.
The exact trace, split-torus grading and direct-sum/tensor-product
interfaces are the EGA/SGA locators retained in NN-C038 and NN-S019.
The only counting input is NN-C038 with $C=3$; its earlier
permanent heavy-row bound supplies the parameter comparison,
while the unmatched-row and rectangular cases are proved above.

\paragraph{Internal color pairing and matrix pullback.}
Replace the leaf matrix coordinates on the affine parameter scheme
by the morphism $(S_J,U_J,V_J)\mapsto S_J+U_JV_J$.
Equations (C055r.3)--(C055r.4) identify the pullback of the
leaf boundary section by an integral contraction on a finite
free module with two additional rank-$d_J$ color spaces.
The pairing is the canonical sum over bijections between
symmetric tensors in mutually dual color spaces.  On a repeated
one-color tensor of degree $a$ it has coefficient $a!$;
it is defined integrally and need not be perfect in positive
characteristic.  Simultaneous changes $U\mapsto UG$,
$V\mapsto G^{-1}V$ for $G\in\operatorname{GL}_d(R)$ preserve
the contraction and the resulting section.

The sparse leaf before contraction is still the specified
coefficient of multiplication trace on the original Boolean
finite etale scheme, over the larger truncated coordinate
algebra.  Every color and boundary coordinate is charged in
the finite free rank bound.  The outer compatible polynomial
matrix morphism and all-ring permanent identity remain unchanged.
Thus the theorem concerns a regular function on an explicitly
parametrized family and its effective matching-character module,
not an algorithm that chooses a hidden factorization or a
positive rule for arbitrary representation multiplicities.
The exact EGA trace and SGA tensor/character locators are those
already recorded for NN-C055, NN-C038 and NN-S019; the only
additional identity is the complete marked-matching argument above.

\paragraph{Local boundary modules and compatible capping.}
The node's actual boundary section is contained in its finite
free degree-truncated symmetric algebra of rank $N_I$; before
a merge the relevant module is the tensor product of the two
child truncations, rather than an unnecessarily enlarged total
simplex.  The integral contractions preserve that module and
the transfer maps land in the stated parent truncation.
Identity-subtree replacement is an explicit polynomial change
of compatible presentation over the same parameter scheme.
It preserves the permanent pullback and the Boolean trace
sections, and uses no assertion about existence of an unspecified
basis or an inverse over the coefficient ring.  The exact EGA,
SGA and integer-contraction source interfaces are those already
recorded for NN-C055.
\end{geometricstatement}
\begin{proof} The displayed dictionary transports the proof of NN-C055. Any additional realization, effectiveness or uniformity hypothesis is retained explicitly.\end{proof}


\section{NN-G001: geometric formulation and dictionary}
\begin{geometricstatement}[NN-G001]
Algebraic geometry and Grothendieck dictionary.:

This is an actual flat affine scheme over $\mathbb Z$, together with a
$\mathbb G_m$ action. The action, equivalently its grading, is the
extra datum in SGA 3, I, Proposition 4.7.3 and Corollary 4.7.3.1.
Forgetting it loses the Hilbert function. SGA 5, III, 4.7 is the
Lefschetz--Verdier trace framework for the local zeta; equality of
zeta functions alone is not an identification of cohomological
degrees, motives, or canonical bases. EGA II, Proposition 2.4.7(i),
additionally shows that taking a Veronese subring preserves
$\operatorname{Proj}$ as a scheme while changing the recorded grading.
There is no implication from the normality of this complex $D_5$
surface to saturation of the support of an arbitrary torus grading:
the example itself refutes that inference.


\end{geometricstatement}
\begin{proof} The displayed dictionary transports the proof of NN-G001. Any additional realization, effectiveness or uniformity hypothesis is retained explicitly.\end{proof}


\section{NN-G002: geometric formulation and dictionary}
\begin{geometricstatement}[NN-G002]
Scope and dictionary.:

The input is P-F31, \path{work/positivity_actual_cubic_pole.tex},
the theorem ``actual cubic pole'' and its denominator consequence;
P-F56 makes the two-adic valuation of a quasi-period exact. This is
an obstruction to an \emph{expanded} Hilbert-series representation.
The integer $M$ itself has only $n$ binary digits, so this proves no
exponential lower bound for arithmetic circuits, succinct rational
expressions, evaluation, or positivity. In geometric language these
are Hilbert characters of the graded invariant coordinate rings;
the grading/action dictionary is SGA 3, I, 4.7.3.1 and the projective
construction is EGA II, 2.4.1--2.4.7. No statement about the size of
a defining system or a Quot/Hilbert parameter space follows from
the length of this particular output format.


\end{geometricstatement}
\begin{proof} The displayed dictionary transports the proof of NN-G002. Any additional realization, effectiveness or uniformity hypothesis is retained explicitly.\end{proof}


\section{NN-G003: geometric formulation and dictionary}
\begin{geometricstatement}[NN-G003]
Geometric translation.:

The series is the Hilbert series of the graded
$k[x,y]$-module
$k[x,y]\oplus k[x,y](-L)^{\oplus(L-1)}$, with
$\deg x=1$, $\deg y=L$. Thus the same failure of inference occurs
for effective coherent graded modules, equivalently equivariant
modules under SGA 3, I, 4.7.3. The associated sheaf construction is
EGA II, Proposition 2.5.2. This uses a module, not a normal domain
or an actual GCT multiplicity algebra; no stronger geometric
realization is asserted.


\end{geometricstatement}
\begin{proof} The displayed dictionary transports the proof of NN-G003. Any additional realization, effectiveness or uniformity hypothesis is retained explicitly.\end{proof}


\section{NN-G004: geometric formulation and dictionary}
\begin{geometricstatement}[NN-G004]
Inputs, geometry and complexity.:

The inputs are ALPOGE\_GCT P-F56d and P-F57, with the exact wave
and dimension estimates in the supplied intake part 4 and the
paired-wave/strict-index proof files. A cyclotomic partial-fraction
expansion is a character computation in a localized representation
ring; the tensor and restriction operations themselves appear in
cote 2, scan 17, and symmetric-group operations in cote 133,
scan 28. A virtual-character cancellation is not a proof that
no different effective realization exists. The inequality gives
no Boolean or algebraic circuit lower bound and no obstruction
to exact integer arithmetic or symbolic cancellation. It makes
the ``near-supported zero'' precise only relative to a specified
expansion and norm.


\end{geometricstatement}
\begin{proof} The displayed dictionary transports the proof of NN-G004. Any additional realization, effectiveness or uniformity hypothesis is retained explicitly.\end{proof}


\section{NN-G005: geometric formulation and dictionary}
\begin{geometricstatement}[NN-G005]
Scope and geometric translation.:

This corrects the meaning of ``support linearization'' in
\path{03_split_zero_defect_preprint.tex}, its support section,
and \path{04_primitive_support_defects_preprint.tex}, the
primitive-support construction. A finite set of locally closed
strata gives an inclusion--exclusion valuation to its constructible
functions; assigning an Euler characteristic, a compactly supported
cohomology class, or a rank to each stratum is a further choice of
the weights $e_p\mapsto a_p$. For finite disjoint rational points
this is literal $K_0$ of their vector spaces. The external/direct
sum and tensor operations are the framework of cote 2, scan 17;
SGA 5, III, 4.7 concerns the subsequent trace realization, not a
canonical trace on the bare lattice. This theorem creates neither
witness multiplicities nor an effective counting algorithm from
Boolean support.


\end{geometricstatement}
\begin{proof} The displayed dictionary transports the proof of NN-G005. Any additional realization, effectiveness or uniformity hypothesis is retained explicitly.\end{proof}


\section{NN-G006: geometric formulation and dictionary}
\begin{geometricstatement}[NN-G006]
Inputs and exact dictionary.:

This repairs and strengthens the injectivity/pole statements in
\path{OURS_CONSULT_a157422a_COUNTING_CONSTRUCTION.md},
Theorem 4(a),(d),(e). The zeta and primitive odd Dirichlet functional
equations have the exact classical locators DLMF 25.4.1--25.4.4
and 25.15.5 (Apostol, \emph{Introduction to Analytic Number
Theory}, equations (12)--(16), pp.\ 259--263); A is required only when
one additionally locates their nontrivial zeros on the lines
$\operatorname{Re}s=j+1/2$. The geometric objects are Tate twists
of the trivial and Gaussian quadratic compatible rank-one systems.
For each $\ell$ their local systems are lisse on
$\operatorname{Spec}\mathbb Z[1/2\ell]$; the factor at $\ell$
is read using a different auxiliary prime. Their representation
ring is $\mathcal A$; SGA 3, I, 4.7.2--4.7.3 describes the
comodule/graded representation formalism. The Witt operations
are explicitly present in cote 43, scan 37. The local factor at
$2$ is omitted for tensor multiplicativity; restoring it as an
ordinary scalar correction does not make an unramified sign
character exist at the ramified prime. Faithful encoding is not
efficient decoding and does not prove PH1 or a counting bound.


\end{geometricstatement}
\begin{proof} The displayed dictionary transports the proof of NN-G006. Any additional realization, effectiveness or uniformity hypothesis is retained explicitly.\end{proof}


\section{NN-G007: geometric formulation and dictionary}
\begin{geometricstatement}[NN-G007]
Geometric consequences.:

At an inert odd prime the free orbit is
$\operatorname{Spec}\mathbb F_{p^2}$: it has no $\mathbb F_p$
points and two $\mathbb F_{p^2}$ points. Its geometric support,
its dimension, and its rational-point count are therefore distinct
invariants. Globally its Artin representation comes from
$\operatorname{Spec}\mathbb Z[i,1/2]$ and has zeta
$\zeta^{(2)}(s)L(s,\chi)$.
Changing $h$ in (G007.1) changes the Euler product by
\[
 \prod_j\zeta_{\mathbb Q(i)}^{(2)}(s-j)^{h_j}.
\]
A puts the nontrivial zeros of all these choices on their
appropriate lines, but selects none of the choices. Negative
Tate powers are legitimate twists of an $\ell$-adic object;
$q^{-1}$ itself is not the point count of a scheme over
$\mathbb F_q$. Thus an effective Artin--Tate representation is
not automatically an effective scheme count.
The exact archive comparison is cote 133, scan 19, where passing
from equivariant sets to representations must respect this
distinction; its subset/exterior-power formula already needs the
correction recorded in ALPOGE\_GCT PL-01. The group-scheme
representation formalism is SGA 3, I, 4.7.2. This gives no way to
recover satisfiability from a virtual cancellation or an
arbitrarily chosen dimension lift.


\end{geometricstatement}
\begin{proof} The displayed dictionary transports the proof of NN-G007. Any additional realization, effectiveness or uniformity hypothesis is retained explicitly.\end{proof}


\section{NN-G008: geometric formulation and dictionary}
\begin{geometricstatement}[NN-G008]
Correction and dictionary.:

Theorem 1(c) of the consult correctly distinguishes rational
descent from integrality; its Theorem 4(c) and final rank-four
summary do not preserve that distinction. A rational function
can also have an even-order pole at $v=1$, so ``even vanishing
order'' requires regularity there. For example rescaling by
$D-4$ gives poles. In the geometric dictionary a generic fibre
and a rational change of basis do not choose an integral lattice
or an effective class in $K_0$; SGA 3, I, 4.7.3 is a grading
equivalence, not an integrality theorem. Exterior powers and
effective direct summands, discussed in cote 133, scan 28 and
SGA 6, VI, Theorem 3.2, require actual vector-bundle classes.
The remaining rank-four gate includes a defined canonical basis,
integral Laurent constants, and an effective compatible lift.
The formal substitution $v=\sqrt q$ constructs neither a CM
Hecke character nor a compatible arithmetic family.


\end{geometricstatement}
\begin{proof} The displayed dictionary transports the proof of NN-G008. Any additional realization, effectiveness or uniformity hypothesis is retained explicitly.\end{proof}


\section{NN-G009: geometric formulation and dictionary}
\begin{geometricstatement}[NN-G009]
Inputs, geometry and complexity.:

The word-trace identity is GCT\_NSRH F42 and counting consult
Theorem 4(c). The new assertion quantifies the necessary size;
it does not rebrand F42 as new. A realization is $40$ copies of
$H^*((\mathbb P^1)^{2k})(k)$, with all cohomological degrees even.
It has a short tensor description even though its total rank is
exponential. Tensor products and trace are the operations of
cote 2, scan 17 and SGA 5, III, 4.7. This explains why a finite
list of multiplication objects can coexist with the absence of
a single bounded-rank Weil object. Neither a short tensor
description nor A gives a polynomial-size positive polytope
formula or rank-four canonical basis.


\end{geometricstatement}
\begin{proof} The displayed dictionary transports the proof of NN-G009. Any additional realization, effectiveness or uniformity hypothesis is retained explicitly.\end{proof}


\section{NN-G010: geometric formulation and dictionary}
\begin{geometricstatement}[NN-G010]
Geometric translation.:

For $X$ a disjoint union of $|B|$ rational points, matrices of
Tate--Artin objects are kernels on $X\times X$. Convolution is
$p_{13!}(p_{12}^*(-)\otimes p_{23}^*(-))$. This is the finite
zero-dimensional instance of the external product and projection
formula on cote 2, scan 17; Artin gradings use SGA 3, I, 4.7.2.
The construction refines counting consult Theorem 5 and
Corollary 8 by stating precisely which Grothendieck group is
proved. It uses the already known structure constants to build
the objects. It cannot establish their positivity without those
constants, nor establish PH0 or PH1 for a canonical geometric
construction.


\end{geometricstatement}
\begin{proof} The displayed dictionary transports the proof of NN-G010. Any additional realization, effectiveness or uniformity hypothesis is retained explicitly.\end{proof}


\section{NN-G011: geometric formulation and dictionary}
\begin{geometricstatement}[NN-G011]
Inputs, certificates and meaning.:

The exact input is counting consult Theorem 2 and
\path{OURS_E4_scripts/e3_constants.json}.
\path{verify_dimension_7276.py} independently builds the
table with integer Laurent dictionaries, checks all 1000
associativity triples, and verifies all coefficient inequalities
against the supplied 132 constants. It does not import the
supplied dimension-ring builder. The full certificate is
\path{dimension_certificate_7276.json}.
NN-G010 converts the result into a finite-kernel monoidal
category, with the same SGA 3/cote 2 dictionary. This lowers
an explicit upper bound for the noncanonical cancellation
cost. It neither selects a canonical lift, proves global
optimality, nor decides PH0/PH1 at rank four. No unimodality
claim is made for this particular smaller lift.


\end{geometricstatement}
\begin{proof} The displayed dictionary transports the proof of NN-G011. Any additional realization, effectiveness or uniformity hypothesis is retained explicitly.\end{proof}


\section{NN-G012: geometric formulation and dictionary}
\begin{geometricstatement}[NN-G012]
Geometry and scope.:

The full sixteen-product certificate is
\path{verify_signpoint_matrix_block.py} and
\path{signpoint_matrix_certificate.json}. The four-dimensional
kernel algebra on a two-point scheme is again the convolution
algebra of cote 2, scan 17; (G012.1) is its explicit basis change.
The trace cyclicity is the finite-dimensional case of SGA 6,
I, section 8 and SGA 5, III. At the sign point, the cyclic
permutation mixes four failed generic coefficients with two
shortest coefficients, as E4 section 6b records. Accordingly
this explains the restored trace symmetry, not generic
unimodality or the existence of a rank-four canonical basis.


\end{geometricstatement}
\begin{proof} The displayed dictionary transports the proof of NN-G012. Any additional realization, effectiveness or uniformity hypothesis is retained explicitly.\end{proof}


\section{NN-G013: geometric formulation and dictionary}
\begin{geometricstatement}[NN-G013]
Exact inputs, signed failure and complexity scope.:

The unary hardness locators are Fischer--Ikenmeyer,
Theorem 3.1, label \path{thm:plethysm-positivity-np-hard},
and Ikenmeyer--Mulmuley--Walter, \texttt{tkronnp} in the
supplied main source. The latter's positive formula for
an NP-hard subclass is \texttt{tkronpos}, so positivity of
a formula is not itself a tractability theorem.
For a signed sum isolation tests the presence of summands:
on $S_2$, $g((1,1),(1,1),(1,1))=(1-1)/2=0$, although
both character summands exist. A and M do not supply the
hashed parity oracle. E can evaluate an explicitly encoded
witness count after a parsimonious bounded-width CNF
encoding, but its exponent is measured in that encoding's
variables, including verifier and hash auxiliary variables;
the word ``polynomial'' does not preserve a linear exponent.
There is no derivation of NP=RP from the stated inputs.


Geometric translation.:

The witness relation and affine hashes are equations over
$\mathbb F_2$ on Boolean coordinates; verification constraints
give a finite affine scheme whose rational points are the
witnesses after parsimonious encoding. Its point-count
parity is additional information beyond support of its
terms. SGA 5, III, 4.7 interprets counts as alternating
Frobenius traces but does not evaluate them in polynomial
time. The finite-set/representation distinction is
NN-G007 and cote 133, scan 19. An arbitrary split-zero
coefficient formula is not such a witness scheme until
this identification is constructed.


\end{geometricstatement}
\begin{proof} The displayed dictionary transports the proof of NN-G013. Any additional realization, effectiveness or uniformity hypothesis is retained explicitly.\end{proof}


\section{NN-G014: geometric formulation and dictionary}
\begin{geometricstatement}[NN-G014]
Reduced gate and dictionary.:

For a GCT positivity algorithm this leaves three precise
tasks: a coefficient-preserving localization, polynomially
bounded certificates for the target exposure and residue,
and a procedure finding such a certificate for every yes
instance. A deterministic polynomial-time complete procedure
would yield P=NP on the known hard families; randomized
one-sided completeness yields NP=RP. A bare unique smallest
signed term does neither. The source is the question left
by the GCT support consult, section on P-F30/P-F23 and
reverse isolation. Geometrically a one-parameter subgroup
of a torus filters an equivariant module by weights; the
associated graded residue retains the exposed piece. The
exact action/graded-module dictionary is SGA 3, I, 4.7.3;
the projective/graded localization construction is EGA II,
2.2--2.5. No degeneration has been asserted to preserve
the desired multiplicity without the coefficient identity.


\end{geometricstatement}
\begin{proof} The displayed dictionary transports the proof of NN-G014. Any additional realization, effectiveness or uniformity hypothesis is retained explicitly.\end{proof}


\section{NN-G015: geometric formulation and dictionary}
\begin{geometricstatement}[NN-G015]
Inputs, border scope and geometry.:

BIP, \emph{No occurrence obstructions in geometric
complexity theory}, Theorem 1.1, gives its stated
$n\ge m^{25}$ support inclusion for padded permanent
and determinant orbit closure. It does not prove
containment and does not eliminate multiplicity
obstructions or other equations. Orbit closure is the
geometric object for border complexity; confusing it
with the orbit or the support set would discard the
defining ideal. In Grothendieck language the needed
object is the equivariant closed subscheme, its
quotient coordinate algebra, and its graded locally
free representation pieces. Cote 133, scan 28 supplies
the Schur/tensor operations on their $K_0$ classes;
cote 2, scan 17 supplies restriction and tensor products.
No FGA representability theorem by itself computes the
ideal or the multiplicity comparison. The example in
this proof can be realized by finite Boolean witness
schemes with trivial group action, using succinct
equations; its explicit vector space may have
exponentially many basis vectors. Thus it makes no
unproved reduction into the BIP orbit-closure family.


\end{geometricstatement}
\begin{proof} The displayed dictionary transports the proof of NN-G015. Any additional realization, effectiveness or uniformity hypothesis is retained explicitly.\end{proof}


\section{NN-G016: geometric formulation and dictionary}
\begin{geometricstatement}[NN-G016]
Source correction, scope and geometry.:

The explicit model is the one in corpus paper 03, lines 850--998,
and paper 04, lines 358--504, with $\ell=\log p$. Their scalar
$h$ cannot by itself lift a rank-one projection: its boundary
value is a multiplier unit in a nonunital algebra. The replacement
$hP$ removes precisely that gap. The classical analytic inputs
are Blackadar, \emph{K-Theory for Operator Algebras}, second edition,
\S9.3.2 (printed page 68), for the exponential connecting map,
and Williams, \emph{Crossed Products of C*-Algebras}, author draft
version 3.1 (6 September 2006), Corollary 4.11 and Remark 4.12
(printed page 126), for the free proper action Morita equivalence.
Both locators were checked in the primary author PDFs.

The group action on all of $Y$ is also free and proper: compact
sets have bounded $y$ coordinates, so only finitely many translates
can meet them. Thus the same free proper action theorem applies
to $E$, compatibly with its open and closed pieces. Transferring
the calculation to a larger arithmetic quotient still requires the
claimed identification of its edge extension with this exact model.
No verification of that global identification is smuggled into
the present model theorem, and no purity or complexity bound follows.

An honest algebraic-geometric comparison is available for the
scalar loop: $(t+i)/(t-i)$ is a rational function on
$\mathbb P^1_{\mathbb C}$ with divisor $[-i]-[i]$ and unit norm
on the real locus. Equivalently its real coordinates are
$((t^2-1)/(t^2+1),2t/(t^2+1))$ on the norm-one torus
$x^2+y^2=1$; the removed degree-two divisor $t^2+1=0$ has no
real points. The zero principal divisor class in the Picard
group and the nonzero winding on the real locus are different
invariants. In Grothendieck language a rank-one corner is an
actual rank-one object before taking $K_0$, not the literal unit
of a multiplier algebra. The supplied archive cote 2, scan 17
provides the analogous pullback/tensor formalism for geometric
objects; it does not identify this analytic $K_1$ class with an
\'{e}tale trace or Weil weight. There is no claimed EGA/SGA theorem
turning the crossed-product model into an algebraic scheme.


\end{geometricstatement}
\begin{proof} The displayed dictionary transports the proof of NN-G016. Any additional realization, effectiveness or uniformity hypothesis is retained explicitly.\end{proof}


\section{NN-G017: geometric formulation and dictionary}
\begin{geometricstatement}[NN-G017]
Geometric dictionary and consequences.:

This is already a theorem about a fixed smooth normal scheme
over $\mathbb Z$ with varying diagonal group-scheme actions.
SGA 3, I, Proposition 4.7.3 and Corollary 4.7.3.1 identify
the actions with the coordinate-ring gradings. The saturated
support is $\mathbb N$: the group generated by $L,L+1$ is
$\mathbb Z$ and their rational cone is $\mathbb R_{\ge0}$.
For the associated semigroup curve
$\operatorname{Spec}\mathbb Z[t^L,t^{L+1}]$, the finite
normalization is $\operatorname{Spec}\mathbb Z[t]$ and the
conductor ideal in $\mathbb Z[t]$ is
$t^{L(L-1)}\mathbb Z[t]$. Indeed $t$ is integral because
$t^L$ belongs to the subring, the two fraction fields agree
because $t=t^{L+1}/t^L$, and $\mathbb Z[t]$ is integrally
closed. The conductor calculation follows termwise from the
already proved description of the holes. This semigroup
curve is an associated encoding of the support; it is not
the smooth scheme $X$ itself.

A makes the nontrivial zeros of the common arithmetic zeta
lie on $\operatorname{Re}s=5/2$, but it supplies no missing
weight spaces and changes none of the holes. A fixed finite
pure Weil realization of an \emph{ungraded} count, even with
its induced connected-group action, therefore cannot stand
in for the grade-indexed trace or witness objects needed by
NSRH and PH1. The two-variable formula here already has
the desired positive counting form; this example refutes a
proposed inference about conductor bounds, not PH1 itself
and not a complexity conjecture about actual plethysm.


\end{geometricstatement}
\begin{proof} The displayed dictionary transports the proof of NN-G017. Any additional realization, effectiveness or uniformity hypothesis is retained explicitly.\end{proof}


\section{NN-G018: geometric formulation and dictionary}
\begin{geometricstatement}[NN-G018]
Geometric dictionary and exact limit of the conclusion.:

The input is P-F31, theorem \texttt{actualcubicpole}, part two,
in \path{positivity_actual_cubic_pole.tex}; its family is
the Cartan multiplicity algebra of cubic $SL_n$ invariants.
Its variable $t$ records the grading in SGA 3, I, 4.7.3.1.
To reinterpret $f_n(r)$ as an extension-field trace in the
SGA 5, III, 4.7 framework requires a new identification of
that grading index with the power of Frobenius. If such an
identification is proposed using one finite cohomological
carrier, (G018.1) bounds its total rank. Neither a Tate twist,
an inserted correspondence nor virtual cancellation removes
the pole count; insertions and signs were already allowed.
It is essential to distinguish this from the Tate--Artin
realizations of individual Laurent-polynomial structure
constants in NN-G006--NN-G011, where the field size and
the basis-word or stretching parameter are separate variables.

The rank obstruction is unconditional relative to the cited
refutation and is unaffected by A or conditional M. It
rules out a polynomial-rank iteration carrier on this actual
plethysm family, but leaves short tensor descriptions,
varying graded carriers, succinct Hilbert rational functions,
and PH1 positive witness formulas open. No arithmetic-circuit
or Boolean lower bound follows from it alone.


\end{geometricstatement}
\begin{proof} The displayed dictionary transports the proof of NN-G018. Any additional realization, effectiveness or uniformity hypothesis is retained explicitly.\end{proof}


\section{NN-G019: geometric formulation and dictionary}
\begin{geometricstatement}[NN-G019]
Geometric translation and smaller remaining gate.:

By NN-G010 both minimizing rings give additive monoidal
categories of kernels on a fixed ten-point scheme with
Tate--Artin coefficient objects. The score counts the added
regular $C_2$ pairs in the structure-constant objects,
coefficient by coefficient. The relevant actual-object and
split-$K_0$ distinction is SGA 3, I, 4.7.2 and archive cote 2,
scan 17; effectivity is not equality in a virtual representation
ring. This optimization selects no canonical geometric functor.

Any improvement below $7276$ must now leave at least one
specific part of (G019.1): use unequal cell $\delta_i$,
a nonmonic or higher-degree $\delta$, $\deg W>2$,
nonintegral or nonconstant ratios $G_i/(W\delta)$,
allow $\Lambda_i$ below the true $\lambda_i$ in some
coefficients while compensating elsewhere in the table,
or change the table shape. These are genuine remaining
possibilities. The result narrows the search gate but
does not prove global minimality, canonical positivity,
PH1, or anything about rank four by itself.


\end{geometricstatement}
\begin{proof} The displayed dictionary transports the proof of NN-G019. Any additional realization, effectiveness or uniformity hypothesis is retained explicitly.\end{proof}


\section{NN-G020: geometric formulation and dictionary}
\begin{geometricstatement}[NN-G020]
Geometric dictionary and the exact failure of zeta recovery.:

The torus action and its multigraded comodule are the objects of
SGA 3, I, 4.7.2--4.7.3 and Corollary 4.7.3.1. Equation (G020.7)
is an actual affine quotient by the indicated split subtorus;
the finite generation and normality proofs above make its scheme
meaning explicit. The ray grading also gives $\operatorname{Proj}R_F$
and its associated graded-module sheaves in the sense of EGA II,
2.4--2.5, especially 2.4.7 and 2.5.2. No identification with a
plethysm or Kronecker multiplicity algebra is asserted.

Put $N=2n+m$. The underlying \emph{ambient} scheme
$\operatorname{Spec}S=\mathbb A^N_{\mathbb Z}$ is independent of
the formula once $n,m$ are fixed. For every prime power $q$ and
integer $r\geq1$ it has $q^{Nr}$ points over $\mathbb F_{q^r}$.
Its only nonzero compact-support $\ell$-adic cohomology group
is $H_c^{2N}=\mathbb Q_\ell(-N)$, and the connected torus action
induces the trivial action on that group. The precise trace and
connected-action locator is SGA 4, XVIII, 2.8.1 and the immediately
following paragraph; the affine-space calculation also follows
from XVII, 5.4.3 and the curve trace calculation XVIII, 1.1.6--1.1.9.
Consequently
\[
 Z(\mathbb A^N_{\mathbb F_q},T)=(1-q^NT)^{-1},\qquad
 \zeta_{\mathbb A^N_{\mathbb Z}}(s)=\zeta(s-N).
 \tag{G020.9}
\]
Under A its nontrivial zeros lie on $\Re s=N+1/2$. These common
point counts and this rank-one pure Weil object do not recover
the degree-one rank, despite the normal ray algebra, known rational
half-point, positive monomial rule and conductor bound two.
For an explicit smallest comparison, the formulas
$(x\vee x\vee x)\wedge(x\vee x\vee x)$ and
$(x\vee x\vee x)\wedge(\neg x\vee\neg x\vee\neg x)$
have degree-one ranks one and zero and the same ambient
$\mathbb A^4_{\mathbb Z}$.

This last assertion concerns the ungraded ambient realization;
no formula-independent zeta or cohomology is claimed for
$\operatorname{Spec}R_F$ itself, or for grade-indexed trace objects.
Likewise the conductor here is an actual support conductor,
not GCT VI's strict off-residue saturation index. The construction
does not contradict a conjecture confined to actual plethysm or
Kronecker families. It proves that normality, a homogeneous
lattice-point model, known moment-cone membership, a uniformly
tiny actual conductor, and a finite pure ambient realization
still leave an NP-hard first-coefficient gate in this explicit
geometric family. Neither A nor the unrelated conditional M
removes that gate; E supplies precisely (G020.8).
\hfill$\square$


\end{geometricstatement}
\begin{proof} The displayed dictionary transports the proof of NN-G020. Any additional realization, effectiveness or uniformity hypothesis is retained explicitly.\end{proof}


\section{NN-G021: geometric formulation and dictionary}
\begin{geometricstatement}[NN-G021]
Geometric meaning and the remaining gate.:

NN-G010 turns each of these effective multiplication tables
into the same type of additive monoidal kernel category on
a ten-point scheme, with effective Artin--Tate multiplicity
objects. Only the actual structure constants specify those
objects. An auxiliary factor such as $\Lambda_2$ need not
be the class of an effective object even when its product
with $\delta$ is. This is the precise distinction between
an effective tensor structure and a chosen polynomial
factorization of its coefficients; no cancellation history
can be read from a supported-zero tag alone.
The comodule dictionary is SGA 3, I, 4.7.2; kernel products
are the construction recorded in archive cote 2, scan 17.
Neither source canonically selects these tables.

The class bound now allows signed $\Lambda_i$ and rational
constant ratios. A score below $7064$ must leave another
specified restriction: a common monic quadratic $\delta$,
degree at most two for $W$, constant ratios $G_i/(W\delta)$,
or the closed table shape. This is a smaller global
optimization gate than NN-G019. It remains a question
about noncanonical realizations, with no consequence for
PH0, PH1 or a rank-four canonical basis.


\end{geometricstatement}
\begin{proof} The displayed dictionary transports the proof of NN-G021. Any additional realization, effectiveness or uniformity hypothesis is retained explicitly.\end{proof}


\section{NN-G022: geometric formulation and dictionary}
\begin{geometricstatement}[NN-G022]
Geometric interpretation and the reduced remaining gate.:

This is a rigidity consequence of the parity of actual effective
Artin--Tate multiplicity objects and integrality of the closed
kernel-composition table. The equality modulo two is the condition
that dimension and signed trace arise from one representation,
as in NN-G007 and SGA 3, I, 4.7.2. The integer roots here are
roots of auxiliary parameter polynomials; they are not asserted
to be Frobenius eigenvalues or points of an arithmetic scheme.
Archive cote 2, scan 17 supplies the honest kernel-product
dictionary through NN-G010, not an arithmetic interpretation
of those roots.

Within monic quadratic cells and positive rational constant
ratios, any realization competitive with $7064$ must have
$\delta_1=\delta_2$, even for higher-degree $W$. The remaining
optimization possibilities are higher-degree $W$ in that common
quadratic class, nonmonic or nonquadratic cells, nonconstant
ratios, or a different table shape. This improves the exact
gate without introducing a canonical construction or a
nonstandard purity theorem.


\end{geometricstatement}
\begin{proof} The displayed dictionary transports the proof of NN-G022. Any additional realization, effectiveness or uniformity hypothesis is retained explicitly.\end{proof}


\section{NN-G023: geometric formulation and dictionary}
\begin{geometricstatement}[NN-G023]
Geometric dictionary and remaining global scope.:

The fourth-root evaluation is the trace of a finite
character of the grading: restrict a $\mathbb G_m$
representation to $\mu_4$ and evaluate its character at
$i$. This is the diagonalizable-group/comodule formalism
of SGA 3, I, 4.7.2--4.7.3. It is not a finite-field
cardinality, a Frobenius eigenvalue, or a zero of an
arithmetic zeta function. Effective graded objects may
have zero trace at this character. What fails in the
hypothetical small-score table is the required divisibility
of one effective character by the specified cell character.
The actual effective multiplicity objects and their
kernel products remain those of NN-G010 and archive
cote 2, scan 17.

The arbitrary-degree gluing gate is closed for monic
quadratic cells and rational constant ratios. A smaller
score must use nonmonic or nonquadratic cells, nonconstant
ratios, or another
table shape. None of these possibilities is dismissed.
The optimum remains noncanonical and has no implication
for a rank-four canonical basis, PH0, PH1 or the missing
nonstandard-purity bounds.


\end{geometricstatement}
\begin{proof} The displayed dictionary transports the proof of NN-G023. Any additional realization, effectiveness or uniformity hypothesis is retained explicitly.\end{proof}


\section{NN-G024: geometric formulation and dictionary}
\begin{geometricstatement}[NN-G024]
Geometric translation and the remaining gate.:

NN-G010 realizes the integral effective multiplicities of
(G024.2) by actual graded sign representations in an
additive monoidal kernel category on a ten-point scheme.
The exact dictionaries remain SGA 3, I, 4.7.2--4.7.3
for graded representations and archive cote 2, scan 17
for finite-kernel composition. There is no object of
fractional dimension attached to $\Lambda_2$; the rational
factor is merely a coordinate factorization in $\mathbb Q(D)$.
The structure constant $\delta\Lambda_2$ is the actual
integral object. The content-three factor is three copies
of its effective cell object, not division by three in
an effective Grothendieck semiring.

For the equal-cell, constant-ratio tables realized by these
chosen kernel generators, the minimum sum over the $132$
coefficient positions and grading pieces of half the excess of dimension over
absolute signed trace is exactly $5932$. This is a minimum
for the specified character-valued multiplication data and
its NN-G010 realizations, not for all monoidal categories
or all generating sets with the same Grothendieck ring.

The same supplied trace algebra therefore admits smaller
composing effective lifts after a purely auxiliary
integrality restriction is removed. This demonstrates
why formal factor parameters cannot be identified with
canonical count objects. It establishes neither a
canonical family nor a rank-four basis, PH0, PH1 or
nonstandard purity. The precise remaining lower-score
gates are unequal general cells, nonconstant ratios,
or a different multiplication-table shape; arbitrary
common-cell degree and content are now covered.


\end{geometricstatement}
\begin{proof} The displayed dictionary transports the proof of NN-G024. Any additional realization, effectiveness or uniformity hypothesis is retained explicitly.\end{proof}


\section{NN-G025: geometric formulation and dictionary}
\begin{geometricstatement}[NN-G025]
Geometric meaning and residual scope.:

The two cells have the same primitive grading character
and different positive multiplicities. Their actual
coefficient objects still lie in the effective graded
sign-representation semiring, and NN-G010 supplies the
same finite-kernel monoidal realization. The precise
classical dictionaries are SGA 3, I, 4.7.2--4.7.3 and
archive cote 2, scan 17. Rational $\Lambda_2$ remains
an auxiliary factor, not a fractional effective object.
Different contents of the two cell characters do not
alter the supplied trace algebra whose effective lift
is being sought; they change the chosen dimension data.

Among the NN-G010 kernel realizations whose two cell characters
are positive integral multiples of the same primitive character
and whose specified ratios are constant, the least total
half-excess of dimension over absolute signed trace is $4520$,
with $3532$ in the bottom-cell coefficients for (G025.1).
These numerical invariants refer to the chosen table and
generators; they are not invariants of every categorification
of the abstract trace algebra.

The exact common-cell minimum and the exact proportional-cell
minimum differ, so equality of cell characters is a
substantive optimization restriction, not a consequence
of the trace constraints. A smaller score must now use
nonproportional cell polynomials, nonconstant ratios,
or another table shape. No canonicality, generic
rank-four positivity, PH0/PH1 or nonstandard purity
conclusion is inferred from this noncanonical optimum.


\end{geometricstatement}
\begin{proof} The displayed dictionary transports the proof of NN-G025. Any additional realization, effectiveness or uniformity hypothesis is retained explicitly.\end{proof}


\section{NN-G026: geometric formulation and dictionary}
\begin{geometricstatement}[NN-G026]
Geometric dictionary and the smaller remaining gate.:

The parity and modulo-four arguments are constraints
on actual effective graded sign representations:
dimension and signed trace agree modulo two, while
the composing table imposes further integral divisions.
SGA 3, I, 4.7.2--4.7.3 provides the diagonalizable-group
character dictionary, and the NN-G010 realization is
the finite-kernel composition recorded in archive
cote 2, scan 17. The evaluations in (G026.7) restrict
the grading character to finite roots of unity;
they are not Frobenius counts or arithmetic zeta zeros.
No rational auxiliary factor is promoted to an
effective object unless its actual coefficient has
been verified integral and nonnegative.

In geometric terms, all effective graded multiplicity objects
for the two cells may now vary independently, subject to this
closed kernel-table presentation and its constant rational
ratio identities. Among the resulting NN-G010 realizations,
the least total half-excess of dimension over the absolute
signed trace coefficients is exactly $4520$. The trace
specialization is fixed, but no canonical choice of objects,
coherence data, or generators is claimed. This optimum is
presentation-dependent, not a universal categorical invariant.

Arbitrary nonproportional cell polynomials and
unbounded support locations have now been treated
within the constant-ratio closed-table class. A
smaller score must use nonconstant ratios
$G_i/(W\delta_i)$ or a different table shape.
This is an exact realization theorem with explicit
scope, not a canonical-basis theorem or a consequence
for PH0, PH1, rank four, or nonstandard purity.


\end{geometricstatement}
\begin{proof} The displayed dictionary transports the proof of NN-G026. Any additional realization, effectiveness or uniformity hypothesis is retained explicitly.\end{proof}


\section{NN-G027: geometric formulation and dictionary}
\begin{geometricstatement}[NN-G027]
Geometric meaning and residual scope.:

The actual polynomials in (G027.2) are effective
Artin--Tate multiplicity objects and define another
NN-G010 finite-kernel monoidal realization of the
same trace table. With its chosen generators, the
sum over multiplication positions and grading pieces
of half the excess of dimension over absolute signed
trace is $3452$ for the first construction, and $3352$
for (G027.4), with bottom-cell contributions $2852$
and $2752$, respectively.
This is an explicit upper bound for a presentation-dependent
quantity, not a universal invariant of all categories
realizing the abstract algebra.

The rational $\Lambda_i$ and the two nonconstant ratios
are generic coordinate factorizations; no object of
fractional dimension or meromorphic sheaf is inferred.
Only the verified actual coefficients define objects.
The precise dictionary is SGA 3, I, 4.7.2--4.7.3 and
archive cote 2, scan 17 through NN-G010. The example
leaves the exact optimum for arbitrary rational ratios
and for other table shapes open. It does not give a
canonical family, a rank-four basis, PH0/PH1 or
nonstandard purity.


\end{geometricstatement}
\begin{proof} The displayed dictionary transports the proof of NN-G027. Any additional realization, effectiveness or uniformity hypothesis is retained explicitly.\end{proof}


\section{NN-G028: geometric formulation and dictionary}
\begin{geometricstatement}[NN-G028]
Geometric translation and the next exact gate.:

Fixing (G028.1) fixes three actual multiplicity objects
in the NN-G010 finite-kernel presentation. Among all
its effective graded sign-representation tables with
those fixed objects and the specified trace specialization,
the sum over multiplication positions and grading pieces
of half the dimension-minus-absolute-trace excess is
exactly $3352$. There are two minimizing character-valued
parameter sets in this presentation; no uniqueness of
categories or coherence data is asserted. This is an
optimization of chosen multiplication data, not an
intrinsic minimum for every categorification of the
abstract trace algebra.

The exact source dictionary remains SGA 3, I,
4.7.2--4.7.3 for effective grading characters and
archive cote 2, scan 17 for kernel composition.
The mass reduction respects actual coefficient
objects and never treats a rational auxiliary factor
as an effective object. A smaller closed-table score
must change at least one of $W,U_1,U_2$ in (G028.1);
other table shapes are a further separate possibility.
No canonicality, rank-four, PH0/PH1 or purity theorem
is inferred from this exact restricted optimum.


\end{geometricstatement}
\begin{proof} The displayed dictionary transports the proof of NN-G028. Any additional realization, effectiveness or uniformity hypothesis is retained explicitly.\end{proof}


\section{NN-G029: geometric formulation and dictionary}
\begin{geometricstatement}[NN-G029]
The finite algorithm and its exact scope.:

There are only finitely many arrays of nonnegative integer
Laurent coefficients with total mass at most $N$ and
exponents in $[-B,B]$. Enumerate them, impose bar symmetry,
the specified unit, coefficient domination and parity,
and test associativity by exact finite polynomial
arithmetic. Optional prescribed zero positions are
retained throughout. The known incumbent ensures the
accepted set is nonempty. Minimize the cost over it.

Any better unbounded-degree lift would compress to an
accepted lift of the same cost in this finite set.
Therefore the computed minimum is the unrestricted
symmetric-lift minimum for the chosen basis and unit,
not merely a degree-window minimum. The search bound is
immense: neither polynomial time nor any useful running
time estimate follows from this argument.

For the numerical specialization, NN-G011 proves the
absolute nonunit trace mass $4424$, and the nineteen
unit positions add $19$, giving $N_0=4443$.
The supplied trace coefficient table has maximum
Laurent exponent five. The checked NN-G027 lift has
mass $11128+19=11147$. Hence $n_*\le5578$ and
$n_*(K+1)4^{n_*-1}\le33468\,4^{5577}<2^{11170}$.
The small certificate
\path{verify_degree_cutoff_constants.py} checks
these input norms and constants; it does not perform
the enormous finite search or substitute for the proof.


Geometric translation and the reduced optimization gate.:

The exponent variables are characters in a common finite
presentation of a diagonalizable group. More explicitly,
start with a free character group on these symbols and
one marked symbol $e_0$. Impose the homogeneous character
relations corresponding to the exponent collisions and
the anchor relations $y_r-r e_0=0$. An integral cocharacter
to this diagonalizable group, with value one on $e_0$,
assigns the integer grading weights. The strict order
conditions select an open rational polyhedral region
in that affine cocharacter slice. The original and the
compressed solutions are two integral cocharacters in
the same finite collision region, both fixing the marked
trace weights.

SGA 3, I, 4.7.2--4.7.3 identifies these character
presentations with diagonalizable-group comodules.
Pulling back their finite multiplicity objects along
the two cocharacters preserves the displayed finite
multiplication relations; NN-G010 and archive cote 2,
scan 17 then give the corresponding finite-kernel
realizations. Thus the actual conclusion is existence
of an equal-cost realization with bounded grading
weights and the same specified trace table. No equivalence
of the two categories or canonical choice of grading
is asserted.

The reassignment is not a homomorphism of the original
whole weight lattice $\mathbb Z$ to itself, and is not
a single Veronese scaling: the weight-one anchor remains
one. It preserves the finite character relations needed
for this multiplication presentation. The unrestricted
symmetric-lift optimum is now a finite effective search
problem, and unary-mass-budget existence has polynomial-size
verifiable certificates. Its value and any practical
optimization method remain open. Nothing here computes actual GCT multiplicities,
proves PH0/PH1, or transfers a nonstandard purity statement.


\end{geometricstatement}
\begin{proof} The displayed dictionary transports the proof of NN-G029. Any additional realization, effectiveness or uniformity hypothesis is retained explicitly.\end{proof}


\section{NN-G030: geometric formulation and dictionary}
\begin{geometricstatement}[NN-G030]
Geometric translation and algorithmic scope.:

At the level of integer congruences, individual effective
graded sign representations require dimension and signed
trace to agree modulo two.
Associativity of their finite-kernel composition imposes
the additional congruences (G030.1). In the symmetric
character ring $\mathbb Z[D]$, they say that the listed
coefficient characters vanish in the explicit quotient
rings $\mathbb Z[D]/(4,2D)$,
$\mathbb Z[D]/(4,2D,D^2)$, and
$\mathbb Z[D]/(4,2D^2(D-1))$, as appropriate.
These are exact scheme-theoretic congruence statements
over the nonreduced characteristic-four base, not
purity or point-counting assertions.

The character interpretation is SGA 3, I, 4.7.2--4.7.3;
$D=4$ is evaluation at the identity grading character,
while $D=0$ is evaluation at $q=-1$.
The composing objects remain the NN-G010 kernel
realization of archive cote 2, scan 17. These conditions
provide inexpensive necessary tests before a bounded
search or NP verification, but are not sufficient for
an effective lift and do not determine its optimum,
canonicality, rank-four positivity, or PH0/PH1.


\end{geometricstatement}
\begin{proof} The displayed dictionary transports the proof of NN-G030. Any additional realization, effectiveness or uniformity hypothesis is retained explicitly.\end{proof}


\section{NN-G031: geometric formulation and dictionary}
\begin{geometricstatement}[NN-G031]
Geometric translation and scope.:

In the NN-G010 realization by graded sign-representation
kernels, $w$ is the ordinary dimension of the gluing
coefficient object. The theorem says that any such
closed-table realization costing at most $3352$
regular-sign pairs must use a gluing object of dimension
at least eight. A dimension-six gluing object has the
unique character $D+2$; its associative kernel products
would force the incompatible character congruences and
mass bounds proved above. This is an actual dimension
obstruction for the effective realization, not a statement
about Frobenius weights or point counts. The character
dictionary is SGA 3, I, 4.7.2--4.7.3; the composition
dictionary is archive cote 2, scan 17, as realized in
NN-G010. The result shrinks the remaining optimization
gate but gives no global optimum or canonical GCT basis.



\end{geometricstatement}
\begin{proof} The displayed dictionary transports the proof of NN-G031. Any additional realization, effectiveness or uniformity hypothesis is retained explicitly.\end{proof}


\section{NN-G032: geometric formulation and dictionary}
\begin{geometricstatement}[NN-G032]
Geometric translation and limitations.:

For a fixed coefficient pattern, the grading weights lie
in the cocharacter lattice of a finitely presented
diagonalizable group, with the associativity collisions
imposing integral linear relations. Domination and the
chosen order of weights specify finitely many lattice
chambers. Eliminating exponent variables projects these
integer lattice conditions; the resulting congruences
record finite-index lattice obstructions. Their character
duals are finite diagonalizable groups such as $\mu_m$.
The actual conclusion is that, for bounded total dimension
mass, admissibility of the prescribed trace weights is
decidable by finitely many such linear inequalities and
congruences, computed uniformly from the mass budget.

This is an arithmetic statement about cocharacter lattices,
not a claim that inequalities define Zariski-open subsets
or that an effective lift is a canonical geometric object.
The grading dictionary is SGA 3, I, 4.7.2--4.7.3, and
the actual composing kernel realization is NN-G010,
using archive cote 2, scan 17. The algorithm gives no
polynomial bound on the mass budget needed by varying
GCT inputs, no practical small value of $F$, and no
rank-four positivity, purity, PH0 or PH1 theorem.



\end{geometricstatement}
\begin{proof} The displayed dictionary transports the proof of NN-G032. Any additional realization, effectiveness or uniformity hypothesis is retained explicitly.\end{proof}


\section{NN-G033: geometric formulation and dictionary}
\begin{geometricstatement}[NN-G033]
Geometric translation and limitations.:

The effective Artin--Tate coefficient objects in
(G033.2), with the NN-G010 kernel-composition realization,
give an actual additive monoidal model of the same
graded sign trace. The excess of ordinary dimension
over absolute signed trace is precisely twice $2744$
regular-sign pairs in its chosen multiplication
presentation, with $1296$ pairs in the bottom cell.
The gluing character has dimension eight and includes
weight-two characters absent from $D$.

This construction concerns actual coefficient objects,
not formal $K_0$ cancellation or a count obtained by
evaluating a virtual class. Its class is governed by
SGA 3, I, 4.7.2--4.7.3 and the archive cote 2, scan 17
composition dictionary implemented in NN-G010.
The rational auxiliary $\Lambda_2$ does not represent
a fractional-dimensional object. The global minimum,
canonicality, rank-four extensions and PH0/PH1 remain
separate questions. The older exact subclass bounds
and the NN-G031 implication for all costs at most
$3352$ remain valid with their stated hypotheses.



\end{geometricstatement}
\begin{proof} The displayed dictionary transports the proof of NN-G033. Any additional realization, effectiveness or uniformity hypothesis is retained explicitly.\end{proof}


\section{NN-G034: geometric formulation and dictionary}
\begin{geometricstatement}[NN-G034]
Geometric translation and remaining gate.:

Fixing (G034.1) fixes the first-cell coefficient object,
its bottom action character, and one gluing product in
the NN-G010 category of graded sign-representation kernels.
Among all effective associative completions in this closed
presentation, the least total excess of dimension over
absolute trace is exactly $2\cdot2744$, achieved by the
single set of coefficient characters of NN-G033.
The proof allows arbitrary grading locations in the other
objects; the mass-ten $D-2$ factor is treated by an actual
integral divisibility argument, not discarded as generic.

The finite-character evaluations $D=0$ and $D=2$
are restrictions at $q=-1$ and $q=i$, under SGA 3,
I, 4.7.2--4.7.3. They are not assertions about Frobenius
eigenvalues. The monoidal composition dictionary is
archive cote 2, scan 17, via NN-G010.
Any smaller closed-table score must change at least one
of the three first-cell data in (G034.1); other table
shapes remain outside this class. This is a strict
remaining gate, with no canonical-basis, rank-four,
PH0/PH1 or purity conclusion.



\end{geometricstatement}
\begin{proof} The displayed dictionary transports the proof of NN-G034. Any additional realization, effectiveness or uniformity hypothesis is retained explicitly.\end{proof}


\section{NN-G035: geometric formulation and dictionary}
\begin{geometricstatement}[NN-G035]
Geometric translation and scope.:

For the actual graded sign-representation kernels of
NN-G010, this theorem bounds the possible ordinary
dimensions of the first-cell coefficient object, its
bottom action object, and its correction product by
the $21$ listed triples. Their grading locations are
still unrestricted. The bottom coefficient object
cannot have exactly the absolute trace character
$4D^2(D-1)^2$ in any competitive realization; integral
character factorization and parity force a dimension
contradiction. The finite list records necessary
dimension data, not actual objects for every row.

The dictionary is SGA 3, I, 4.7.2--4.7.3 and archive
cote 2, scan 17 via NN-G010. No equality of formal
$K_0$ classes is treated as an effective realization.
The result shrinks the closed-table search to explicit
mass data, while NN-G029--NN-G032 separately address
the wider unrestricted presentation. Neither statement
supplies a canonical GCT basis, rank-four positivity,
or PH0/PH1.



\end{geometricstatement}
\begin{proof} The displayed dictionary transports the proof of NN-G035. Any additional realization, effectiveness or uniformity hypothesis is retained explicitly.\end{proof}


\section{NN-G036: geometric formulation and dictionary}
\begin{geometricstatement}[NN-G036]
Geometric translation and scope.:

Within the NN-G010 realization by graded sign-representation
kernels, a competitive first-cell bottom-action object
of ordinary dimension $64$ determines its entire character
and, in turn, the unique $2744$ completion in this closed
presentation. An improvement requires dimension at least
$72$ for that object and initially obeys $20$
dimension triples, reduced to $15$ by the addendum below. This statement has no bound on the
grading support of the other surviving objects.
The irreducibility argument is a statement in the
integral grading-character ring, not a claim of
irreducibility of the representations themselves.

The dictionary is SGA 3, I, 4.7.2--4.7.3; the character
evaluation at $q=-1$ and the parity reductions are
integral restrictions of this ring. Actual composition
uses archive cote 2, scan 17 through NN-G010.
The proof neither identifies a canonical GCT basis nor
applies this lower bound to the general extra-position
problem. At this intermediate stage only a lower bound was established.
The later exclusions below close the entire closed-table gate;
the extra-position optimization remains separate.



Scope and geometric content.:

These four base-row exclusions concern the ordinary
dimensions of actual graded $C_2$-representation
structure objects in the matrix-kernel category of
NN-G010, with the same symmetric closed multiplication
table and character constraints. In particular, within
the mass-eight gluing stratum the actual smallest
effective lift is the $2744$-pair object of NN-G033,
uniquely within that presentation. The integral torus
character interpretation is SGA~3, I, 4.7.2--4.7.3;
kernel composition is the supplied archive transcription,
cote~2, scan~17. No new claim about a canonical GCT
representation family or a different multiplication
table follows.


Geometric conclusion at the intermediate stage.:

Within the specified closed presentation of composing graded
$C_2$-representation objects from NN-G010, every effective
lift has at least $2604$ cancelling trivial/sign pairs and
one with $2744$ pairs exists. In the gluing-dimension-eight
stratum, that $2744$ object is the unique minimum in this
presentation. A strict improvement can occur only in the
seventeen displayed ordinary-dimension rows, whose actual
character realization remains a further condition.
The bounds concern genuine effective representation objects,
using SGA~3, I, 4.7.2--4.7.3 and the archive kernel-product
transcription at cote~2, scan~17. They are not statements
about other multiplication-table shapes or canonical GCT
multiplicities.


% NN-G036 further uniform refinement; root whole-proof PASS 2026-11-22T00:06:53Z.

Meaning of this refinement.:

The closed-table optimum is consequently attained at
gluing dimension eight unless a strictly smaller object
occurs in the remaining dimension-twelve stratum.
The proof uses restriction of integral torus characters
at $q=-1$ and $q=i$, together with ordinary dimensions
at $q=1$; it assumes no Weil-weight estimate and no
effectivity of $C,H/W$, or $K$.
Its geometric objects and source dictionary remain the
actual composing graded sign representations of NN-G010,
SGA~3, I, 4.7.2--4.7.3, and the supplied kernel-product
transcription at archive cote~2, scan~17.


Geometric scope.:

The divisibility lemma removes ordinary-dimension
patterns for effective graded sign-representation
objects in the specified closed kernel presentation.
It does not assume positive polynomial quotients:
content and integrality alone supply the character
divisibilities. The forward-row exclusion further
uses actual coefficientwise effectivity of the
structure objects. These are the same mathematical
objects of NN-G010 and SGA~3, I, 4.7.2--4.7.3,
with composition in the supplied archive
cote~2, scan~17 transcription. No change of
multiplication-table presentation is permitted
in this conclusion.


Exact closed-table optimum and geometric conclusion.:

The complete necessary mass gate of NN-G035,
followed by all the uniform exclusions of NN-G036,
leaves exactly the unique $2744$-pair realization
of NN-G033--NN-G034. Therefore the minimum number
of cancelling trivial/sign pairs in the specified
symmetric closed-table presentation is exactly
$2744$, and its polynomial parameter realization
is unique within that presentation.
All cell degrees and contents, all grading support
locations, and all rational auxiliary parameters
were allowed throughout, subject only to integrality
and effectivity of the actual structure constants.
The final certificate projects every surviving
ordinary-dimension row; it does not enumerate
bounded grading degrees to justify the exclusions.

In the NN-G010 category this is an exact minimum
for actual effective graded sign-representation
objects composing by the fixed kernel table.
The grading is the diagonalizable-group dictionary
of SGA~3, I, 4.7.2--4.7.3, and composition is
the supplied archive cote~2, scan~17 transcription.
The theorem concerns this presentation, not all
possible effective multiplication tables.
The unrestricted extra-position optimization still
has the actual upper bound $2744$, the finite
equal-score cutoff of NN-G029/NN-G033, and the
total-mass FPT algorithm of NN-G032; its numerical
minimum remains open.



\end{geometricstatement}
\begin{proof} The displayed dictionary transports the proof of NN-G036. Any additional realization, effectiveness or uniformity hypothesis is retained explicitly.\end{proof}


\section{NN-G037: geometric formulation and dictionary}
\begin{geometricstatement}[NN-G037]
Geometric translation and limits.:

For the graded sign-representation structure objects of
NN-G010, forgetting both grading and the sign action
gives their ordinary nonnegative integer dimensions.
Tensor composition and direct sum make these dimensions
an associative multiplication table.
The initial certificate forces at least $61$ trivial/sign cancellation pairs
among the structure objects of every such effective lift,
even when zero-trace structure positions are populated.
This is a lower bound on actual effective multiplicities,
not on virtual Grothendieck-group representatives.
The grading dictionary is SGA~3, I, 4.7.2--4.7.3;
the kernel-composition interpretation is the supplied
archive cote~2, scan~17 transcription.
No categorical rigidity, canonical-basis property, rank-four
result, or complexity-class separation is inferred.
The shared-quadratic refinement below improves this lower bound;
the unrestricted optimum is not asserted.


Geometric conclusion and remaining unrestricted gate.:

For every effective graded sign-representation lift of the prescribed
trace, ordinary rank gives an associative nonnegative multiplication
table.  The polynomial map $E\mapsto(E_i,E_iE_j)$ places this table
in the corresponding quadratic graph in affine space; its associator
equations are linear after adjoining these product coordinates.
Nonnegative even ranks and the proposed total-rank cap imply
(G037.6), whose exact weighted consequence excludes that entire
quadratic-graph slice.  Thus actual structure objects across all
$810$ allowed positions require at least $104$ cancelling pairs,
even if zero-trace positions are populated.  The SGA~3 I4.7.2--4.7.3
grading dictionary and NN-G010 actual composing-kernel interpretation
(supplied archive cote~2, scan~17 transcription) are unchanged.
At this stage the unrestricted interval is $[104,2744]$;
the subsequent exact certificate below strengthens its lower endpoint.  The stronger
original-$132$-support interval $[900,2744]$ and the exact symmetric
closed-table minimum $2744$ concern different loci and are not
used to prove this unrestricted bound.


Exact geometric conclusion.:

For every effective graded sign-representation multiplication
table with the normalized unit and prescribed trace, all $810$
nonunit-input positions are still permitted.  Its ordinary ranks
determine a point on the graph $Q_{ij}=E_iE_j$ satisfying the
linearized associator equations.  The integer combination above
excludes the whole slice of that graph with at most $210$ added
rank units.  Thus these actual structure objects require at least
$106$ trivial/sign cancelling pairs.  This is a rank obstruction
for the same NN-G010 composing-kernel realization and the same
SGA~3 I4.7.2--4.7.3 grading dictionary (supplied archive cote~2,
scan~17 transcription).  It asserts neither a stronger original
support condition nor a geometric realization of a floating
relaxation point; the original-$132$ bound $900$ and the closed
presentation optimum $2744$ remain separate conclusions.


Geometric meaning and scope.:

Ordinary rank of the effective structure object landing in the
unit basis position is divisible by four; as its trace is zero,
its total number of trivial/sign pairs is even.
This conclusion uses the actual direct-sum and tensor-product
dimensions of the NN-G010 kernel realization.  It is an ordinary-rank
statement, not coefficientwise divisibility of its Laurent character.
The rational basis rescaling above is an algebraic integrality
statement; it does not construct fractional representation objects,
a categorical functor, or a new trace-preserving lift.


Geometric meaning of the graded charge.:

Restrict the grading torus to its subgroup $\mu_2$ over a
characteristic-zero field.  The two grading-parity summands
of each effective unit-output structure object have dimensions
divisible by four.  Since its prescribed sign trace is zero,
each summand consists of equally many trivial and sign objects,
and that common multiplicity is even.  Equivalently, the
effective cancelling-pair character reduces to zero in
$\mathbb F_2[q^{\pm1}]/((q+1)^2)$, the characteristic-two
coordinate ring of this grading subgroup.  This coefficient-ring
calculation retains the possibly nonreduced special fibre; it
does not replace it by its single geometric point.  The
diagonalizable-group grading dictionary is SGA~3 I4.7.2--4.7.3,
with the NN-G010 composing-kernel interpretation and supplied
archive cote~2, scan~17 transcription.  No numerical improvement is claimed from this lemma alone.


Geometric conclusion and scope.:

The actual effective sign-representation structure objects give
nonnegative even rank corrections $E_i$.  Their shared tensor
products lie on the same quadratic graph $Q_{ij}=E_iE_j$.
The difference expressions in (G037.D1) are nonnegative on this
integral rank locus, although they need not be nonnegative on
all real points of the quadratic graph.  Their exact weighted
combination excludes its entire total-rank slice
$\sum E_i\le212$.  Thus every actual lift, including those with
additional zero-trace positions, needs at least $107$ cancelling
pairs.  The grading and actual-kernel dictionaries remain
SGA~3 I4.7.2--4.7.3 and NN-G010 (supplied archive cote~2,
scan~17 transcription).  This integer-rank obstruction does not
assert a new Laurent attainer, geometric positivity of a real
relaxation, or the optimality of the upper bound $2744$.


% Fully reviewed weighted-difference refinement; see CHECK_EVENTS.tsv.

Geometric conclusion and scope.:

The actual effective sign-representation structure objects have
even nonnegative rank corrections, so each integral linear form
$aE_i-bE_j$ belongs to $2\mathbb Z$.  Their shared tensor ranks
lie on the same quadratic graph.  The integer polynomial
(G037.W1) is nonnegative on that rank locus; this is an arithmetic
constraint on ranks, not an assertion that its formal difference
is an effective representation.  The exact combination excludes
the whole rank slice $\sum E_i\le218$ for all additional
zero-trace positions as well as the original positions.  Thus
every normalized-unit effective lift requires at least $110$
cancelling pairs, while the actual $2744$-pair construction remains
the upper bound.  The precise grading and actual-kernel dictionary
is SGA~3 I4.7.2--4.7.3 and NN-G010, with supplied archive cote~2,
scan~17 transcription.  No new Laurent attainer or unrestricted
optimality assertion follows from this rank obstruction.


% Reviewed strengthened certificate: root whole proof/checker PASS; exact replay 2026-11-22T04:52:53Z.

Geometric conclusion and scope.:

The actual effective sign-representation structure objects have
even nonnegative rank corrections, so each integral linear form
$aE_i-bE_j$ belongs to $2\mathbb Z$.  Their shared tensor ranks
lie on the same quadratic graph.  The integer polynomial
(G037.W7) is nonnegative on that rank locus; this is an arithmetic
constraint on ranks, not an assertion that its formal difference
is an effective representation.  The exact combination excludes
the whole rank slice $\sum E_i\le222$ for all additional
zero-trace positions as well as the original positions.  Thus
every normalized-unit effective lift requires at least $112$
cancelling pairs, while the actual $2744$-pair construction remains
the upper bound.  The precise grading and actual-kernel dictionary
is SGA~3 I4.7.2--4.7.3 and NN-G010, with supplied archive cote~2,
scan~17 transcription.  No new Laurent attainer or unrestricted
optimality assertion follows from this rank obstruction.

\subsubsection*{Signed lattice-form certificates for all 810 positions}

This is a necessary-inequality and exact-certificate interface, not a
new numerical bound without a displayed independently checked certificate.
Retain every nonunit-input structure position of NN-G037.  Its added
ordinary rank is $E_i\in s_i\mathbb Z_{\ge0}$, where $s_i=4$ for a
unit-output position and $s_i=2$ otherwise, by the previously proved
unit-output charge lemma and the effective $C_2$ coefficient parity.
No original-support, Laurent-degree or symmetry restriction is imposed.


Geometric conclusion and scope.:

The ranks are those of the actual effective coefficient objects
in the NN-G010/NN-G037 Artin--Tate kernel realization and the
SGA~3 I4.7.2--4.7.3 grading dictionary.  Unit-output divisibility
restricts their ordinary-rank image to the displayed lattice.
Every signed form is an integral function on that rank lattice;
the absence of an integer between $k$ and $k+1$ proves its
nonnegative square-product inequality.  Thus any exact positive-gap
certificate excludes the specified bounded-rank image of the full
normalized-unit effective associativity locus.  This is a necessary
arithmetic condition, not a presentation of that scheme or a
positive-basis construction.  The moment variables and covariance
selection device do not assert a geometric probability measure or
an unprovided trace realization.  A new numerical bound will be
stated only after its full integer certificate has been reviewed.

\subsubsection*{A reduced-associator parity space for all 810 positions}

This is a further necessary condition on the entire unrestricted lift
class, with no degree, symmetry or original-support restriction.
It is not a new numerical lower bound by itself.  Number the prescribed
basis as
\[
 (0,1,2,3,4,5,6,7,8,9)
 =(\mu,\Sigma,g_{1,1},g_{1,12},g_{1,2},g_{1,21},
                 g_{2,1},g_{2,12},g_{2,2},g_{2,21}).
\]
Write $t_{ab}^z\in R=\mathbb Z[q^{\pm1}]$ for its trace table,
with normalized unit $0$.  Let $d$ be any effective associative
dimension table with this trace and this normalized unit.  Put
\[
 d_{ab}^z=|t_{ab}^z|+2e_{ab}^z=t_{ab}^z+2h_{ab}^z,
 \qquad h_{ab}^z=(t_{ab}^z)^-+e_{ab}^z,
 \qquad E_{ab}^z=2e_{ab}^z(1).
\tag{G037.H1}
\]
Here absolute value and negative part are coefficientwise, and
$e_{ab}^z\in\mathbb N[q^{\pm1}]$.  All coefficients with a unit
input have $h=e=0$.  The free positions are the $810$ triples
$a,b\in\{1,\ldots,9\}$, $z\in\{0,\ldots,9\}$, in lexicographic
order, numbered $0$ through $809$.


Standalone geometric conclusion.:

The characteristic-two nonunit-input multiplication coefficients
of the actual graded trace algebra vanish to order at least two
along $q=1$ in the grading torus.  Dividing its coefficients by $(q+1)^2$ and
specializing extracts their second Hasse coefficient.  The
first variation of associativity in $d=t+2h$, on the output
coordinates where the unit terms vanish, then gives the exact
linear equations (G037.H4).  Ordinary-rank specialization of
the effective cancelling-pair objects therefore lands in the
$92$-dimensional linear subspace $\mathcal V\subset\mathbb A^{810}_{\mathbb F_2}$.
In particular the ordinary ranks at the stated $421$ positions
are even, and the corresponding added dimension ranks are
divisible by four.  This conclusion uses the second normal graded piece at the
identity of the grading torus, visible in the quotient by $(q-1)^3$;
it cannot be recovered from its underlying geometric point
alone.  It is a necessary image constraint, not an identification
of a deformation space or a computation of a full Hochschild
cohomology group.  The diagonalizable-group grading dictionary
is SGA~3, Expose~I, 4.7.2--4.7.3, with the actual effective
kernel objects of NN-G010 and the already checked unit-graded
special-fibre calculation.  All remaining finite linear
algebra and Hasse coefficient formulas have been proved here.


\end{geometricstatement}
\begin{proof} The displayed dictionary transports the proof of NN-G037. Any additional realization, effectiveness or uniformity hypothesis is retained explicitly.\end{proof}


\section{NN-G038: geometric formulation and dictionary}
\begin{geometricstatement}[NN-G038]
Algebraic-geometric translation and the remaining gate.:

On the locus where all $113$ specified nonunit coefficients
are invertible, the associativity scheme is explicitly an
open subset of the nineteen-dimensional split torus:
the nineteen recovered coordinates and formulas
(G038.2)--(G038.4) are mutually inverse.
The additional open conditions invert the thirty-two
within-cell bottom numerators; all other required
coefficients are products or ratios of already invertible
coordinates.
This assertion is valid over arbitrary commutative base
rings on that invertible locus, since every cancellation
used in the proof is by an explicitly inverted element.
In affine terms its coordinate ring is a localization of
$\mathbb Z[t_1^{\pm1},\ldots,t_{19}^{\pm1}]$.
An integral lift gives a $K$-point whose reconstructed
coefficients lie in $R$; it need not give an $R$-point
of this localized scheme, since the nineteen coefficients
need not be units in $R$.

The integral grading-character interpretation is SGA~3,
I, 4.7.2--4.7.3; actual effective coefficients give the
matrix-kernel objects of NN-G010, with composition in
the supplied archive cote~2, scan~17 transcription.
The open torus describes the algebraic associativity
parameters, not the coefficientwise positive integer
lattice. The exact $2744$ optimum from NN-G036 applies
to its smaller fixed-$D$ subfamily only.
The remaining support-preserving optimization is now
precisely the nineteen-parameter integral-effectivity
problem above. Tables populating any additional position
remain a separate part of the full $810$-position problem.


Geometric meaning and limitation of the anchor result.:

Within the nineteen-dimensional parameter torus of NN-G038,
(G038.5) specifies a three-coordinate fiber over $\mathbb Q(D)$.
On the integral effective-character points of this fiber,
the square coordinate is forced to (G038.6), and both word-action
coefficient objects have rank at least $384$ after forgetting
the grading. Equality of the two word actions was never assumed.
The restrictions $q=-1,i$ are restrictions of grading characters
to $\mu_2$ and $\mu_4$, in the SGA~3 I4.7.2--4.7.3 dictionary;
actual coefficient objects compose as in NN-G010 and the supplied
archive cote~2, scan~17 transcription. This fiber result does not
restrict other choices of $S$ or the two generator actions, and
it asserts neither a global optimum nor uniqueness of the word
characters of mass $384$.


Geometric translation and exact scope of the global bound.:

The coordinate open torus from NN-G038 parametrizes all
associative tables with this support. On all of its integral
effective-character points with the prescribed trace, the
ordinary-rank functional now has excess at least $688$
over the absolute trace objects, hence at least $344$
copies of the cancelling pair $1+\epsilon$ in total.
The proof first uses the parity characters over
$\mathbb F_2[q,q^{-1}]$ and then the ordinary-dimension
specialization of the actual coefficient objects. SGA~3
I4.7.2--4.7.3 and NN-G010, with the supplied archive cote~2,
scan~17 transcription, give respectively the grading and
composing-kernel dictionary. The scalar dynamic program
is a necessary numerical image of this effective locus;
it neither identifies that image exactly nor replaces the
objects by virtual $K_0$ cancellation. The $344$ bound applies
to every point with the original support, including nonsymmetric
gradings and arbitrary parameter values. It has no asserted
extension to tables using additional support positions.


Geometric meaning and the smaller remaining gate.:

On every integral effective-character point of the full
nineteen-parameter original-support locus, the ordinary-rank
excess is at least $1120$, or $560$ cancelling pairs after
forgetting the grading. The proof uses the generic associative
cell multiplication together with the exact ordinary-rank
images of its coefficients; it imposes no anchor fiber and
no Laurent involution. The scalar table (G038.21) is a point
of the ordinary-dimension associativity locus satisfying these
numerical constraints. Whether it lies in the image of the
integral graded effective locus is a further lifting question.
The grading and actual composing-kernel interpretation remains
SGA~3 I4.7.2--4.7.3, NN-G010, and the supplied archive cote~2,
scan~17 transcription. Formal $K_0$ subtraction is not being
used to realize that missing graded point. This moves the
original-support lower bound from $344$ to $560$ while leaving
the all-$810$-position bound and the $2744$ graded attainer unchanged.


Geometric interpretation.:

The order $r(f)$ is the multiplicity at the characteristic-two
point $q=1$ of the reduction of a nonzero character polynomial
on the grading torus. The two-adic valuation of $f(1)$ records
its specialization to the identity section. Unique factorization
shows that when this valuation is one, a unique integral
irreducible divisor carries that specialization, so product
identities preserve its reduction-order multiset. Applied to
the actual cell coefficient identities, this proves the ordinary
rank alternatives (G038.22) for every point of the integral
effective-character locus. The SGA~3 I4.7.2--4.7.3 grading
interpretation and NN-G010 composing-kernel dictionary remain
as before. The auxiliary factors are polynomial divisors and
need not be effective representation objects; no tensor-factor
or motivic decomposition is inferred from them.


Geometric conclusion and remaining gate.:

The geometric conclusion is now that every integral
effective-character point in the full nineteen-parameter
original-support locus has ordinary-rank excess at least
$1672$, or $836$ cancelling pairs. The improved obstruction
comes from characteristic-two divisor multiplicities before
ordinary dimension is taken. The preceding $344$ and $560$
bounds are historical stages. At this stage the original-support interval is $[836,2744]$,
strengthened below to $[900,2744]$, while the all-$810$-position
minimum remains in $[112,2744]$ and the smaller symmetric
closed presentation retains its exact unique $2744$ optimum.


Geometric translation and limitation.:

The saturated coefficient is a fixed effective character of the
grading torus.  Its integral character polynomial has a finite
prime-divisor decomposition; allocation of those divisors across
the multiplication identity must respect both specialization at
the identity and the order of the characteristic-two divisor there.
These are necessary conditions on all integral character points
of the NN-G038 associativity scheme, with no degree cutoff.
The irreducible polynomial factors used in the proof are not
asserted to be effective characters, subrepresentations, tensor
factors, or geometric summands.  Their use is solely in the UFD
of integral Laurent polynomials.  The geometric conclusion is an
actual exclusion of the displayed ordinary-dimension point from
the integral character-point image, not an exclusion of a virtual
$K_0$ class or a statement about the full 810-position scheme.


Geometric conclusion and the remaining image problem.:

The ordinary-rank excess on every integral effective-character point
of the nineteen-parameter original-support locus is now at least
$1800$, or $900$ cancelling pairs.  This conclusion holds for arbitrary
Laurent gradings.  The saturated strata have fixed integral divisors
on the grading torus, and their finite factor allocations must be
compatible with specialization at the identity and the multiplicity
of the characteristic-two divisor there.  Applying this constraint
before taking dimensions strengthens the numerical image bound.
The scalar table (G038.SF9) belongs to the ordinary associative locus
and passes the stated necessary filters, but whether it lies in the
image of the effective integral character locus remains open.
The intermediate original-support interval $[900,2744]$ is
strengthened below to $[908,2744]$;
the all-$810$-position interval remains $[112,2744]$ and the symmetric
closed presentation retains its exact unique $2744$ optimum.
The grading dictionary is SGA~3 I4.7.2--4.7.3; the actual
composing-kernel interpretation is NN-G010 and the supplied archive
cote~2, scan~17 transcription.  The UFD input is precisely the Stacks
Project, Algebra, Lemmas 10.120.7 and 10.120.10,
\url{https://stacks.math.columbia.edu/tag/034O}; it supplies no
effectivity assertion about the auxiliary polynomial factors.


% Fully reviewed bottom-saturation refinement; see CHECK_EVENTS.tsv.

Geometric meaning and precise gate.:

Saturation of an effective structure object's rank at the absolute
trace rank leaves no cancelling summand, so the two relevant
objects have the fixed grading characters $D$ and $D^2$.
Tensor-composition associativity then forces the neighboring
bottom action character to be $D$ times the generator action.
The remaining relation, reduced in the characteristic-two
grading-coordinate domain and evaluated at the identity, gives
the mod-four obstruction (G038.B1) to actual effective
representation objects.  The calculation is on the full
nineteen-parameter original-support locus; it does not apply to
tables with additional structure positions.  Its dictionaries are
SGA~3 I4.7.2--4.7.3 and NN-G010, with supplied archive cote~2,
scan~17 transcription.  The next projection imposes this
necessary condition before retaining any scalar cell minimum.


Strengthened geometric conclusion.:

The image under ordinary rank of every effective integral-character
point of the full nineteen-parameter original-support locus has
excess at least $1816$.  On a saturated bottom stratum, the two
fixed effective characters force the adjacent action relation;
reducing that relation in the characteristic-two character ring
gives the additional mod-four restriction on the rank image.
The new scalar point (G038.B8) belongs to the ordinary associative
locus and passes the necessary image tests of that intermediate relaxation.  Its lift to
effective integral character data remains unresolved.  The actual
original-support interval is consequently $[908,2744]$; this
argument does not change the separate all-$810$-position bound.
The exact geometric locators remain SGA~3 I4.7.2--4.7.3,
NN-G010 and the supplied cote~2, scan~17 transcription, together
with the previously verified Laurent-UFD lemmas for the retained
saturated-factor tests.


% Fully reviewed structural grading obstruction; numerical projection separate.

Geometric conclusion and scope.:

The saturated word coefficient fixes its grading character and
forces a double divisor at the two fourth roots of unity.  The
characteristic-two order of the mixed-generator coefficient,
together with its rank modulo four, prevents that coefficient
from absorbing even one copy of this divisor.  Associativity
therefore forces the mixed-word character to vanish to order
two there.  Evaluation on the explicit first-order neighborhood
$q=i+\varepsilon$, $\varepsilon^2=0$, is
$P(i)+\varepsilon P'(i)$.  An effective correction of rank four
to $P_0$ cannot cancel both of these terms: its value saturates
the triangle inequality and leaves its derivative purely
imaginary.  Thus the corresponding scalar rank stratum is absent
from the effective integral-character image.  This is a
first-jet obstruction on the grading torus, using SGA~3
I4.7.2--4.7.3 and the actual-kernel dictionary NN-G010
(supplied cote~2, scan~17 transcription).  It neither creates
a Frobenius sign interpretation nor constrains the additional
positions in the separate unrestricted all-$810$ problem.


% Fully reviewed structural grading obstruction; numerical projection separate.

Geometric conclusion.:

The saturated bottom-square character fixes the entire integral
grading character.  The cross-cell tensor identity, reduced twice
in characteristic two, forces divisibility of the first-cell
action character by four and fixes its remaining reduction.
Specialization at the fourth root of unity then forces the
bottom self-square character into $4\mathbb Z[i]$.
The special-fiber map $i\mapsto1$ in characteristic two turns
that constraint into the rank congruence (G038.SQ1).  Thus the
specified ordinary-rank stratum is absent from the effective
integral-character image.  All maps and divisibilities are
explicit in the character ring; the dictionary is SGA~3
I4.7.2--4.7.3 and NN-G010, with supplied cote~2, scan~17
transcription.  No extra-position associativity table, motive,
or trace-purity inference is being assumed.

% Reviewed numerical refinement: full proof/code root review 2026-11-22T04:15:42Z; independent replay 04:18:31Z.

Geometric conclusion and remaining gate.:

The ordinary-rank image of the effective integral-character
locus on the full nineteen-parameter original-support scheme
has excess at least $1864$.  The first new restriction comes
from a forced first-order vanishing condition at a fourth root
of unity; the second comes from the integral Gaussian
specialization of a saturated bottom-square character and its
characteristic-two reduction.  Their finite rank consequences
exclude every scalar slice below $932$ pairs, including the
previous scalar points at $908$ and $910$.  The point (G038.J3)
passes the new necessary tests, but its realization by actual
effective integral characters remains open.  Thus the original
support interval becomes $[932,2744]$.  The separate all-$810$
interval remains $[112,2744]$, and the symmetric closed-table
presentation retains its exact unique $2744$ optimum.
The grading and actual-kernel dictionaries are SGA~3
I4.7.2--4.7.3 and NN-G010 (supplied cote~2, scan~17
transcription); the explicit first-order and Gaussian
specializations are proved in (G038.DR1)--(G038.SQ4).

% NN-G039: whole proof reviewed by root05:22:32 and analytic05:22:39 UTC.

\end{geometricstatement}
\begin{proof} The displayed dictionary transports the proof of NN-G038. Any additional realization, effectiveness or uniformity hypothesis is retained explicitly.\end{proof}


\section{NN-G039: geometric formulation and dictionary}
\begin{geometricstatement}[NN-G039]
Algebraic geometry and the finite trace-form witness.:

For fixed real $a$ and $N$, let
\[
        A_{N,a}=\mathbb R[X]/(J_{N,a}(X)),\qquad
        X_{N,a}=\operatorname{Spec}A_{N,a}.
 \tag{G039.29}
\]
The leading coefficient $N!b_N(a)$ is positive, so division
by it gives a monic equation. Consequently this is an actual
finite locally free real scheme of rank $N$.
The canonical symmetric trace form is
$(f,g)\mapsto\operatorname{Tr}_{A_{N,a}/\mathbb R}(fg)$.
Its exact Grothendieck locator is EGA IV$_4$, 18.2.1--18.2.2,
printed p.~111, especially (18.2.1.1)--(18.2.1.2).
Proposition~18.2.3, p.~112, identifies the nondegenerate
trace locus with the finite etale locus.
The original scan is
\url{https://webusers.imj-prg.fr/~leila.schneps/grothendieckcircle/EGA/EGAIV.4.pdf}.
This is the trace pairing, not a citation for a norm definition.

Every simple nonreal pair supplied above cuts out a degree-two
coprime real factor of $J_{N,a}$. Chinese remaindering gives
an actual direct factor $\mathbb C$ of $A_{N,a}$.
Take the element which is $i$ in this factor and zero in the
complement. It is represented by a real polynomial of degree
less than $N$, and its trace square is exactly
\[
                   \operatorname{Tr}_{A_{N,a}/\mathbb R}(h^2)=-2.
 \tag{G039.30}
\]
For example the corresponding idempotent is obtained by Bezout
between the real quadratic and the remaining polynomial;
within the quadratic factor $(X-\Re\lambda)/\Im\lambda$
represents $i$.
Thus (G039.3) bounds the rank of a finite real algebra with a
negative trace-square witness constructed from the heat
coefficients. Equivalently the finite matrix
$[\operatorname{Tr}(X^{i+j})]_{0\le i,j<N}$ is not positive
semidefinite. This matrix is a trace matrix of the Jensen
algebra, not the original coefficient Toeplitz matrix.

For the hyperbolic degrees in (G039.4), all roots are real.
After splitting into its real primary factors, the trace of
multiplication by $f^2$ is a sum $m_\lambda f(\lambda)^2$,
so the same trace form is positive semidefinite; nilpotent
directions may make it degenerate.
The degree-two complex factor at the failing degrees is etale,
although other factors of the full algebra need not be.
The local scaled family (G039.28), after the quadratic base
change $\varepsilon=q^2$, has the explicit reduced Hermite
branches off $q=0$; it supplies the nonreal quadratic
components underlying (G039.6).

The finite-degree operators in (G039.20) act on
$\mathbb R[z]_{\le N}$, the dehomogenized symmetric-power
representation $\operatorname{Sym}^N(\mathbb R^2)$.
Their displayed symbols are determinants of diagonal pencils
$(z+w)I+C$ with $C$ positive semidefinite, by the two negative
roots computed there. This is a concrete representation-space
positivity statement. The classical representation-ring and
symmetric/exterior-power dictionary is SGA~6, Expos\'e~V,
Example~2.2(b) and Lemma~2.2.1; it does not identify these
real analytic coefficients with integral multiplicities.


Scope addendum: Jensen degree two does not bound PF detecting order.:

The Jensen degree in NN-G039 cannot control the original Toeplitz
detecting order on the class of positive-coefficient entire functions
of order $1/2$, even after the usual coefficient normalizations.
This is a statistic-specific obstruction, separate from the angle-only
counterexample. It does not concern a member of the actual heat family.

For every integer $r\ge2$, put
\[
 \begin{aligned}
 \theta_r&=\frac{\pi}{2r},\qquad m=r+1,\qquad
 x_j=e^{i\theta_r(m+1-2j)}\quad(1\le j\le m),\\
 g_r(z)&=\prod_{j=1}^{m}(1+x_jz),
 \qquad F_r(z)=g_r(z)\cosh\sqrt z.
 \end{aligned}
 \tag{G039.PF1}
\]
Then $F_r$ is real entire of order $1/2$, $F_r(0)=1$, and
every Taylor coefficient is strictly positive. Every allowed
Toeplitz minor of order at most $r$ is strictly positive.
Nevertheless its least nonhyperbolic Jensen degree is exactly two.
In fact, writing $F_r(z)=\sum_{k\ge0}b_kz^k$, one has
\[
 \begin{aligned}
 b_1&=\cot\theta_r+\frac12>0,\qquad
 b_1^2-2b_2=-\frac56,\\
 \operatorname{Disc}\bigl(1+2b_1X+2b_2X^2\bigr)&=-\frac{10}{3}.
 \end{aligned}
 \tag{G039.PF2}
\]
The least negative Toeplitz-minor order is finite and greater
than $r$. The least failing finite unitary rank for
$U\mapsto\det(F_r(U)/F_r(1))$ is likewise finite and greater
than $r$. All positivity and detecting-index conclusions persist under
$\widehat F_r(z)=F_r(z/F_r'(0))$, which also imposes
$\widehat F_r'(0)=1$.

Here and below an allowed Toeplitz minor means one whose ordered
row and column indices are not forced to vanish by upper
triangularity. The coefficients of negative index are zero.
The least index window containing a negative minor also tends
to infinity along this family: a window with at most $r$ rows
and columns cannot contain a minor of order greater than $r$.


Standalone geometric and representation conclusion.:

The alphabet in (G039.PF1) is a diagonal unitary element of
$GL_{r+1}(\mathbb C)$. Formula (G039.PF3) proves that its
characters in the entire partition-width-$r$ rectangle are
nonnegative, and the independent $Q$ specialization makes
the associated allowed skew-character tests strictly positive.
The exact representation-ring exterior/symmetric-power
dictionary is SGA~6, Expos\'e~V, Example~2.2(b) and
Lemma~2.2.1. These real character values are not asserted to
be integral multiplicities or a canonical GCT basis.

For the same function the finite Jensen algebra
\[
        \mathbb R[X]/(1+2b_1X+2b_2X^2)
 \tag{G039.PF6}
\]
is a degree-two field isomorphic to $\mathbb C$, by the
strictly negative discriminant. It is finite etale over
$\mathbb R$ and its trace form already has a square of
trace $-2$, by the explicit construction (G039.30).
The exact trace and etale locators are EGA IV$_4$
18.2.1--18.2.3, pp.~111--112. Thus a rank-two negative
trace-form witness can coexist with strict positivity of
every original Toeplitz/character test through an arbitrarily
large finite rank. A negative Jensen trace form is not itself
one of those original tests.


Standalone geometric and trace-form consequence.:

The head polynomial $P_a$ is a genuine finite analytic zero divisor,
constructed from contour power sums without labeling its ramified
roots. Removing the nonvanishing complementary analytic unit does
not alter its leading weighted local heat jet. After $a=\gamma q^2$
and the finite Jensen interpolation $\varepsilon=q^2$, that jet at
the length-$m$ endpoint $-x^2$ is the explicit Hermite polynomial
with coefficient $c_x(\gamma)=4\gamma x^2-x^4/2$. For
$\gamma<x^2/8$ its normalized local branches are distinct and real;
for $\gamma>x^2/8$ and $m\ge2$ it has a nonreal conjugate pair.
When $m\ge2$, the threshold jet at equality is nonreduced, and no
assertion about its next-order splitting is made.

In the actual rank-$N$ real algebra $\mathbb R[X]/(J_{N,a})$ from
the base proof, all $N\le\gamma_0/a$ below the first threshold have
positive-semidefinite trace pairing. Above it, the supplied simple
complex pair cuts out a $\mathbb C$ factor with a square of trace
$-2$. Hence the least rank in this \emph{specified Jensen family}
at which such a nonreal trace obstruction occurs has exact scaled
limit $x_*^2/8$, or infinity in the reduced-endpoint case. The
trace/base-change dictionary is the same exact EGA IV$_4$,
18.2.1--18.2.3, printed pp.~111--112, used in NN-G039; the analytic
head construction is not asserted to descend to an arithmetic
scheme. No Toeplitz-rank conversion, effective access to the
unknown endpoint multiplicities, or bit-complexity bound follows.


Geometric meaning and limitation.:

The Jensen degree is the rank of its finite real quotient algebra
and trace-form witness from NN-G039.  The other quantities are the
Toeplitz-minor order, finite window and actual unitary testing rank
defined in NN-A029.  Formula (G039.C1) compares their least detecting
sizes along this one specified analytic heat family, using the
proved Hermite local degenerations and the complete tail arguments.
It constructs no map between these different finite test objects,
no degree-preserving trace functor and no universal conversion for
arbitrary entire functions.  In particular it is consistent with
the strict $PF_r$ examples of Jensen detecting degree two in the
preceding scope obstruction.  No new original-PF gate movement is
claimed from combining the two already proved limits.



\end{geometricstatement}
\begin{proof} The displayed dictionary transports the proof of NN-G039. Any additional realization, effectiveness or uniformity hypothesis is retained explicitly.\end{proof}


\section{NN-S001: geometric formulation and dictionary}
\begin{geometricstatement}[NN-S001]
The skeleton indexes a diagram of ordinary modules, hence of
quasi-coherent sheaves when the modules are localized over $\operatorname{Spec}R$.
The fixed fibre is the kernel of the endomorphism $u-1$ of that sheaf.
This is an exact module calculation, not a claim that the skeleton is
an additional vector-space summand. The ordinary affine-module
dictionary is EGA I, Section 1.3; exact edition locators are recorded
in the source register. The nonadditive skeleton lies outside that
ordinary quasi-coherent category.
\end{geometricstatement}
\begin{proof} The displayed dictionary transports the proof of NN-S001. Any additional realization, effectiveness or uniformity hypothesis is retained explicitly.\end{proof}


\section{NN-S002: geometric formulation and dictionary}
\begin{geometricstatement}[NN-S002]
This theorem is already a scheme-theoretic statement: it is the
equalizer of identity and the order-two automorphism of the vector
group. The equalizer itself \emph{does} commute with base change,
as its displayed equations show. What fails is replacing it by the
zero section using its characteristic-zero points and then reducing.
The coordinate calculation provides the full proof. The rank-one
example is already recorded in SGA 3, Expos\'e I, Remark 4.7.1.2,
printed page 29 of the corrected edition dated 13 October 2024
(the editorial addition is identified by its note 39). The new
application here is to the split-semimodule claim; the classical
fixed-scheme phenomenon is not claimed as new. Equalizers and base
change use SGA 3, I, Sections 1.2 and 1.6.
For an $R$-module the analogous issue is the kernel of multiplication
by $2$; its change after tensoring is measured by the associated
$\operatorname{Tor}_1$ sequence. No assertion about a finite-field
Frobenius trace follows merely from the sign-fixed Boolean pair.
\end{geometricstatement}
\begin{proof} The displayed dictionary transports the proof of NN-S002. Any additional realization, effectiveness or uniformity hypothesis is retained explicitly.\end{proof}


\section{NN-S003: geometric formulation and dictionary}
\begin{geometricstatement}[NN-S003]
The amplitude rank-one examples are Tate Weil objects; their
Frobenius eigenvalues give the standard weights of Weil II,
Definition 1.2.2. The distinction between additive group completion
and the Grothendieck group of an exact category is essential in
SGA 6, Expos\'e IV (the definition of $K_0$ and its functorial
operations). No realization functor to that exact category is
constructed by a semiring support map.
\end{geometricstatement}
\begin{proof} The displayed dictionary transports the proof of NN-S003. Any additional realization, effectiveness or uniformity hypothesis is retained explicitly.\end{proof}


\section{NN-S004: geometric formulation and dictionary}
\begin{geometricstatement}[NN-S004]
The arithmetic scheme seen by these finite Bourne quotients is
exactly $\operatorname{Spec}\mathcal O_K$; ring reflection kills
the full idempotent lattice. The extra lattice-prime stratum in
corpus 16's spectral ordinal sum is a semiring spectrum, not an
extra stratum of the ordinary EGA scheme. For abelian $K$, A and
abelian Artin factorization give RH for the common classical
sheet. This equality of analytic invariants imposes no reconstruction
of the arbitrarily chosen support lattice. This is an explicit
failure of a proposed inference from the invariant, not a logical
independence proof about P versus NP. It supplies no efficient
algorithm for an implicitly presented support lattice.
\end{geometricstatement}
\begin{proof} The displayed dictionary transports the proof of NN-S004. Any additional realization, effectiveness or uniformity hypothesis is retained explicitly.\end{proof}


\section{NN-S005: geometric formulation and dictionary}
\begin{geometricstatement}[NN-S005]
Each finite amplitude quotient is the arithmetic closed
subscheme $\operatorname{Spec}(\mathcal O_K/J)$; the lattice
changes its chosen quotient-cardinality statistic from $N(J)$ to
$N(J)+\ell-1$. The theorem is an analytic statement about that
statistic. It does not identify it with an Euler product or with
a trace of a finite Weil object. In the ordinary scheme dictionary
of EGA I, Section 4 (closed subschemes), $J$ remains the same
ideal sheaf; adjoining support classes is extra semiring data.
The residue term is analytic boundary data, not an additional
effective summand in $K_0$.
\end{geometricstatement}
\begin{proof} The displayed dictionary transports the proof of NN-S005. Any additional realization, effectiveness or uniformity hypothesis is retained explicitly.\end{proof}


\section{NN-S006: geometric formulation and dictionary}
\begin{geometricstatement}[NN-S006]
At each finite jet order $r$, the completed local coordinate
algebra is $\mathbb C[z]/(z^{r+1})$ with the semilinear involution
$z\mapsto-z$, conjugating scalars. With $w=-iz$ its invariant
real form is $\mathbb R[w]/(w^{r+1})$, and scalar extension to
$\mathbb C$ recovers the original algebra. Identity (S006.1)
is exactly descent of the reduced logarithmic section to this
real Artin scheme. The proof is coefficientwise descent, so no
unproved effectivity assertion is needed. The infinitesimal
neighbourhood convention is EGA I, Section 4, and finite locally
free descent is the usual affine faithfully flat descent
dictionary. This is a real form and a reflection symmetry;
it supplies no finite-field Frobenius operator and no positive
canonical basis for a nonstandard quantum group.
\end{geometricstatement}
\begin{proof} The displayed dictionary transports the proof of NN-S006. Any additional realization, effectiveness or uniformity hypothesis is retained explicitly.\end{proof}


\section{NN-S007: geometric formulation and dictionary}
\begin{geometricstatement}[NN-S007]
Constructibility requires a finite algebraic stratification and
finite-rank stalks in the SGA 4 constructible category. The first
proof excludes the proposed literal support; the second excludes
the proposed literal stalk at infinity. Ordinary finite-dimensional
Weil representations and the weight definition of Weil II 1.2.2
cannot be assigned to that infinite tail without an additional
finite quotient or realization. The Boolean skeleton's scalar
linearization instead vanishes by NN-S003. These are distinct
operations, neither of which creates the missing finite object.
\end{geometricstatement}
\begin{proof} The displayed dictionary transports the proof of NN-S007. Any additional realization, effectiveness or uniformity hypothesis is retained explicitly.\end{proof}


\section{NN-S008: geometric formulation and dictionary}
\begin{geometricstatement}[NN-S008]
The Connes--Consani input is \emph{Hochschild homology, trace map
and $\zeta$-cycles}, equation (2.2), the common-zero assertion in
Section 2, and Proposition 4.2. Its spectral values become
$-\gamma^2-1/4$ under A. Real negative eigenvalues alone do not
prove self-adjointness or semisimplicity: a two-by-two Jordan
block with a real negative eigenvalue is a counterexample.
In geometric language the reduced support loses infinitesimal
and equivariant data; the real Artin schemes of NN-S006 retain
the latter. The passage from a scheme to its reduction is EGA I,
Section 5.1; no finite-field realization or canonical-basis
positivity follows from this reduction.
\end{geometricstatement}
\begin{proof} The displayed dictionary transports the proof of NN-S008. Any additional realization, effectiveness or uniformity hypothesis is retained explicitly.\end{proof}


\section{NN-S009: geometric formulation and dictionary}
\begin{geometricstatement}[NN-S009]
In the Grothendieck group of representations or vector bundles,
$[V]-[V]=0=[0]$ by its defining relations. The effective cone is
the image of actual objects and does not record a chosen
presentation of a virtual class. SGA 6, Expos\'e IV, Definition
equation (1.5.2) and Definition 2.2, provide this exact group-completion
dictionary. The support annotation is extra grading or
stratification data; it does not restore a lost exact-sequence
or expression history. The conclusion concerns both orbit-ring
multiplicities and representation-valued moonshine coefficients.
\end{geometricstatement}
\begin{proof} The displayed dictionary transports the proof of NN-S009. Any additional realization, effectiveness or uniformity hypothesis is retained explicitly.\end{proof}


\section{NN-S010: geometric formulation and dictionary}
\begin{geometricstatement}[NN-S010]
For a finite reduced scheme of $n$ rational points, the two
statistics are nonemptiness and the trace $n\bmod p$ on its
degree-zero coefficient space. The theorem identifies the
smallest semiring of joint values. The finite-fibre trace is
the sum over fibres in SGA 4, Expos\'e XVII, Theorem 6.2.3;
the coherent mod-$p$ point-count relation used elsewhere in
this catalogue is SGA 7 II, Expos\'e XXII, Theorem 3.1 and
Corollary 3.2. The abstract quotient itself neither gives a
polynomial-size cohomology presentation nor computes its trace.
\end{geometricstatement}
\begin{proof} The displayed dictionary transports the proof of NN-S010. Any additional realization, effectiveness or uniformity hypothesis is retained explicitly.\end{proof}


\section{NN-S011: geometric formulation and dictionary}
\begin{geometricstatement}[NN-S011]
The ordinary determinant is the action on the top exterior power
of a free module; its multiplicativity is an identity in that
additive category. The Boolean permanent records a different
incidence relation, and the counterexample proves that this
support refinement is not a determinant functor with the same
composition axiom. The exterior-power operations are SGA 6,
Expos\'e V, Section 2. The specialized evaluation algorithm above
does not give a small monotone arithmetic circuit for the
universal support polynomial; it uses a graph algorithm. Thus
monotone circuit lower bounds for matching are compatible with
efficient evaluation of this supported-zero pair. No analogous
algorithm for arbitrary sparse quadratic satisfiability follows.
\end{geometricstatement}
\begin{proof} The displayed dictionary transports the proof of NN-S011. Any additional realization, effectiveness or uniformity hypothesis is retained explicitly.\end{proof}


\section{NN-S012: geometric formulation and dictionary}
\begin{geometricstatement}[NN-S012]
The source is \path{OURS_TRANSLATION_SPLIT_ZERO_20261121.md},
Section 9, items 1--6. Its support-shadow assertion is correct.
The phrase ``no supported zeros'' must specify model 2 or model 3.
In characteristic two the displayed OR formula
$x\oplus y\oplus xy$ has an internal supported zero at $x=y=1$.
Thus its proposed monotone equivalence is correct for output-only
absence, but false for absence at every primitive gate over
$\mathbb F_2$. Passing to $\mathbb F_4$ repairs that precise gate
restriction with (S012.2). It does not prove a Boolean circuit lower
bound: unrestricted computation over either fixed finite field still
has the same Boolean circuit complexity.

For a monotone family $f_n$, let $M_n$ be its monotone circuit size
and $C_n$ its general Boolean circuit size. The exact corrected
comparison is between $M_n$ and $C_n$, or the equivalent split models
specified above. In particular, $\mathrm{CLIQUE}\in\mathrm P/\mathrm{poly}$
is equivalent to a polynomial upper bound in the unrestricted amplitude
model. A statement merely that the ratio $M_n/C_n$ is exponential
does not provide that upper bound: as numerical growth conditions,
$M_n=2^{n^a}$ and $C_n=2^{n^a/2}$ have an exponential ratio and neither
is polynomial. This example is an explanation of the missing logical
premise, not an assertion of these actual CLIQUE complexities.
Any claimed equivalence using ``exponential savings'' must include
the polynomial endpoint, or supply an additional theorem proving it.
The same caution applies to E: its proved subexponential running times
give subexponential circuits for the encoded computations, not a
polynomial endpoint. NN-C001 and NN-C019 give the exact supported-zero
removal and instance-dependent isolation consequences available from E.

Over $\mathbb F_2$, put
\[
 X=\operatorname{Spec}
 \mathbb F_2[x_1,\ldots,x_n]/(x_1^2-x_1,\ldots,x_n^2-x_n).
\]
This is the finite disjoint union of the $2^n$ rational points indexed
by Boolean assignments. The amplitude at a gate is a section of its
finite free function algebra; its support is an additional upward-closed
subset of that indexing Boolean poset. In model 3 the support is exactly
the nonvanishing locus of the section. Two nonempty upward-closed
supports contain the all-ones point, and their sections both equal one
there. Their sum vanishes there in characteristic two while the support
tag remains present. This is precisely the addition obstruction proved
above. It is a statement about the actual presentation by functions and
support tags, not about a virtual Grothendieck class or a Frobenius weight.

The finite-free algebra and its base-change trace are the setting of
EGA IV$_4$, 18.2.1, especially equation (18.2.1.2), and 18.2.2,
original printed page 111. The original page image was read in this
run. Base change to $\mathbb F_4$ keeps the same split finite scheme
and supplies the additional constant $\lambda$ used in (S012.2).
Thus every-gate zero avoidance can change under a finite scalar
extension even though the Boolean function and its support poset do
not change. The support tags are extra data beyond this algebra;
NN-S003 explains why additive ringization forgets them. For a general
ring, absence of the zero element at an $R$-valued input is not
silently identified with invertibility or with a base-change-stable
open condition. The general-ring theorem was proved directly in
$G(R)$, and the displayed scheme dictionary is specifically the
finite-field case. Neither purity nor A or M imposes the missing
polynomial upper or lower circuit bound.
\hfill$\square$
\end{geometricstatement}
\begin{proof} The displayed dictionary transports the proof of NN-S012. Any additional realization, effectiveness or uniformity hypothesis is retained explicitly.\end{proof}


\section{NN-S013: geometric formulation and dictionary}
\begin{geometricstatement}[NN-S013]
On the split Boolean finite scheme $X$ of NN-S012, an output fibre
is a clopen union of Boolean points. The inversion proof computes
emptiness and a smallest point of fibres of maps whose coordinate
pullbacks are conjunction monomials. A Boolean amplitude map has
a corresponding morphism between the split Boolean schemes, while
the support map retains an additional clopen tag at each coordinate.
The locus $B_i$ is the clopen intersection where that tag is present
and the amplitude section is zero. Uniform input measure is the
normalized counting measure on the $2^n$ rational points. Thus
(S013.1) is equivalently a lower bound on the rank of the finite
coordinate algebra of one such $B_i$, divided by $2^n$.

Finite-free trace and its base change are exactly EGA IV$_4$,
18.2.1--18.2.2, original printed p.~111; the finite-fibre sheaf trace
is SGA 4, XVII, 6.2.3. Every such finite rational-point scheme has
Frobenius acting as the identity on its degree-zero cohomology.
That purity does not construct its potentially exponentially large
point-indexed basis from the circuit. The missing effectiveness is
the same distinction as NN-C006. The finite extension from
$\mathbb F_2$ to $\mathbb F_4$ changes the available zero-free circuit
presentations; it leaves the underlying Boolean finite set and its
counting distribution unchanged. Neither the trace identity nor
the sign-fixed skeleton supplies a polynomial inverter for the
monotone maps allowed over $\mathbb F_4$.
\hfill$\square$
\end{geometricstatement}
\begin{proof} The displayed dictionary transports the proof of NN-S013. Any additional realization, effectiveness or uniformity hypothesis is retained explicitly.\end{proof}


\section{NN-S014: geometric formulation and dictionary}
\begin{geometricstatement}[NN-S014]
For each input length, the Boolean input algebra is the coordinate
algebra of the finite split scheme with $2^n$ points. Each $B_i$ is a
clopen subscheme, and its finite free rank is $\#B_i$.
The conclusion is therefore a lower bound on the rank of at least one
supported-zero output fibre, with probability measured by rank divided
by $2^n$. The finite-algebra trace convention and its base-change
identity are EGA~IV$_4$, Section~18.2.1, especially equation~(18.2.1.2),
and Section~18.2.2, printed p.~111. This is an explicit finite-scheme
dictionary; no l-adic weight or purity theorem is used to obtain the
probability bound.

The construction of the clopen fibre is implicit in the circuit;
its exponentially many points need not be enumerated. The theorem
improves a necessary cancellation condition and supplies a uniform
attack when that condition fails. It gives neither a lower bound on
the total number of circuit gates nor a proof of P versus NP or
VP versus VNP. E's separate subexponential inversion consequences
retain their original sparsity and locality hypotheses.
\end{geometricstatement}
\begin{proof} The displayed dictionary transports the proof of NN-S014. Any additional realization, effectiveness or uniformity hypothesis is retained explicitly.\end{proof}


\section{NN-S015: geometric formulation and dictionary}
\begin{geometricstatement}[NN-S015]
The finding is a theorem about a smooth projective first-order
deformation over a length-two local base, its relatively free
coherent cohomology line, and the extension class of its absolute
Frobenius after forgetting to the residue field. The exact
dictionary is the based hypersurface sequence and (S015.2).
The \'etale invariance theorem explains precisely which realization
forgets this infinitesimal information. The Grothendieck-group
relation forgets the nonsplit extension for a different reason:
it imposes additivity along exact sequences.

All fibres in a fixed dimension have the same pure special-fibre
Weil data, yet computing this coherent extension class is
parity-counting complete from the succinct equation. This refutes
the unrestricted inference that a known pure finite Weil object,
its zeta function, or its $K_0$ class supplies an efficient
computation of all related deformation data. It does not identify
the extension with a GCT canonical-basis multiplicity, nor decide
P versus NP or VP versus VNP. Such an identification and an efficient
coefficient-evaluation construction would be additional hypotheses.
\end{geometricstatement}
\begin{proof} The displayed dictionary transports the proof of NN-S015. Any additional realization, effectiveness or uniformity hypothesis is retained explicitly.\end{proof}


\section{NN-S016: geometric formulation and dictionary}
\begin{geometricstatement}[NN-S016]
The multipliers form an integral Nullstellensatz refutation of the
quadratic verifier equations with their Boolean axioms. They also
supply low-degree polynomial coordinates for the same smooth cubic
model's arithmetic affine-space isomorphism. This refines the
Boolean-quotient identity of NN-C009 by explicitly constructing the
ordinary ideal multipliers in this special empty Boolean-scheme
case. No polynomial-size expanded certificate, short arithmetic
circuit, or efficient isomorphism-finding algorithm follows from
this degree calculation alone. NN-C029 separately analyzes what E's
symbolic partition routine supplies for circuit descriptions.

The morphism $X_\Phi\to\operatorname{Spec}\mathbb Z$ is a smooth
geometrically integral rational sparse cubic model. For each field fiber,
the cycle classes of its two affine components above each Boolean section
satisfy exactly one relation, $[D_{a,0}]+[D_{a,1}]=0$, and no others.
Its integral Chow ring is the square-zero extension of $\mathbb Z$ by
the free rank-$N$ group in codimension $v$, and its rational $\ell$-adic
cycle map is an isomorphism. The positive-codimension effective class
monoid equals the whole Chow group; it is not an ordered representation
semiring with a pointed cone. On the complex fiber the corresponding
integral homotopy type is the wedge of $N$ even spheres. The proof is
the localization-and-purity calculation and the explicit local
quadric argument above, rather than an inference of topology from
point counts alone.

This places the statement in Grothendieck's cohomology with supports and
Poincar\'e duality, SGA~4 XVI~3.7/3.10(a), XVIII~3.2.5/3.2.6.2, and
the cycle class of SGA~4\,$\frac12$, \emph{Cycle}~2.2.2/2.3.2.
The primary SGA~4 re-edition is
\url{https://www.normalesup.org/~forgogozo/SGA4/16/16.pdf}, pp.~11--13,
and \url{https://www.normalesup.org/~forgogozo/SGA4/18/18.pdf},
pp.~77--78 (original pp.~585--586 for duality).
In particular the purity twist is $(-v)$; the minus sign is explicit
in the primary re-edition's Remark~3.10(a), p.~13.
The localization dictionary for cycles is Stacks Tags~02RX, 02TS and~02TY.
The completeness assertions refer to succinct input equations and basis
rank extraction, while the individual realized cycle classes have the
explicit geometric descriptions just given.
\end{geometricstatement}
\begin{proof} The displayed dictionary transports the proof of NN-S016. Any additional realization, effectiveness or uniformity hypothesis is retained explicitly.\end{proof}


\section{NN-S017: geometric formulation and dictionary}
\begin{geometricstatement}[NN-S017]
The ingredients are smooth integral schemes, a framed regular closed
immersion, its normal cone, and the deformation-to-the-normal-cone
purity construction. The normal-cone/intersection-theory interface is
the same geometric framework as SGA~6, Expos\'e~VII, regular immersions;
the exact homotopy theorem used is the later theorem cited above, not
a theorem attributed to SGA~6. The resulting Tate decomposition is
stronger than the SGA~4/SGA~5 realization calculation in NN-C027.
It supplies actual integral motives after passage to a perfect field,
without a Tate conjecture or a deduction of motives from zeta functions.

Existence of a pure Tate realization, rationality, algebraicity of all
classes, and now an explicit unstable motivic sphere type do not supply
an efficient multiplicity extractor. E evaluates this multiplicity at
its source counting bound; it does not make it polynomial-time.
Neither A nor M is used. No canonical GCT multiplicity identification,
polynomial algorithm for SAT, or separation of VP and VNP follows.
\end{geometricstatement}
\begin{proof} The displayed dictionary transports the proof of NN-S017. Any additional realization, effectiveness or uniformity hypothesis is retained explicitly.\end{proof}


\section{NN-S018: geometric formulation and dictionary}
\begin{geometricstatement}[NN-S018]
Append constant-zero verifier variables, with their Boolean equation
and the equation setting them to zero, until $v\ge3$ is odd.
At most three variables are needed; all old assignments extend
uniquely, all equations retain bounded degree and locality, and
the integral verifier scheme remains the same disjoint union of
$N$ sections in its enlarged affine space. Thus the new cubic
still satisfies every hypothesis of NN-C027 and NN-S017.
Formula~\eqref{eq:S018-torsion} now identifies the indicated
$\mathbb F_2$-dimension with the original exact satisfying count.
This proves all complexity claims, including membership on the
specified syntactic family. E supplies its source counting time,
with the padding contributing only polynomial overhead.

In Grothendieck's language this is an explicit augmented
representation-style $\lambda$-ring of vector bundles, with every
Schur operation determined by its exterior powers and with its
integral gamma filtration computed exactly. Its rational graded
pieces have only gamma degree $v$, corresponding to Tate twist
$v$, while lower
integral graded pieces can have finite torsion of counting-hard
rank. The underlying $K_0$ stays free, as do the ordinary Chow groups
over the perfect fields just specified.
The calculation cautions against replacing the integral gamma
filtration by its rational realization when extracting parity
or multiplicity information. It does not identify these groups
with GCT's canonical multiplicities, prove PH1, or settle P versus
NP or VP versus VNP. Neither A nor M is used.
\end{geometricstatement}
\begin{proof} The displayed dictionary transports the proof of NN-S018. Any additional realization, effectiveness or uniformity hypothesis is retained explicitly.\end{proof}


\section{NN-S019: geometric formulation and dictionary}
\begin{geometricstatement}[NN-S019]
Let $T$ be the split torus over $\mathbb Z$ with independent
characters indexed by all the weights $y,z$. On the finite
etale Boolean scheme
$Y=(\operatorname{Spec}\mathbb Z\sqcup\operatorname{Spec}\mathbb Z)^V$,
give the free rank-one summand at section $x$ the character
equal to its displayed monomial. Its pushforward to
$\operatorname{Spec}\mathbb Z$ is an actual free $T$-module
$\mathcal V_G$ of rank $2^n$, with character $P_G$.
The dictionary between split-torus representations and
character gradings is SGA~3, Expos\'e~I,
4.7.2--4.7.3; direct sum and tensor product induce the
addition and multiplication of the representation
Grothendieck ring, as in SGA~6, Expos\'e~IV,
(1.5.2) and Definition~2.2.

A circuit building this class solely by direct sums,
tensor products, and nonnegative integer multiplicities
from the specified one-dimensional characters is a
monotone arithmetic circuit. It therefore needs
$2^{\kappa n}$ gates. E and NN-C034 instead give a
$2^{o(n)}$ expression in the representation ring using
virtual differences and shared subexpressions. The
resulting class is still the actual effective class
$[\mathcal V_G]$. This proves an exponential cost for
removing subtraction from these precise representation
expressions, despite their bounded-locality counting origin.

The lower bound concerns this arithmetic model. It does
not cover formulas with arbitrary Schur-functor primitives,
lattice-point-counting oracles, coefficient extraction,
hashing operations, arbitrary dual-character generators, or
the entire finite pushforward as a unit-cost primitive. It does not count intermediate
split-semiring values equal to the supported zero: by
NN-S009, a value alone does not retain cancellation
provenance. It gives no lower bound for unrestricted
Boolean circuits and no P versus NP or VP versus VNP
resolution. In particular, the upper bound remains
subexponential, not polynomial. Neither A nor M is used.

For each fixed allowed-tuple bit string $a$, let
\[
 B_n=\mathbb Z[x_v:v\in V]/(x_v^2-x_v:v\in V)
                   \simeq\prod_{x\in\{0,1\}^V}\mathbb Z.
\]
For a vertex or edge factor $f$, let $p_{f,a}\in B_n$ be its
Boolean interpolation polynomial, taking value $a_{f,x|I_f}$
on section $x$.  Explicitly the unary interpolation is
$a_{v,0}(1-x_v)+a_{v,1}x_v$, and the binary interpolation is
the corresponding sum of four degree-two Boolean indicator
monomials multiplied by their allowed bits.  Each $p_{f,a}$
is an idempotent.  Put $e_a=\prod_f p_{f,a}$ and
\[
                 Y_a=\operatorname{Spec}(B_n/(1-e_a)).
\]
This is the finite etale scheme over $\mathbb Z$ consisting of
exactly the allowed assignment sections.  Equivalently one may
use all equations $p_{f,a}=1$ together with the Boolean equations:
for idempotents the ideals $(1-e_a)$ and $(1-p_{f,a}:f)$ agree.
One inclusion follows by expanding $1-\prod_f p_{f,a}$ as a
telescoping sum; the other follows from
$(1-p_{f,a})e_a=0$.

For every field $K$, the scheme $Y_a\times\operatorname{Spec}K$
is nonempty if and only if $\operatorname{Supp}_G(a)=1$.
Thus the nonemptiness of this explicitly given finite etale
subscheme is computable by the displayed $O(n^3)$ monotone
Boolean circuit in its allowed-tuple parameters.  Its rank is
$P_G(a)$ as an ordinary nonnegative integer, while the full
torus character before specialization is the exponentially
costly positive arithmetic object described above.  The map
from positive formal polynomials to Boolean support is a
semiring map with idempotent addition.  It is not an evaluation
of invertible torus characters at zero in the representation
ring, and it does not retain multiplicities.  This distinction
is exactly why the arithmetic lower bound and the Boolean
upper bound coexist.

The integral scheme
\[
 Y_G=\operatorname{Spec}
 \mathbb Z[x_v:v\in V]/
           (x_v^2-x_v, x_ux_v:uv\in E_G)
\]
is the disjoint union of one copy of $\operatorname{Spec}\mathbb Z$
for every independent set.  Put $B=\mathbb Z[z_{v,b}]$.
On the finite locally free $B$-algebra
$B\otimes_{\mathbb Z}\Gamma(Y_G,\mathcal O)$, multiplication
by $\prod_v((1-x_v)z_{v,0}+x_vz_{v,1})$ has trace $J_G(z)$.
This follows directly from its values in the orthogonal
Boolean-idempotent basis.  The trace and arbitrary base-change
interface is EGA~IV4, 18.2.1--18.2.2, printed p.~111, including
equation~18.2.1.2.  The foregoing circuit bound is a uniform
bound for this regular trace function in the compact graph
presentation; it does not assume that the full idempotent
basis is part of the input.

Equivalently, let $T=(\mathbb G_m)^{2n}$ over $\mathbb Z$,
with the specified rail characters.  The actual free
$T$-module
\[
 \mathcal W_G=
 \bigoplus_{I\,\text{independent}}
       \left(\bigotimes_{v\in I}L_{v,1}\right)
          \otimes
       \left(\bigotimes_{v\notin I}L_{v,0}\right)
\]
has character $J_G$.  In the model with these one-dimensional
generators, direct sums, tensor products, nonnegative integer
multiplicities and shared subexpressions, constructing its
class requires $2^{\Omega(n)}$ positive arithmetic gates.
The integral representation Grothendieck ring admits a
$2^{o(n)}$ expression for that same effective class using
virtual differences, by E and NN-C034.  The classical
dictionary is SGA~3, Expos\'e~I, 4.7.2--4.7.3, for torus
gradings, and SGA~6, Expos\'e~IV, (1.5.2) and Definition~2.2,
for the sum, tensor product and Grothendieck-ring operations.
No arbitrary representation functor, entire finite pushforward,
or coefficient-extraction oracle is counted as a unit-cost gate.

For Boolean unary availability data $a$, discard from $Y_G$
the independent-set sections using a forbidden rail.
This is an open-and-closed integral subscheme, whose nonemptiness
is exactly equation~\eqref{eq:S019-independent-support}.
The linear monotone formula thus decides whether an effective
summand is available, while the full character has the proved
exponential positive arithmetic complexity.  Here Boolean
evaluation is the support map on the positive polynomial
semiring; setting a character weight to zero is not asserted
to define a homomorphism of the torus representation ring.
The result distinguishes these explicit arithmetic, Boolean
and representation models.  It gives neither a polynomial
integral circuit nor a lower bound for unrestricted Boolean
circuits, and uses neither A nor M.
\end{geometricstatement}
\begin{proof} The displayed dictionary transports the proof of NN-S019. Any additional realization, effectiveness or uniformity hypothesis is retained explicitly.\end{proof}
